\documentclass[a4paper,	11pt]{report}
% Include preambles
% ---------------- Packages ----------------
%\usepackage{hyphenat}

\usepackage[utf8]{inputenc}
\usepackage[T1]{fontenc}
\usepackage{lmodern}
\usepackage{graphicx}
% \usepackage[fleqn]{amsmath}
\usepackage{amsmath}
\usepackage{amsthm}
\usepackage{mdframed}
\usepackage{listings}
\usepackage{booktabs}
\usepackage{hyperref}
% \usepackage{dtsyntax}
% \usepackage[colorinlistoftodos,prependcaption,textsize=tiny]{todonotes}
\usepackage[normalem]{ulem}
\usepackage{pdflscape}
\usepackage{amssymb}
\usepackage{longtable}
\usepackage{tikz}
\usetikzlibrary{positioning, arrows.meta, babel, shapes.multipart, decorations.pathreplacing}
%\usepackage{algorithm}
\usepackage[authoryear,round]{natbib}
\usepackage{mathpartir}
\usepackage{multicol}
\usepackage{centernot}
\usepackage{hyperref}
\usepackage{array}
\usepackage{tabularx}
\renewcommand{\chapterautorefname}{Capítulo}
\renewcommand{\sectionautorefname}{Sección}
\renewcommand{\figureautorefname}{Figura}
\renewcommand{\tableautorefname}{Tabla}
\renewcommand{\arraystretch}{1.4}

\usepackage{placeins}

\usepackage[a4paper,top=2.5cm,bottom=3cm,inner=2.5cm,outer=2.5cm]{geometry}

\usepackage{needspace}
\usepackage{pifont}

\newcommand{\cmark}{\textcolor{green!50!black}{\ding{51}}}
\newcommand{\xmark}{\textcolor{red!70!black}{\ding{55}}}

\usepackage{multirow}
% For \includegraphics with valign
\usepackage[export]{adjustbox}

\usepackage{tablefootnote}

% M: Tabla en vez de Cuadro
\usepackage[spanish,es-tabla]{babel}

\usepackage{xargs} % Use more than one optional parameter in a new commands

% Algorithms boxes
\usepackage{algorithm}
\usepackage{algorithmic}

% For \qtyrange and \hertz
\usepackage{siunitx}
\usepackage{textcomp}
\DeclareSIUnit{\pixel}{px}
\DeclareSIUnit{\fps}{fps}

\usepackage{subfig}
% \usepackage{subcaption}

% F* code
\usepackage{listings}
\usepackage{xcolor}
\renewcommand{\lstlistingname}{Código}

\definecolor{fstarnavy}{HTML}{254B6E}
\definecolor{fstargreen}{HTML}{356B4A}
\definecolor{fstargray}{HTML}{666666}

% Replace complete keyword tokens without changing longer identifiers.
\makeatletter
\newcommand{\fstarsymbol}[1]{%
  \color{black}%
  \lst@token{\ensuremath{#1}}%
  \lst@length=1\relax
}
\makeatother

\lstdefinelanguage{FStar}{
  sensitive=true,
  alsoletter={'},
  morekeywords={
    module,open,include,
    let,rec,and,in,
    type,noeq,unfold,
    match,with,when,
    if,then,else,
    function,
    val,assume,
    requires,ensures,decreases,
    Lemma,Tot,GTot,Type,Type0,prop,
    squash
  },
  morekeywords=[2]{exists},
  morekeywords=[3]{forall},
  morekeywords=[4]{fun},
  morecomment=[l]{//},
  morecomment=[s]{(*}{*)},
  morestring=[b]"
}

\lstdefinestyle{fstar}{
  language=FStar,
  literate=*
    {/\\}{{\color{black}\ensuremath{\land}}}1
    {\\/}{{\color{black}\ensuremath{\lor}}}1
    {->}{{\color{black}\ensuremath{\rightarrow}}}1
    {^->}{{\color{black}\ensuremath{\rightsquigarrow}}}1
    {==>}{{\color{black}\ensuremath{\Rightarrow}}}1
    {<==>}{{\color{black}\ensuremath{\Leftrightarrow}}}1,
  basicstyle=\ttfamily\small,
  breakindent=1.5em,
  keywordstyle=\color{fstarnavy},
  keywordstyle=[2]\fstarsymbol{\exists},
  keywordstyle=[3]\fstarsymbol{\forall},
  keywordstyle=[4]\fstarsymbol{\lambda},
  commentstyle=\color{fstargray}\itshape,
  stringstyle=\color{fstargreen},
  %numbers=left,
  %numberstyle=\tiny\color{fstargray},
  %numbersep=10pt,
  %stepnumber=1,
  frame=none,
  breaklines=true,
  breakatwhitespace=true,
  showstringspaces=false,
  keepspaces=true,
  columns=fullflexible,
  tabsize=2,
  xleftmargin=1.8em,
  framexleftmargin=1.8em,
  captionpos=b
}

\lstset{style=fstar}

\mdfdefinestyle{barraizquierda}{
  hidealllines=true,
  leftline=true,
  linecolor=black,
  linewidth=0.8pt,
  innerleftmargin=10pt,
  innerrightmargin=0pt,
  innertopmargin=3pt,
  innerbottommargin=3pt,
  skipabove=10pt,
  skipbelow=10pt
}

% ---------------- Macros / Custom commands ----------------
\newcommand{\fstar}[0]{\texorpdfstring{F$^\star$}{F*}}
\newcommand{\nmapsto}{\centernot \mapsto}

% Title page
\newcommand*{\titleGP}{\begingroup
    \centering
    \vspace*{\baselineskip}
    \thispagestyle{empty}

    \rule{\textwidth}{1.6pt}\vspace*{-\baselineskip}\vspace*{2pt}
    \rule{\textwidth}{0.4pt}\\[\baselineskip]

    {\LARGE Formalización de Lógica de Incorrectitud de Separación en \fstar } \\[0.2\baselineskip]

    \rule{\textwidth}{0.4pt}\vspace*{-\baselineskip}\vspace{3.2pt}
    \rule{\textwidth}{1.6pt}\\[\baselineskip]

    \vspace*{2\baselineskip}

    \includegraphics[width=0.3\textwidth]{img/logo-unr.png}

    \vspace*{2\baselineskip}

    \scshape
    Tesina de Grado\\[\baselineskip]
    Licenciatura en Ciencias de la Computación\par

    \vspace*{2\baselineskip}

    {\Large Autora: María Agustina Torrano \\}
    {Legajo: T-2989/1 }
    \vspace*{2\baselineskip}
    {\Large \\ Director: Guido Martínez } \vfill

    \vspace*{2\baselineskip}
    {\scshape \today} \\[0.3\baselineskip]

\endgroup}

\newcommand{\triple}[3]{ \left\{ #1 \right\}~#2~\left\{ #3 \right\} }
\newcommand{\itriple}[3]{ \left\lbrack #1 \right\rbrack~#2~\left\lbrack #3 \right\rbrack }
\theoremstyle{definition}
\newtheorem{definicion}{Definición}[chapter]
\theoremstyle{plain}
\newtheorem{teorema}[definicion]{Teorema}
\newtheorem{lema}[definicion]{Lema}
\newtheorem{corolario}[definicion]{Corolario}
\newtheorem{ejemplo}{Ejemplo}[chapter]
\newtheorem{observacion}[definicion]{Observación}
\surroundwithmdframed[style=barraizquierda]{definicion}
\surroundwithmdframed[style=barraizquierda]{teorema}
\surroundwithmdframed[style=barraizquierda]{lema}
\surroundwithmdframed[style=barraizquierda]{corolario}


\usepackage{combelow}
\usepackage{stmaryrd}

% Comentarios
\usepackage{xcolor}
\newif\ifdraft\drafttrue
\newcommand{\comm}[3]{\ifdraft{{\color{#1}[#2: #3]}}\fi}
\definecolor{darkgreen}{rgb}{0 0.5, 0}
\newcommand{\changed}[1]{\textcolor{magenta}{#1}}
\newcommand{\agustina}[1]{\comm{darkgreen}{Agustina}{#1}}
\newcommand{\guido}[1]{\comm{blue}{Guido}{#1}}
\newcommand{\remove}[1]{\ifdraft{\color{gray}{\sout{#1}}}\fi}
\newcommand{\cambio}[2]{{\color{gray}{\sout{#1}}\color{purple}{#2}}}

% Document settings
% ---------------- Document settings and layout adjustments ----------------

% Page layout settings
\renewcommand\floatpagefraction{.9}
\renewcommand\topfraction{.9}
\renewcommand\bottomfraction{.9}
\renewcommand\textfraction{.1}
\setcounter{totalnumber}{50}
\setcounter{topnumber}{50}
\setcounter{bottomnumber}{50}



\begin{document}

% Title Page
\titleGP
\newpage
\thispagestyle{empty}

% Table of Contents
\tableofcontents

% Abstract
\chapter*{Resumen}
\addcontentsline{toc}{chapter}{Resumen}
Desde los inicios de la verificación formal de software, la atención se ha centrado en asegurar la correctitud de los programas, es decir, la ausencia total de errores.
Con este fin, sistemas como la Lógica de Hoare y la Lógica de Separación proveen las bases matemáticas para demostrar formalmente que un programa se comporta según su especificación.
Sin embargo, la aplicación de estos sistemas a programas de gran tamaño requiere mucho esfuerzo de especificación y prueba.

Frente a estas limitaciones, la industria ha usualmente preferido usar herramientas como analizadores estáticos, herramientas que hacen un mejor esfuerzo para encontrar errores de forma automática.
A diferencia de los demostradores de correctitud, estos analizadores suelen emplear abstracciones y aproximaciones de los comportamientos posibles del programa, perdiendo precisión en pos de la escalabilidad.
Aunque estas técnicas pueden estar respaldadas por fundamentos formales, suelen reportar muchos falsos positivos, es decir, advertencias sobre errores que en realidad no son alcanzables en ninguna ejecución real del programa.
Esto reduce la confianza y productividad de los desarrolladores, quienes deben invertir tiempo en descartar manualmente errores mal advertidos.

La Lógica de Incorrectitud surge como propuesta para invertir el paradigma tradicional al centrarse en la demostración de errores alcanzables,
garantizando la ausencia de falsos positivos en los errores reportados.
Recientemente, esta teoría se combinó con el razonamiento local sobre la memoria de la Lógica de Separación, introduciendo la Lógica de Incorrectitud de Separación.
La mecanización de esta lógica constituye todavía un área de desarrollo reciente.

Presentamos una formalización de la Lógica de Incorrectitud de Separación utilizando \fstar{}, un lenguaje de programación de tipos dependientes y asistente de pruebas.
Por el lado teórico, demostramos las metapropiedades fundamentales de la lógica dentro del asistente de pruebas, haciendo especial énfasis en probar su solidez (\textit{soundness}) matemática para asegurar que cada error detectado sea verdaderamente alcanzable.
Por el lado concreto, utilizamos la lógica formalizada para demostrar la existencia de errores en programas no triviales que hacen uso de estado mutable.

\chapter*{Agradecimientos}
\addcontentsline{toc}{chapter}{Agradecimientos}
No podía faltar este espacio para todas las personas que me acompañaron a lo largo de mi vida,
y especialmente durante estos años,
desde que empecé la facultad hace 6 años hasta hoy,
que presento mi tesis.

En primer lugar,
le agradezco a Guido por tanto tiempo dedicado y por su guía,
porque,
a pesar de la distancia,
siempre estuvo disponible para resolver mis dudas y acompañarme hasta en el más mínimo detalle.

A los profesores de la LCC,
por sus enseñanzas durante estos 6 años y por ayudarme a convertirme en la profesional que soy hoy.

A mis amigos, TODOS.
A los de la facu,
que conocí recién en segundo año,
pero de quienes desde ese entonces no me separé nunca más,
compartiendo mates, birras, juegos, tardes de estudio, rendidas, alegrías y tristezas.
A las de hockey,
porque hacen que un entrenamiento a las 11 de la noche valga la pena y porque celebran cada logro,
sin importar qué tan pequeño sea.
Por último,
y no menos importante,
a los de la escuela,
que son como mi familia y me conocen desde siempre,
que se bancan todas mis excusas cuando no me puedo juntar,
pero que hoy y siempre están.

Le agradezco a Bauti,
por mirarme siempre con amor,
hasta en mis momentos de mayor estrés;
por entenderme y quererme como soy,
y enseñarme tantas cosas que no entiendo cómo entran en un cerebro humano.

Le agradezco a toda mi familia
(abuelos, tíos y primos),
que me confirma cada día que la unión de la familia prevalece por sobre todas las cosas.

A Juanse y Sofi,
mis hermanos,
por escucharme y bancarse todos mis problemas,
porque compartir la vida con ellos la hace más divertida,
incluso en los momentos difíciles.

Por último,
les agradezco a los dos que me enseñan cosas todos los días.
A mi mamá,
por impulsarme siempre a ser mejor y enseñarme la importancia del trabajo;
por motivarme a seguir y superarme o,
en el caso de esta tesina,
a encontrarle un límite.
Y a mi papá,
por su curiosidad e interés en todo lo que hago,
y por enseñarme a ser constante y a pelear siempre por lo que uno quiere.

Todos ellos me convirtieron en la persona que soy hoy,
así que esta tesina es por y para ellos.
Simplemente,
gracias.

%%%%%%%%%%%%%%%%%%%%%%%%%%%%%%%%%%%%%%%%%%%%%%%%%%%%%%%%%%%%%%%%%%%%%%%%%%%%%%%%%%%%
%%%%%%%%%%%%%%%%%%%%%%%%%%%%%%%%%%%% CAPÍTULO 1 %%%%%%%%%%%%%%%%%%%%%%%%%%%%%%%%%%%%
%%%%%%%%%%%%%%%%%%%%%%%%%%%%%%%%%%%%%%%%%%%%%%%%%%%%%%%%%%%%%%%%%%%%%%%%%%%%%%%%%%%%

\chapter{Introducción}
\label{cap:introduccion}
\begin{quote}
``The price of reliability is the pursuit of the utmost simplicity.'' \\
\raggedleft --- C. A. R. Hoare
\end{quote}

\section{Motivación}
\label{sec:motivacion}
La verificación formal busca establecer propiedades sobre el comportamiento de los programas mediante modelos y demostraciones matemáticas.
Con los trabajos fundacionales de \citet{floyd67} y \citet{hoare},
la mayor parte de esta disciplina se ha concentrado en demostrar la \emph{correctitud} de los programas,
entendida como el cumplimiento de sus especificaciones.
Estas especificaciones permiten expresar, entre otras propiedades, que determinadas clases de errores no pueden producirse.

A partir de estas primeras propuestas, se desarrollaron numerosos sistemas lógicos capaces de verificar propiedades cada vez más complejas.
Por ejemplo, la lógica de \citet{concurrence} para programas concurrentes,
la Lógica Temporal de \citet{templog} para la evolución de las ejecuciones,
la Lógica Dinámica de \citet{dinlog},
la Lógica de Hoare Relacional~\citep{hlrel} para comparar ejecuciones,
y diversas lógicas probabilísticas aplicadas,
entre otros casos, a la verificación de construcciones criptográficas~\citep{cripto}.

La aplicación de estos sistemas lógicos a programas de gran escala,
sin embargo, requiere un esfuerzo considerable.
En particular, los métodos deductivos orientados a la correctitud, como la
Lógica de Hoare (HL), la Lógica de Separación (SL, \cite{sl}) y
todas las lógicas mencionadas en el párrafo anterior,
requieren especificar formalmente las propiedades que debe cumplir el programa
y construir demostraciones que cubran todos sus distintos caminos de ejecución.
Además, las demostraciones deben mantenerse y adaptarse a medida que el programa cambia.

La magnitud de este esfuerzo puede observarse en proyectos como seL4~\citep{sel4}.
Aunque se trata de un microkernel de aproximadamente 10.000 líneas de código,
su prueba de correctitud demandó cerca de veinte años-persona.
Al momento de la publicación,
el conjunto de pruebas y especificaciones formales alcanzaba unas 480.000 líneas.
Por otro lado, el Project Everest~\citep{everest} reunió un equipo numeroso de
investigadores e ingenieros para desarrollar una implementación completamente
verificada del ecosistema HTTPS, incluidos TLS, sus varios algoritmos
criptográficos y componentes de parsing y serialización.
Durante sus cinco años de duración, se estima que el proyecto involucró un
esfuerzo acumulado de entre cincuenta y cien años-persona.
Asimismo, aunque este proyecto produjo varios componentes verificados, el
resultado final fue más acotado que la meta inicial, no llegando a verificar
el \emph{handshake} de TLS, pero sí sus componentes criptográficos y de parsing/serialización.
%
Estos proyectos muestran que la verificación de sistemas reales es posible, pero
también que requiere una gran inversión de recursos, lo que ha hecho que haya
poca aplicación en la industria general.

Frente a este costo, la atención ha estado más enfocada en herramientas
de búsqueda de bugs (\emph{bug finders}).
Los analizadores estáticos son una familia de buscadores de bugs
que ha
tenido una adopción mucho más amplia que la verificación formal.
%
Éstos son herramientas automatizadas que
examinan un
programa sin necesidad de ejecutarlo sobre entradas concretas,
con el objetivo de encontrar posibles errores comunes.
En general, los analizadores utilizan aproximaciones que priorizan la escalabilidad y la automatización,
pero pueden perder precisión y dar lugar a falsos positivos.
La necesidad de aproximar está relacionada con la indecidibilidad de las propiedades semánticas no triviales de los programas,
establecida por el Teorema de Rice~\citep{riceteo}.
Por lo tanto, los analizadores estáticos tienen una tensión inevitable
entre reportar todos los posibles errores y evitar falsas advertencias.

En este contexto, eliminar por completo los falsos positivos resulta importante,
ya que un exceso de reportes incorrectos puede reducir la confianza de los desarrolladores en el analizador
y hacer que sus advertencias sean ignoradas.
Por este motivo,
las herramientas utilizadas en la práctica
buscan producir resultados con rapidez y
minimizar las falsas advertencias.
Para alcanzar estos objetivos, los analizadores suelen ajustar sus heurísticas de manera \emph{ad hoc},
sin contar necesariamente con una justificación lógica de por qué cada advertencia corresponde a un error real.

Este interés por eliminar falsos positivos
motivó a \citet{inclog} a proponer una lógica de \emph{incorrectitud},
conocida como Lógica de Incorrectitud (IL).
En lugar de intentar describir todos los comportamientos posibles de un programa
para demostrar que ninguno produce un error,
busca demostrar que ciertos errores pueden alcanzarse mediante ejecuciones concretas.
De este modo, una advertencia puede estar acompañada por un certificado formal
que justifica la existencia del error reportado.

Este razonamiento resulta especialmente útil cuando el programa presenta comportamiento no determinista.
Por ejemplo, si un error depende del resultado de una elección aleatoria o no predecible,
ejecutar el programa varias veces con la misma entrada puede no revelar el fallo.
La IL permite construir una demostración de que existe una ejecución que alcanza el error,
aunque otras ejecuciones puedan terminar normalmente.
Esta demostración puede comprobarse sin ejecutar el programa ni depender de que se repitan las elecciones que producen el fallo.

La IL presentada por O'Hearn es análoga a la HL clásica,
pero se orienta a demostrar que ciertos resultados pueden alcanzarse,
en particular aquellos que corresponden a errores.
%ya que se define sobre un lenguaje de comandos cuyo estado está representado con un \emph{store},
%pero sin un \emph{heap}.
%\guido{al final no estaba mal esto?}
%Por lo tanto, no permite operaciones sobre memoria dinámica y estructuras de datos mutables.
En el ámbito de la correctitud,
la SL fue desarrollada como una extensión a la HL que soporta razonamiento composicional sobre programas con estados complejos.
%A diferencia de la HL, su modelo de estado incorpora,
%además del \emph{store}, un \emph{heap},
%y sus aserciones permiten describir localmente el fragmento de memoria relevante para el comando que se analiza.
La SL incorpora aserciones espaciales que permiten describir localmente el fragmento de memoria relevante para el comando que se analiza.
Recientemente, \citet{isl} introdujeron la Lógica de Incorrectitud de Separación (ISL),
cambiando el foco hacia la incorrectitud pero manteniendo la modularidad brindada por la SL.
Representamos las diferencias mencionadas en la \autoref{tab:relacion-logicas}.
\begin{table}[h]
\centering
\small
\begin{tabularx}{\textwidth}{
    >{\centering\arraybackslash}X
  | >{\centering\arraybackslash}X
  | >{\centering\arraybackslash}X
}
  %\hline
   & \textbf{Correctitud} & \textbf{Incorrectitud} \\
  \hline
  Lógica base & Lógica de Hoare (HL) & Lógica de Incorrectitud (IL) \\
  \hline
  Extensión con razonamiento espacial & Lógica de Separación (SL) & Lógica de Incorrectitud de Separación (ISL) \\
  %\hline
\end{tabularx}
\caption{Relación entre las lógicas.}
\label{tab:relacion-logicas}
\end{table}

La ISL combina, entonces, el razonamiento sobre errores alcanzables de la IL
con el razonamiento local sobre memoria dinámica de la SL.

Sin embargo, para confiar en las conclusiones obtenidas con este sistema,
no basta con definir sus aserciones y reglas de inferencia;
también es necesario demostrar sus \emph{metapropiedades}, es decir,
propiedades del propio sistema lógico y de su relación con la semántica de los programas.
Una de las más importantes es la \emph{solidez} (\emph{soundness}),
que garantiza, en este caso, que toda conclusión obtenida mediante las reglas de la lógica corresponde a
un comportamiento admitido por la semántica.

Estas propiedades pueden demostrarse matemáticamente sobre un papel.
Sin embargo, este tipo de razonamiento puede contener errores u omisiones difíciles de identificar.
Una forma de incrementar la confianza en estos resultados consiste en mecanizarlos dentro de un asistente de pruebas,
de modo que cada definición y cada paso de una demostración sean comprobados por el sistema.
En esta tesina se usa \fstar{}, un lenguaje de programación orientado a pruebas que combina tipos dependientes,
especificaciones lógicas y automatización mediante el solucionador SMT (\emph{Satisfiability Modulo Theories}) Z3.

\section{Resultados del trabajo}
\label{sec:resultados}
El trabajo presenta una formalización mecanizada de la ISL en el asistente de pruebas \fstar{},
junto con una demostración de su solidez (\emph{soundness}) respecto de una semántica operacional explícita.
Esta formalización se desarrolló de manera incremental.

En una primera etapa, se mecaniza la IL sobre un lenguaje imperativo simplificado,
sin operaciones sobre memoria dinámica.
La memoria se representa mediante un \emph{store},
que asocia variables con valores,
y la semántica distingue entre terminaciones normales y con error.
Esta primera formalización permite estudiar de forma aislada las características propias del razonamiento subaproximado,
el no determinismo, la composición secuencial y el razonamiento sobre bucles.
Sobre este modelo se demuestra la solidez de la IL respecto de la semántica operacional.

En una segunda etapa, se extiende el modelo para incorporar memoria mutable y razonamiento espacial.
Para ello, el estado, representado inicialmente con un \emph{store},
se amplía con un \emph{heap}, 
que representa la memoria dinámica,
manteniendo la distinción entre terminaciones normales y con error.
Además, se incorporan operaciones de lectura, escritura, liberación y reserva de memoria.
Sobre esta extensión se formaliza la ISL,
se demuestra su solidez y se mecaniza el Teorema de la Huella,
que permite identificar qué parte de la memoria es necesaria para justificar una ejecución.

A partir de las formalizaciones desarrolladas en estas dos etapas,
se verifican programas ilustrativos que incluyen errores de memoria,
elecciones no deterministas, bucles y estructuras enlazadas.
Para estos programas se construyen derivaciones comprobadas por \fstar{}.
Los teoremas de solidez garantizan que los resultados descritos por estas derivaciones pueden alcanzarse
mediante ejecuciones del programa.
Al proporcionar además estados de error concretos,
se certifica la existencia de ejecuciones que los alcanzan.

\section{Estructura de la tesina}
\label{sec:estructura}
El resto de la tesina se organiza de la siguiente manera.

El \autoref{cap:conceptos-previos} presenta los conceptos teóricos necesarios para el desarrollo del trabajo.
En primer lugar, introduce las nociones de estado, aserción y relación semántica.
Luego, presenta la HL, la IL, la SL y la ISL.
Finalmente, describe las características principales de \fstar{} necesarias para entender la mecanización.

El \autoref{cap:formalizacion-il} desarrolla la formalización de la IL.
Se definen el lenguaje imperativo, su semántica operacional y la demostración de su solidez.

El \autoref{cap:formalizacion-isl} extiende esta formalización para incorporar memoria dinámica.
Se presenta el modelo de \emph{heaps}, aserciones espaciales, \emph{heaps} negativos,
la conjunción separadora y la Regla del Marco.
Además,
se demuestran la solidez de la ISL y el Teorema de la Huella,
y se presenta la construcción de certificados de existencia de error.

El \autoref{cap:evaluacion} presenta distintos ejemplos verificados mediante la formalización,
incluyendo casos relacionados con errores de memoria, no determinismo, bucles y estructuras enlazadas.

Finalmente, el \autoref{cap:conclusiones} resume los resultados obtenidos, los trabajos relacionados,
las limitaciones de la propuesta y posibles líneas de trabajo futuro.

%La formalización y ejemplos se desarrollaron sobre \fstar{} version XXX
%y se encuentran en el repositorio XXX, commit XXX
El código de la formalización y de los ejemplos se desarrolló sobre \fstar{} versión xxx.
Está disponible en el repositorio público
\url{https://github.com/agustorrano/isl-fstar} y el commit xxx corresponde a la versión utilizada en esta tesina.

%%%%%%%%%%%%%%%%%%%%%%%%%%%%%%%%%%%%%%%%%%%%%%%%%%%%%%%%%%%%%%%%%%%%%%%%%%%%%%%%%%%%
%%%%%%%%%%%%%%%%%%%%%%%%%%%%%%%%%%%% CAPÍTULO 2 %%%%%%%%%%%%%%%%%%%%%%%%%%%%%%%%%%%%
%%%%%%%%%%%%%%%%%%%%%%%%%%%%%%%%%%%%%%%%%%%%%%%%%%%%%%%%%%%%%%%%%%%%%%%%%%%%%%%%%%%%

\chapter{Conceptos previos}
\label{cap:conceptos-previos}
\begin{quote}
``Program testing can be quite effective for showing the presence of bugs, but is hopelessly inadequate for showing their absence.'' \\
\raggedleft --- Edsger W. Dijkstra
\end{quote}

En este capítulo se describen los fundamentos teóricos de las lógicas formalizadas a lo largo de este trabajo,
abordando sus conceptos centrales y reglas de inferencia.
Asimismo, se presenta una introducción a \fstar, el lenguaje y asistente de pruebas utilizado para el desarrollo práctico.

En primer lugar, se introduce un marco teórico común,
basado en estados, aserciones y relaciones, que permite representar el comportamiento de los programas.
A partir de esto, se presenta la Lógica de Hoare, la Lógica de Incorrectitud y la Lógica de Separación,
junto con las definiciones formales y las particularidades de cada sistema.
Luego, se introduce la Lógica de Incorrectitud de Separación.

\section{Estados, aserciones y relaciones}
Para describir formalmente el comportamiento de un programa es necesario representar tanto la información
disponible durante su ejecución como las transformaciones que el programa puede producir sobre ella.
Con esto en mente, sea $\Sigma$ un conjunto de estados posibles de un programa.
Un estado $\sigma \in \Sigma$ representa el contenido de la memoria del programa en un momento de su ejecución.

Para describir propiedades de estos estados utilizamos aserciones.
Una aserción sobre estados es un predicado $p : \Sigma \rightarrow \mathsf{Prop}$,
donde $\mathsf{Prop}$ representa el universo de las proposiciones lógicas.
Por lo tanto, para cada estado $\sigma$, la expresión $p(\sigma)$ afirma una propiedad que puede ser verdadera o falsa.
En particular, una aserción puede identificarse con el conjunto de estados que la satisfacen:
$$\llbracket p \rrbracket = \{\sigma \in \Sigma \mid p(\sigma)\}.$$
En lo que sigue se utilizará esta identificación de manera implícita.
Por lo tanto, las expresiones $\sigma \in \llbracket p \rrbracket$ y $p(\sigma)$ se considerarán equivalentes.
Cuando no haya ambigüedad,
escribiremos también $p$ para el conjunto de estados que satisfacen esa aserción.

\subsection*{Relaciones entre estados}
Además de describir conjuntos de estados, es necesario expresar cómo un estado puede transformarse en otro
mediante la ejecución de un programa.
Para esto se utilizan relaciones binarias sobre $\Sigma$.
Sea $R \subseteq \Sigma \times \Sigma$ una relación binaria entre estados.
La pertenencia $(\sigma, \sigma') \in R$ indica que el estado $\sigma$ está relacionado mediante $R$
con el estado $\sigma'$.
Dado un conjunto de estados iniciales,
la postimagen mediante $R$ reúne los estados que pueden alcanzarse desde alguno de ellos a través de esta relación.

\begin{definicion}[Postimagen]
\label{def:post-img}
Sea $R \subseteq \Sigma \times \Sigma$ una relación entre estados y sea $p \subseteq \Sigma$ una aserción.
El operador de postimagen $\operatorname{post}(R) : \mathcal{P}(\Sigma) \rightarrow \mathcal{P}(\Sigma)$
asocia a cada aserción $p$ el conjunto:

$$\operatorname{post}(R)(p) = \{\sigma' \mid \exists \sigma \in p. \; (\sigma, \sigma') \in R\}.$$

Este conjunto contiene exactamente los estados relacionados mediante $R$ con algún estado perteneciente a $p$.
\end{definicion}

\begin{lema}[Monotonía y distributividad de la postimagen]
\label{lem:prop-pi}
Sea $R \subseteq \Sigma \times \Sigma$ una relación y sean $p_1, p_2 \subseteq \Sigma$.
Entonces:
$$p_1 \subseteq p_2 \implies \operatorname{post}(R)(p_1) \subseteq \operatorname{post}(R)(p_2)$$
y
$$\operatorname{post}(R)(p_1 \cup p_2) = \operatorname{post}(R)(p_1) \cup \operatorname{post}(R)(p_2).$$
La primera propiedad expresa la monotonía de la postimagen respecto de la inclusión,
mientras que la segunda expresa su distributividad respecto de la unión.
\end{lema}
\emph{Demostración.} La monotonía se sigue directamente de que todo estado perteneciente a $p_1$ también pertenece a $p_2$.
Para la segunda propiedad,
un estado pertenece a la postimagen de $p_1 \cup p_2$ si posee un estado inicial en $p_1$ o en $p_2$,
lo que equivale a pertenecer a alguna de las dos postimágenes.
$\ensuremath{\blacksquare}$

\subsection*{Programas como relaciones}
Las relaciones anteriores proporcionan el marco matemático necesario para
representar el comportamiento de los programas.
En particular, la ejecución de un comando puede interpretarse como una relación entre
sus posibles estados iniciales y finales.
Sin embargo, además del estado final, es interesante distinguir la manera en que concluye la ejecución.

Para ello, consideramos los \emph{modos de terminación} $\varepsilon \in \{ok, er\}$,
donde $ok$ representa una terminación normal y $er$, una terminación provocada por un error.
Para cada comando $C$ y modo de terminación $\varepsilon$, se define una relación semántica:
$$[\![C]\!]_\varepsilon \subseteq \Sigma \times \Sigma.$$
La pertenencia $(\sigma, \sigma') \in [\![C]\!]_\varepsilon$ indica que existe una ejecución completa de $C$ que empieza en el estado
$\sigma$ y termina en el estado $\sigma'$ con modo de terminación $\varepsilon$.
En este sentido, $[\![C]\!]_\varepsilon$ constituye una representación relacional de una semántica de estilo \emph{big-step},
que relaciona directamente los estados iniciales y finales sin describir los pasos intermedios.

Por ejemplo,
un comando que no modifica el estado y termina normalmente relaciona cada estado inicial consigo mismo mediante $[\![C]\!]_{ok}$.
En cambio,
un comando que señala un error sin modificar el estado establece esa misma relación mediante $[\![C]\!]_{er}$.
Así, el modo de terminación permite distinguir ambos comportamientos,
aunque los estados inicial y final coincidan.

Aplicando la definición de postimagen a la relación semántica de un comando,
se construye la siguiente definición:
\begin{definicion}
\label{def:post-comado}
Sea $C$ un comando, $\varepsilon$ un modo de terminación y $p$ una aserción.
Definimos la postimagen de $p$ mediante $C$ con modo de terminación $\varepsilon$ como
$$\operatorname{post}_\varepsilon(C)(p) = \operatorname{post}([\![C]\!]_\varepsilon)(p).$$
\end{definicion}

Desarrollando esta definición, obtenemos
$$\operatorname{post}_\varepsilon(C)(p) = \{\sigma' \in \Sigma \mid \exists \sigma \in p. \; (\sigma, \sigma') \in [\![C]\!]_\varepsilon\}.$$
Por lo tanto, $\operatorname{post}_\varepsilon(C)(p)$ contiene los estados que pueden obtenerse mediante
una ejecución de $C$ que comienza en un estado que satisface $p$
y termina con el modo $\varepsilon$.

A partir de ahora, en expresiones como $\operatorname{post}_{\varepsilon}(C)(p)$,
se omitirá la escritura explícita de la relación $[\![C]\!]_\varepsilon$,
ya que el comando determina la relación semántica considerada.
El modo de terminación $\varepsilon$ se indica explícitamente para distinguir el tipo de terminación.

Esta representación mediante relaciones permite que un comando
alcance distintos estados finales desde un mismo estado inicial.
También puede ocurrir que no exista un resultado para un modo de terminación determinado.
Por ejemplo, si una ejecución produce una división por cero o desreferencia un puntero nulo,
puede no existir un resultado asociado al modo $ok$, aunque sí uno asociado al modo $er$.
Además, puede ocurrir que no exista ningún resultado para ninguno de los modos considerados,
por ejemplo, cuando todas las ejecuciones divergen y continúan indefinidamente.

\section{Lógica de Hoare (HL)}
\label{sec:hoare}
Esta sección presenta un recordatorio de los conceptos de la Lógica de Hoare
que se utilizarán posteriormente para contrastarla con la Lógica de Incorrectitud.
Se supone que el lector posee cierta familiaridad con los fundamentos de esta lógica,
por lo que no se desarrolla su sistema deductivo completo.

La HL, introducida por \citet{hoare}, es un sistema formal para mostrar la correctitud de programas respecto de una especificación.
Su unidad básica es la \emph{tripleta de Hoare}:
$$\triple{p}{C}{q}$$
donde $p$ y $q$ son aserciones y $C$ es un comando.
La aserción $p$ se denomina \emph{precondición}, mientras que $q$ se denomina \emph{postcondición}.
La tripleta expresa que, si la ejecución de $C$ comienza en un estado que satisface $p$ y termina normalmente,
el estado final debe satisfacer $q$.

El significado de la tripleta también depende de qué se exige sobre la terminación de las ejecuciones.
Existen dos variantes de correctitud en Lógica de Hoare: \emph{total} y \emph{parcial}.
La correctitud total exige que todas las ejecuciones iniciadas desde la precondición terminen,
mientras que la correctitud parcial admite ejecuciones divergentes.
En este trabajo se considera una noción de correctitud parcial con ausencia de terminaciones erróneas:
ninguna ejecución iniciada en un estado que satisface $p$ puede terminar con un error.

Dado que la semántica distingue entre terminaciones normales y terminaciones con error,
la validez de una tripleta se define mediante las postimágenes correspondientes a ambos modos de terminación,
introducidas en la Definición~\ref{def:post-comado}:
$$\models_{\mathrm{HL}} \triple{p}{C}{q} \iff \operatorname{post}_{ok}(C)(p) \subseteq q \land \operatorname{post}_{er}(C)(p) = \varnothing.$$
En esta expresión,
la primera condición impone el cumplimiento de la postcondición para las terminaciones normales,
mientras que la segunda excluye las terminaciones erróneas.

La definición no exige que todas las ejecuciones terminen.
En particular,
una ejecución divergente no produce estados finales asociados a ninguno de los modos considerados y,
por lo tanto,
no invalida la tripleta bajo esta noción de correctitud parcial.
Por el contrario,
si existe una ejecución que comienza en un estado que satisface $p$ y termina con un error,
la tripleta no es válida.

La inclusión
$$\operatorname{post}_{ok}(C)(p) \subseteq q$$
expresa que $q$ es una \emph{sobreaproximación} de los estados finales normales alcanzables.
Por lo tanto, $q$ debe contener todos esos estados,
aunque también puede incluir estados que el programa no alcanza.
Para observar esta propiedad, se presenta un ejemplo.

\begin{ejemplo}
\label{eje:sobreaprox}
Consideremos el comando
$$C \equiv x := x + 2$$
y la precondición $p \equiv x > 0$.
Luego de la asignación, todos los estados finales satisfacen $x > 2$.
Sin embargo, también puede elegirse la postcondición $q \equiv x > 1$,
ya que contiene todos los estados que pueden alcanzarse,
aunque además incluya estados donde $x = 2$,
que no son resultados de una ejecución donde se comienza con $x > 0$.

En consecuencia,
$$\operatorname{post}_{ok}(C)(p) \subseteq \{\sigma \in \Sigma \mid \sigma(x) > 1\}.$$
Además, como una asignación no produce terminaciones erróneas,
se cumple que
$$\operatorname{post}_{er}(C)(p) = \varnothing.$$
Por lo tanto,
se satisfacen ambas condiciones de la definición de validez y la tripleta $\triple{p}{C}{q}$ es válida en la HL.
\end{ejemplo}

El ejemplo anterior muestra la validez de una tripleta desde el punto de vista semántico.
Sin embargo, comprobar directamente la inclusión $\operatorname{post}_{ok}(C)(p) \subseteq q$
puede ser difícil,
ya que requiere considerar todos los estados finales normales que el comando puede alcanzar desde la
precondición y verificar que cada uno satisfaga la postcondición.
Para razonar sobre esta inclusión de manera estructurada,
la HL proporciona un sistema deductivo compuesto por axiomas y reglas de inferencia,
que permite construir tripletas a partir de la estructura de los comandos y
de las relaciones lógicas entre sus aserciones.

Para definir qué significa que una tripleta haya sido construida de esta manera,
introducimos la noción de derivabilidad.

\begin{definicion}
\label{def:deriv-hl}
Escribimos
$$\vdash_{\mathrm{HL}} \triple{p}{C}{q}$$
cuando la tripleta puede obtenerse mediante un número finito de aplicaciones de los axiomas y las reglas de inferencia de la HL.
\end{definicion}

Esta noción permite distinguir entre la validez semántica de una tripleta y
la forma en que esta tripleta se construye dentro del sistema lógico.
A continuación, presentamos un axioma y una regla que permiten obtener tripletas derivables en la HL.

El primer axioma que consideramos es el de asignación.
Este establece que:
\begin{mathpar}
\inferrule []
  {}
  {\vdash_{\mathrm{HL}} \triple{q[e / x]}{x := e}{q},}
\end{mathpar}
donde la expresión $q[e / x]$ representa la sustitución de las apariciones libres de $x$ en $q$ por la expresión $e$.

Este axioma de asignación permite obtener una tripleta cuya postcondición se elige de antemano y
cuya precondición queda determinada por la sustitución de la expresión asignada en dicha postcondición.
Sin embargo, esa precondición no siempre coincide con la que se desea utilizar en el razonamiento.
La Regla de Consecuencia permite adaptar la tripleta a una precondición más restrictiva o una postcondición más general:
\begin{mathpar}
\inferrule []
  {p' \implies p \\ \vdash_{\mathrm{HL}} \triple{p}{C}{q} \\ q \implies q'}
  {\vdash_{\mathrm{HL}} \triple{p'}{C}{q'}}
\end{mathpar}
Esta regla permite
fortalecer la precondición y debilitar la postcondición de una tripleta ya demostrada.

\begin{ejemplo}
Retomamos el comando $C$ del Ejemplo~\ref{eje:sobreaprox}.
Aplicando el axioma de asignación con $q \equiv x > 1$, se obtiene:
$$\vdash_{\mathrm{HL}} \triple{x + 2 > 1}{C}{x > 1}.$$
Como $x > 0 \implies x + 2 > 1$, la Regla de Consecuencia permite derivar:
\begin{mathpar}
\inferrule []
  {x > 0 \implies x + 2 > 1 \\ \vdash_{\mathrm{HL}} \triple{x + 2 > 1}{C}{x > 1}
    \\ x > 1 \implies x > 1}
  {\vdash_{\mathrm{HL}} \triple{x > 0}{C}{x > 1}}
\end{mathpar}
\end{ejemplo}

La HL proporciona así un mecanismo deductivo para construir y adaptar tripletas derivables.
Sobre este mismo marco teórico se define la IL, aunque tiene un objetivo distinto.

\section{Lógica de Incorrectitud (IL)}
\label{sec:loginc}
La IL, introducida por \citet{inclog},
parte del mismo marco relacional, pero cambia el objetivo del razonamiento:
en lugar de demostrar que ciertos comportamientos no pueden ocurrir,
busca demostrar que determinados resultados son alcanzables mediante ejecuciones reales.

Antes de definir formalmente las tripletas de la IL,
es necesario especificar el lenguaje de comandos sobre el cual se razona.
El lenguaje utilizado contiene comandos básicos para modificar variables,
filtrar estados y producir errores,
junto con construcciones para componer, elegir y repetir comandos.

\subsection*{Lenguaje y sintaxis}
Sean $x$ una variable, $e$ una expresión y $B$ una expresión booleana.
La sintaxis de los comandos considerados se define mediante la siguiente gramática:

\begin{definicion}[Sintaxis de los comandos]
Sea $C$ un comando del lenguaje de la IL:
$$C ::= \texttt{skip} \mid \texttt{error} \mid \texttt{assume(}B\texttt{)}
\mid x := e \mid x := \texttt{nondet} \mid C_1; C_2
\mid C_1 + C_2 \mid C^\star$$
\end{definicion}

El comando \texttt{skip} no modifica el estado y termina normalmente.
El comando \texttt{error} representa una terminación provocada por un error.
Por su parte, \texttt{assume(}$B$\texttt{)} conserva solo los estados en los que la expresión booleana $B$ es verdadera.

La asignación $x := e$ evalúa la expresión $e$ y almacena su resultado en la variable $x$.
En cambio, $x := \texttt{nondet}$ asigna a $x$ un valor elegido de manera no determinista.

La composición secuencial $C_1; C_2$ ejecuta primero $C_1$ y,
si este termina normalmente,
continúa con $C_2$.
Por último, $C_1 + C_2$ representa una elección no determinista entre ambos comandos.

El lenguaje utilizado incluye el operador de iteración de Kleene,
denotado por $C^\star$,
que representa la ejecución de un comando $C$ una cantidad finita y no determinista de veces,
incluyendo cero repeticiones.

A partir de estas construcciones pueden expresarse comandos habituales.
Por ejemplo,
un bucle condicionado puede expresarse como:
$$\texttt{while(}B\texttt{)} \; C \equiv \left(\texttt{assume(}B\texttt{)}; C\right)^\star; \texttt{assume(}\neg B\texttt{)}.$$

La expresión anterior permite repetir el comando $C$ mientras $B$ sea verdadera
y finalizar cuando dicha condición deje de cumplirse.
Del mismo modo,
una instrucción condicional puede expresarse como:
$$\texttt{if} \; B \; \texttt{then} \; C_1 \; \texttt{else} \; C_2 \equiv \left(\texttt{assume(}B\texttt{)}; C_1\right) +
\left(\texttt{assume(}\neg B\texttt{)}; C_2\right).$$

Asimismo, puede definirse el comando de aserción mediante:
$$\texttt{assert(}B\texttt{)} \equiv \left(\texttt{assume(}\neg B\texttt{)}; \texttt{error}\right) + \texttt{assume(}B\texttt{)}.$$
En este caso,
la primera rama representa una terminación errónea cuando $B$ es falsa,
mientras que la segunda representa la terminación normal cuando $B$ es verdadera.

\subsection*{Semántica relacional de la IL}
Con la sintaxis anterior,
podemos ahora explicitar la semántica relacional de cada comando.
Para cada comando $C$ y modo de terminación $\varepsilon \in \{ok, er\}$,
la relación $[\![C]\!]_\varepsilon$ describe los pares de estados inicial y final
conectados por una ejecución de $C$ que termina con el modo $\varepsilon$.

En esta sección,
los estados se representan mediante funciones $\sigma : \mathsf{Var} \to \mathsf{Val}$,
donde $\mathsf{Var}$ es el conjunto de variables y $\mathsf{Val}$ es un dominio común de valores que puede tomar cualquiera de ellas.
La notación $[\![e]\!]_\sigma$ denota el valor de la expresión $e$ en el estado $\sigma$.
La actualización $\sigma[x \mapsto v]$ coincide con $\sigma$ en todas las variables salvo en $x$,
a la que se asigna el valor $v$.

Para definir la semántica de $C^\star$,
utilizamos la notación auxiliar $C^n$ para denotar $n$ repeticiones de $C$:
$$C^0 \equiv \texttt{skip}, \qquad C^{n + 1} \equiv C^n; C.$$
La \autoref{fig:semantica-il} reúne las definiciones semánticas;
en particular,
interpreta $C^\star$ mediante la unión de las relaciones asociadas a $C^n$ para $n \in \mathbb{N}$.

\begin{figure}[htbp]
\centering
\begin{minipage}{0.98\textwidth}
\small
$$
\begin{aligned}
\relax[\![\texttt{skip}]\!]_{ok} &= \{(\sigma,\sigma) \mid \sigma \in \Sigma\},
& [\![\texttt{skip}]\!]_{er} &= \varnothing, \\[4pt]
[\![\texttt{error}]\!]_{ok} &= \varnothing,
& [\![\texttt{error}]\!]_{er} &= \{(\sigma,\sigma) \mid \sigma\in\Sigma\}, \\[4pt]
[\![\texttt{assume}(B)]\!]_{ok} &= \{(\sigma,\sigma) \mid [\![B]\!]_\sigma = \texttt{true}\},
& [\![\texttt{assume}(B)]\!]_{er} &= \varnothing, \\[4pt]
[\![x := e]\!]_{ok} &= \{(\sigma,\sigma[x \mapsto [\![e]\!]_\sigma]) \mid \sigma \in \Sigma\},
& [\![x := e]\!]_{er} &= \varnothing, \\[4pt]
[\![x := \texttt{nondet}]\!]_{ok} &= \{(\sigma,\sigma[x \mapsto v]) \mid \sigma \in \Sigma,\ v \in \mathsf{Val}\},
& [\![x := \texttt{nondet}]\!]_{er} &= \varnothing.
\end{aligned}
$$
$$
\begin{aligned}
\relax[\![C_1 + C_2]\!]_{\varepsilon} &= [\![C_1]\!]_{\varepsilon} \cup [\![C_2]\!]_{\varepsilon}, \\[4pt]
[\![C_1;C_2]\!]_{\varepsilon} &= \{(\sigma_1,\sigma_3) \mid
  \exists \sigma_2. \, (\sigma_1,\sigma_2) \in [\![C_1]\!]_{ok} \land (\sigma_2,\sigma_3) \in [\![C_2]\!]_{\varepsilon}\} \\
  &\quad \cup
  \begin{cases}
    [\![C_1]\!]_{er}, & \text{si } \varepsilon = er,\\
    \varnothing, & \text{si } \varepsilon = ok,
  \end{cases} \\[4pt]
[\![C^\star]\!]_{\varepsilon} &= \displaystyle \bigcup_{n \in \mathbb{N}} [\![C^n]\!]_{\varepsilon}.
\end{aligned}
$$
\vspace{0.1em}
\end{minipage}
\rule{\linewidth}{0.4pt}
\caption{Semántica relacional del lenguaje de la IL.}
\label{fig:semantica-il}
\end{figure}

Estas definiciones explicitan los comportamientos relevantes para el razonamiento de la IL.
La elección no determinista se interpreta mediante la unión de las relaciones de sus ramas,
mientras que la composición secuencial distingue entre las ejecuciones en las que el primer comando termina normalmente
y aquellas en las que termina con un error antes de ejecutar el segundo.
En este último caso,
el error se propaga directamente al resultado de la secuencia.

En el caso de la estrella de Kleene,
la interpretación depende del modo de terminación.
Una ejecución de $C^\star$ termina normalmente si realiza cero o más iteraciones de $C$,
todas con terminación normal.
Por eso,
$[\![C^\star]\!]_{ok}$ coincide con la clausura reflexiva transitiva de $[\![C]\!]_{ok}$.
En cambio,
una ejecución termina con error cuando,
tras cero o más iteraciones normales,
una iteración termina con error e impide continuar.
Por lo tanto,
$[\![C^\star]\!]_{er}$ no se define como la clausura reflexiva transitiva de $[\![C]\!]_{er}$.
Las ejecuciones divergentes no quedan representadas en ninguna de estas relaciones,
ya que no alcanzan una terminación normal ni errónea.

Las definiciones anteriores también permiten interpretar las construcciones derivadas
\texttt{if}, \texttt{while} y \texttt{assert},
cuya semántica se obtiene expandiéndolas en términos de los comandos primitivos anteriores.

\subsection*{Tripletas y validez semántica}
Sobre esta semántica se definen las tripletas de la IL,
que permiten especificar los estados iniciales y finales de una ejecución junto con su modo de terminación.
Estas tienen la forma:
$$\itriple{p}{C}{\varepsilon : q},$$
donde $p$ es la \emph{presunción},
nombre utilizado en la IL para la aserción que ocupa el lugar de la precondición en la HL,
$C$ es un comando,
$\varepsilon$ es el modo de terminación considerado y $q$ es la postcondición.
En la HL, el modo de terminación no aparece en la sintaxis de la tripleta,
aunque interviene en su validez semántica.
En cambio, en la IL aparece explícitamente,
ya que permite distinguir si el estado alcanzado corresponde a una terminación normal o a una terminación con error.

Una vez introducida la notación de las tripletas,
definimos formalmente su validez semántica:
$$\models_{\mathrm{IL}} \itriple{p}{C}{\varepsilon : q} \iff q \subseteq \operatorname{post}_\varepsilon(C)(p).$$
Esta inclusión expresa una \emph{subaproximación}:
cada estado descrito por $q$ debe estar respaldado por al menos una ejecución de $C$
que comienza en un estado que satisface $p$ y termina con el modo $\varepsilon$.
La postcondición puede describir solo una parte de los estados alcanzables;
sin embargo, todos los estados que describe deben ser alcanzables.

A diferencia de la HL,
la IL puede describir explícitamente estados producidos por terminaciones con error,
mientras que la HL los excluye al exigir $\operatorname{post}_{er}(C)(p) = \varnothing$.
Sin embargo,
ambas lógicas pueden coincidir cuando la postcondición es igual a la postimagen.

\begin{ejemplo}
\label{eje:comp-HL-IL0}
Tomamos el mismo comando
$$C \equiv x := x + 2$$
y la misma aserción inicial $p \equiv x > 0$ que en el Ejemplo~\ref{eje:sobreaprox}.

En este caso,
$$\operatorname{post}_{ok}(C)(p) = \{\sigma \in \Sigma \mid \sigma(x) > 2\}.$$
Además, como la asignación no produce terminaciones con error,
$\operatorname{post}_{er}(C)(p) = \varnothing.$
Por lo tanto, si $q \equiv x > 2$,
la postcondición coincide con la postimagen.
Se cumplen entonces ambas inclusiones:
$$\operatorname{post}_{ok}(C)(p) \subseteq q \quad \text{y} \quad q \subseteq \operatorname{post}_{ok}(C)(p).$$
En consecuencia, son válidas tanto la tripleta de Hoare
$$\triple{x > 0}{C}{x > 2},$$
como la de incorrectitud
$$\itriple{x > 0}{C}{ok : x > 2}.$$
\end{ejemplo}

El caso anterior muestra que ambas lógicas pueden validar sus respectivas tripletas
cuando la postcondición coincide con la postimagen.
Los siguientes ejemplos muestran cómo,
al modificar la postcondición,
puede obtenerse una tripleta válida en una de las lógicas,
pero no en la otra.

\begin{ejemplo}
\label{eje:comp-HL-IL}
Consideremos nuevamente el comando
$$C \equiv x := x + 2$$
y la aserción inicial $p \equiv x > 0$.

En la IL, es válida, por ejemplo, la tripleta
$$\itriple{x > 0}{C}{ok : x = 3},$$
ya que el estado final $x = 3$ puede obtenerse ejecutando el comando desde un estado inicial donde $x = 1$.
Sin embargo, la tripleta de Hoare correspondiente:
$$\triple{x > 0}{C}{x = 3}$$
no es válida.
En la HL, la postcondición debe cumplirse para toda ejecución normal que comience en un estado que satisfaga la precondición.
Sin embargo, el estado inicial $x = 2$ satisface $x > 0$ y produce el resultado $x = 4$,
que no satisface la postcondición $x = 3$.
\end{ejemplo}

\begin{ejemplo}
\label{eje:comp-HL-IL2}
En el Ejemplo~\ref{eje:sobreaprox}, vimos que la tripleta
$$\triple{x > 0}{C}{x > 1}$$
es válida en la HL.
En cambio, la tripleta
$$\itriple{x > 0}{C}{ok : x > 1}$$
no es válida como tripleta de IL.

Esto se debe a que la postcondición $q \equiv x > 1$ incluye el estado final $x = 2$.
Para alcanzar dicho estado mediante el comando $C$, la ejecución debería comenzar en un estado donde $x = 0$.
Por lo tanto,
el estado final donde $x = 2$ no es alcanzable desde ningún estado que satisfaga la presunción $x > 0$.
\end{ejemplo}

Si se observa en términos de conjuntos,
para un mismo comando $C$ y una misma aserción inicial $p$,
si $q_{\mathrm{IL}}$ es una postcondición elegida como subaproximación de los resultados alcanzables y
$q_{\mathrm{HL}}$ es una postcondición elegida como sobreaproximación,
entonces se cumple
$$q_{\mathrm{IL}} \subseteq \operatorname{post}_{ok}(C)(p) \subseteq q_{\mathrm{HL}}.$$
Por transitividad, también se cumple
$$q_{\mathrm{IL}} \subseteq q_{\mathrm{HL}}.$$
Esta inclusión compara las postcondiciones elegidas,
pero no implica que una tripleta de IL sea automáticamente válida como tripleta de HL utilizando la misma postcondición.
La IL requiere que la postcondición esté contenida en la postimagen,
mientras que la HL requiere que la postimagen esté contenida en la postcondición.

Los tres casos analizados se resumen en la \autoref{tab:comparacion-hl-il}.
En todos ellos se considera el comando $C \equiv x := x + 2$ y la aserción inicial $p \equiv x > 0$.

\begin{table}[ht]
  \centering
  \small
  \renewcommand{\arraystretch}{1.3}
  \begin{tabular}{@{}lccccc@{}}
    \toprule
    Caso
    & Aserción inicial
    & $post_{ok}(C)(p)$
    & Postcondición
    & HL
    & IL \\
    \midrule

    Postcondición exacta
    & $x>0$
    & $x>2$
    & $x>2$
    & \cmark
    & \cmark \\

    Postcondición restringida
    & $x>0$
    & $x>2$
    & $x=3$
    & \xmark
    & \cmark \\

    Postcondición amplia
    & $x>0$
    & $x>2$
    & $x>1$
    & \cmark
    & \xmark \\

    \bottomrule
  \end{tabular}
  \caption{Comparación entre las condiciones de validez de la HL y la IL.}
  \label{tab:comparacion-hl-il}
\end{table}

Hasta ahora, los ejemplos se refirieron únicamente a ejecuciones que terminan normalmente,
es decir, con el modo $ok$.
La IL permite además describir ejecuciones que terminan con un error utilizando el modo $er$.
En ese caso, una tripleta de la forma
$$\itriple{p}{C}{er : q}$$
afirma que cada estado descrito por la aserción $q$ puede alcanzarse mediante una ejecución de $C$
que empieza en un estado que satisface $p$ y termina con el modo $er$.
De este modo, la IL puede utilizarse tanto para certificar resultados normales
como para demostrar la alcanzabilidad de estados finales producidos por errores.

Habiendo considerado ambos modos de terminación,
la definición de validez semántica puede expresarse de manera equivalente en términos de ejecuciones concretas.
La siguiente caracterización hace explícito que cada estado descrito por la postcondición debe estar respaldado por una ejecución real.

\begin{lema}[Caracterización]
\label{lem:caracterizacion}
Los siguientes son equivalentes:
\begin{enumerate}
  \item $\models_{\mathrm{IL}} \itriple{p}{C}{\varepsilon : q}$.
  \item Todo estado que satisface $q$ puede alcanzarse desde algún estado que satisface $p$:
  $\forall \sigma_q \in q. \exists \sigma_p \in p. (\sigma_p, \sigma_q) \in [\![C]\!]_\varepsilon$.
\end{enumerate}
\end{lema}
\emph{Demostración.} Por la definición de validez semántica,
$$\models_{\mathrm{IL}} \itriple{p}{C}{\varepsilon : q} \iff q \subseteq \operatorname{post}_\varepsilon(C)(p).$$
La inclusión anterior equivale a afirmar que,
para todo $\sigma_q \in q$, se cumple
$$\sigma_q \in \operatorname{post}_\varepsilon(C)(p).$$
Por la Definición~\ref{def:post-img},
esto ocurre si y solo si existe un estado $\sigma_p \in p$ tal que
$$(\sigma_p, \sigma_q) \in [\![C]\!]_\varepsilon.$$
Por lo tanto,
$$\forall \sigma_q \in q. \exists \sigma_p \in p. (\sigma_p, \sigma_q) \in [\![C]\!]_\varepsilon.$$
$\ensuremath{\blacksquare}$

Esta caracterización permite vincular la validez semántica con las reglas deductivas que se presentan a continuación.

\subsection*{Reglas de inferencia de la IL}
Las definiciones anteriores caracterizan la validez de una tripleta en términos de la semántica relacional del programa.
Sin embargo, demostrar directamente la inclusión $q \subseteq post_\varepsilon(C)(p)$ puede ser difícil,
ya que exige justificar la alcanzabilidad de cada estado descrito por $q$.
Para organizar este razonamiento,
la IL proporciona,
al igual que la HL,
un sistema deductivo compuesto por reglas de inferencia.
Estas reglas permiten construir derivaciones a partir de la estructura de los comandos y de las relaciones entre aserciones.
Antes de presentarlas, definimos formalmente la noción de derivabilidad.

\begin{definicion}[Derivabilidad en la IL]
\label{def:derivabilidad}
Escribimos
$$\vdash_{\mathrm{IL}} \itriple{p}{C}{\varepsilon : q}$$
cuando la tripleta puede obtenerse mediante un número finito de aplicaciones de las reglas de inferencia de la IL.
\end{definicion}

La derivabilidad es una noción sintáctica,
ya que depende de las reglas del sistema y de la existencia de una derivación finita.
La validez, en cambio, es una noción semántica definida mediante la postimagen.

A continuación, se presentan las reglas de inferencia de la IL,
resumidas en la \autoref{fig:reglas-il}.
Algunas describen comandos básicos, como \texttt{skip}, \texttt{error},
las asignaciones y \texttt{nondet}.
Otras permiten combinar derivaciones,
como las reglas de \textsc{Consequence}, \textsc{Disjunction} y \textsc{Seq}.
Finalmente, las reglas asociadas al operador de Kleene permiten razonar sobre repeticiones finitas de un comando.

\begin{figure}[htbp]
\centering
\begin{mathpar}
\inferrule [Consequence]
  {p \implies p' \\ \vdash_{\mathrm{IL}} \itriple{p}{C}{\varepsilon : q} \\ q' \implies q}
  {\vdash_{\mathrm{IL}} \itriple{p'}{C}{\varepsilon : q'}}

\inferrule [Disjunction]
  {\vdash_{\mathrm{IL}} \itriple{p_1}{C}{\varepsilon : q_1} \\ \vdash_{\mathrm{IL}} \itriple{p_2}{C}{\varepsilon : q_2}}
  {\vdash_{\mathrm{IL}} \itriple{p_1 \lor p_2}{C}{\varepsilon : q_1 \lor q_2}}

\inferrule [Skip]
  {}
  {\vdash_{\mathrm{IL}} \itriple{p}{\texttt{skip}}{ok : p}}

\inferrule [Seq]
  {\vdash_{\mathrm{IL}} \itriple{p}{C_1}{ok : q} \\ \vdash_{\mathrm{IL}} \itriple{q}{C_2}{\varepsilon : r}}
  {\vdash_{\mathrm{IL}} \itriple{p}{C_1; C_2}{\varepsilon : r}}

\inferrule [SeqEr]
  {\vdash_{\mathrm{IL}} \itriple{p}{C_1}{er : r}}
  {\vdash_{\mathrm{IL}} \itriple{p}{C_1; C_2}{er : r}}

\inferrule [Kleene zero]
  {}
  {\vdash_{\mathrm{IL}} \itriple{p}{C^\star}{ok : p}}

\inferrule [Kleene non-zero]
  {\vdash_{\mathrm{IL}} \itriple{p}{C^\star; C}{\varepsilon : q}}
  {\vdash_{\mathrm{IL}} \itriple{p}{C^\star}{\varepsilon : q}}

\inferrule [Kleene Variant]
  {\vdash_{\mathrm{IL}} \itriple{p(n) \land \operatorname{nat}(n)}{C}{ok : p(n + 1) \land \operatorname{nat}(n)}}
  {\vdash_{\mathrm{IL}} \itriple{p(0)}{C^\star}{ok : \exists n. p(n) \land \operatorname{nat}(n)}}

\inferrule [Choice ($i$ = 1 o 2)]
  {\vdash_{\mathrm{IL}} \itriple{p}{C_i}{\varepsilon : q}}
  {\vdash_{\mathrm{IL}} \itriple{p}{C_1 + C_2}{\varepsilon : q}}

\inferrule [Error]
  {}
  {\vdash_{\mathrm{IL}} \itriple{p}{\texttt{error}}{er : p}}

\inferrule [Assume]
  {}
  {\vdash_{\mathrm{IL}} \itriple{p}{\texttt{assume(}B\texttt{)}}{ok : p \land B}}

\inferrule [Assign]
  {}
  {\vdash_{\mathrm{IL}} \itriple{p}{x := e}{ok : \exists x'. p[x' / x] \land x = e[x' / x]}}

\inferrule [Nondet]
  {}
  {\vdash_{\mathrm{IL}} \itriple{p}{x := \texttt{nondet}}{ok : \exists x'. p[x'/x]}}

\end{mathpar}
\rule{\linewidth}{0.4pt}
\caption{Reglas de inferencia de la IL. Adaptadas de~\citep{inclog}.}
\label{fig:reglas-il}
\end{figure}

Las reglas de \textsc{Consequence} y \textsc{Disjunction} reflejan propiedades generales de la postimagen.
La primera permite ampliar la presunción y restringir la postcondición de una tripleta ya derivada.
La segunda combina derivaciones obtenidas a partir de distintos conjuntos de estados iniciales.
Esta última regla se apoya en la preservación de uniones de la postimagen,
establecida en el Lema~\ref{lem:prop-pi}.
La interpretación particular de la Regla de Consecuencia en el contexto subaproximado se analizará en la siguiente subsección.

La regla \textsc{Skip} expresa que el comando \texttt{skip} termina normalmente y preserva la aserción inicial.
Por el contrario, la regla \textsc{Error} expresa que el comando \texttt{error} termina en el modo $er$
y conserva el estado descrito por la presunción.

Para el caso de secuenciación hay dos reglas.
La primera regla (\textsc{Seq}) contempla el caso en que el
primer comando termina normalmente y su postcondición se usa como presunción
del segundo comando.
En este caso, el modo de terminación final puede ser $ok$ o $er$.
%
La segunda regla (\textsc{SeqEr}) contempla el caso en que el
primer comando termina con un error; por lo tanto, el segundo comando no llega a
ejecutarse.

La regla \textsc{Choice} permite construir una derivación seleccionando una de las dos ramas del comando no determinista $C_1 + C_2$.
Esto es suficiente porque la IL busca certificar la alcanzabilidad de resultados determinados
y no describir simultáneamente todos los comportamientos posibles.
Por su parte, \textsc{Assume} restringe los estados iniciales considerados a aquellos que,
además de satisfacer $p$, satisfacen $B$.
Por eso, la postcondición es $p \land B$.

En la regla de \textsc{Assign},
$x'$ es una variable auxiliar que representa el valor que tenía $x$ antes de ejecutar el comando.
De este modo, la postcondición puede relacionar el valor final de $x$ con su valor inicial,
evitando perder información al sobrescribir la variable.
La variable auxiliar $x'$ se elige fresca,
de modo que no aparezca en la aserción $p$ ni en la expresión $e$;
las sustituciones se realizan evitando la captura de variables.
Por ejemplo, para la asignación $x := x + 2$ y la presunción $x > 0$,
la regla permite obtener la postcondición
$$\exists x'. \; x' > 0 \land x = x' + 2,$$
que es equivalente a $x > 2$.

Por su parte, la regla \textsc{Nondet} modela una asignación de un valor arbitrario.
Al igual que en la regla anterior,
la variable auxiliar $x'$ representa el valor que tenía $x$ antes de la asignación.
También aquí se elige una variable auxiliar fresca.
Al cuantificar existencialmente ese valor,
la regla elimina de la postcondición las restricciones que dependían del valor anterior de $x$.
Esto significa que el nuevo valor de $x$ puede ser cualquiera y no depende del valor que tenía antes.

El análisis de $C^\star$ se realiza mediante tres reglas.
La regla \textsc{Kleene zero} representa la ejecución de cero iteraciones.
Si $p$ describe el estado inicial,
entonces ese mismo estado también es el estado final de una ejecución de $C^\star$ que no ejecuta $C$.

La regla \textsc{Kleene non-zero} permite agregar una iteración a una ejecución ya establecida.
La regla parte de una derivación para la secuencia $C^\star; C$,
que representa una cantidad finita de iteraciones seguida de una iteración adicional.
Si esta secuencia alcanza una postcondición $q$ con modo de terminación $\varepsilon$,
entonces también puede derivarse que $C^\star$ alcanza $q$ con ese mismo modo.

La forma $C^\star; C$ resulta particularmente útil para razonar sobre errores.
Las iteraciones representadas por $C^\star$ pueden terminar normalmente,
mientras que la iteración adicional puede ser la que termine con error.
En ese caso, el modo de terminación de toda la secuencia es $er$.

Finalmente, la regla \textsc{Kleene Variant} permite razonar sobre la cantidad de iteraciones ejecutadas.
Para ello, se utiliza una familia de aserciones $p(n)$,
donde $p(n)$ describe los estados obtenidos después de $n$ iteraciones.
Además,
se utiliza la aserción $\operatorname{nat}(n)$,
la cual indica que $n$ es un número natural,
incluido el cero.
La regla exige que,
para cada número natural $n$,
todo estado que satisface $p(n + 1)$ pueda alcanzarse mediante una ejecución de $C$ con terminación normal desde algún estado que satisface $p(n)$.
Bajo esta premisa,
todo estado que satisface $p(n)$ para algún $n$ puede alcanzarse mediante $C^\star$ desde un estado que satisface $p(0)$,
también con terminación normal.

\subsection*{La Regla de Consecuencia en la IL}
Una de las diferencias más contraintuitivas entre la IL y la HL aparece en la Regla de Consecuencia.
En la HL, la Regla de Consecuencia permite debilitar la postcondición.
De esta manera, una propiedad que describe todos los estados finales alcanzables
puede reemplazarse por otra más general que también los contenga.
Por ejemplo, a partir de
$\triple{p}{C}{q_1 \land q_2}$ puede obtenerse $\triple{p}{C}{q_1}$.
Esto es válido porque, si todos los estados alcanzables satisfacen $q_1 \land q_2$, también satisfacen $q_1$.

En la IL, este debilitamiento de la postcondición, en general, no es válido.
Al reemplazar $q_1 \land q_2$ por $q_1$,
la nueva postcondición puede incluir estados que satisfacen $q_1$,
pero que no satisfacen $q_2$.
Por lo tanto, algunos de los estados incorporados podrían ser inalcanzables.

La Regla de Consecuencia de la IL produce el efecto contrario:
permite fortalecer la postcondición y conservar solo una parte de los estados alcanzables,
concentrándose en los comportamientos que se desean certificar.
\begin{mathpar}
\inferrule []
  {p \implies p' \\ \vdash_{\mathrm{IL}} \itriple{p}{C}{\varepsilon : q} \\ q' \implies q}
  {\vdash_{\mathrm{IL}} \itriple{p'}{C}{\varepsilon : q'}}
\end{mathpar}

Interpretando las aserciones como conjuntos,
las dos implicaciones que acompañan la regla equivalen a $p \subseteq p'$ y $q' \subseteq q$.
Por lo tanto, la regla permite ampliar el conjunto de estados iniciales y restringir el conjunto de estados finales.

La dirección de estas inclusiones puede justificarse utilizando las definiciones anteriores y el Lema~\ref{lem:prop-pi}.
Este lema establece que, si se amplía la presunción,
también puede ampliarse el conjunto de estados finales alcanzables:
$$p \subseteq p' \implies post_\varepsilon(C)(p) \subseteq post_\varepsilon(C)(p').$$
En otras palabras,
toda ejecución que comienza en un estado perteneciente a $p$ también comienza en un estado perteneciente a $p'$,
ya que $p'$ contiene todos los estados de $p$.
Por lo tanto, los resultados alcanzables desde $p$ también están incluidos entre los estados alcanzables desde $p'$.

Si la tripleta original es semánticamente válida, entonces:
$$q \subseteq post_\varepsilon(C)(p).$$
Luego,
$$q' \subseteq q \subseteq post_\varepsilon(C)(p) \subseteq post_\varepsilon(C)(p').$$
Por transitividad,
$$q' \subseteq post_\varepsilon(C)(p').$$
Esta inclusión es la condición de validez semántica de la tripleta $\itriple{p'}{C}{\varepsilon : q'}$.

\begin{ejemplo}
Supongamos que una derivación obtiene la postcondición $q_1 \lor q_2$,
donde $q_1$ y $q_2$ describen los resultados de dos ramas posibles.
Como $q_1 \implies q_1 \lor q_2$,
la Regla de Consecuencia permite derivar una tripleta cuya postcondición es solamente $q_1$:
\begin{mathpar}
\inferrule []
  {\vdash_{\mathrm{IL}} \itriple{p}{C}{\varepsilon : q_1 \lor q_2} \\ q_1 \implies q_1 \lor q_2}
  {\vdash_{\mathrm{IL}} \itriple{p}{C}{\varepsilon : q_1}}
\end{mathpar}
De este modo, la derivación conserva solo los resultados descritos por $q_1$.
Cuando $q_1$ corresponde a una rama particular del programa,
esto equivale a ignorar las demás ramas.

Esta selección no es realizada automáticamente por la lógica.
La rama que se conserva tiene que quedar especificada explícitamente en la derivación.
\end{ejemplo}

El ejemplo anterior muestra la regla de manera abstracta.
Se observa ahora por qué, en un programa con ramas condicionales,
esas condiciones deben conservarse en la postcondición.

\begin{ejemplo}
Denotamos por $C$ al comando representado por el siguiente fragmento:
  \begin{verbatim}
  if (x is even) {
      if (y is odd) {
          z := 42;
      }
  }
  \end{verbatim}
y consideramos la presunción $p = \{\sigma \in \Sigma \mid \sigma(z) = 11\}$.

¿La aserción $q = \{\sigma \in \Sigma \mid \sigma(z) = 42\}$
es una subaproximación de los estados alcanzables por este código partiendo de la presunción?

Para que la tripleta $\itriple{p}{C}{ok : q}$ sea válida se exige:
$$q \subseteq post_{ok}(C)(p).$$
Sin embargo, el estado $\sigma = [z \mapsto 42, x \mapsto 1, y \mapsto 2]$ pertenece a $q$,
pero no pertenece a $post_{ok}(C)(p)$.
En efecto, como el programa no modifica $x$ ni $y$,
sus valores se mantienen durante la ejecución.
Para que se ejecute la asignación $z := 42$,
$x$ debe ser par e $y$ debe ser impar.
En el estado considerado, $x = 1$ es impar e $y = 2$ es par,
por lo que la asignación no puede ejecutarse.
Por lo tanto,
$$q \nsubseteq post_{ok}(C)(p),$$
y la tripleta no es válida.

Una postcondición subaproximada válida debe conservar las condiciones que justifican la ejecución de la asignación.
Por ejemplo:
$$\itriple{z = 11}{C}{ok: z = 42 \land (x \bmod 2 = 0) \land (y \bmod 2 \ne 0)}.$$
\end{ejemplo}

\subsection*{Razonamiento sobre bucles}
El análisis de los bucles iterativos constituye uno de los principales desafíos de la verificación formal de software.
En la HL,
razonar sobre un bucle requiere encontrar un invariante
que se preserve entre iteraciones y sea suficientemente fuerte para demostrar la propiedad buscada.
Esta tarea no es trivial y a menudo requiere invención humana,
lo que dificulta la escalabilidad y automatización de las herramientas.
Aunque existen técnicas automáticas,
no hay un procedimiento que resuelva esta tarea para programas arbitrarios.

La IL no necesita encontrar un invariante que describa todas las ejecuciones del bucle
cuando el objetivo es certificar una ejecución particular.
En ese caso, puede concentrarse en una cantidad finita de iteraciones.

Una forma de aplicar este enfoque consiste en \emph{desenrollar el bucle}, es decir,
reemplazarlo por una cantidad finita de repeticiones de su cuerpo mediante
las reglas \textsc{Kleene zero} y \textsc{Kleene non-zero}.
Si la derivación obtenida de este modo produce una tripleta con modo de terminación $er$ y postcondición satisfacible,
entonces certifica la existencia de una ejecución real que alcanza ese error,
sin necesidad de describir los demás comportamientos posibles del bucle.

Para los casos en que se requiere un razonamiento paramétrico,
la regla \textsc{Kleene Variant} puede interpretarse hacia atrás.
La regla parte de una derivación de la forma:
$$\vdash_{\mathrm{IL}} \itriple{p(n) \land \operatorname{nat}(n)}{C}{ok : p(n + 1) \land \operatorname{nat}(n)}.$$
Semánticamente, si la tripleta $\itriple{p(n) \land \operatorname{nat}(n)}{C}{ok : p(n + 1) \land \operatorname{nat}(n)}$ es válida,
entonces cada estado que satisface $p(n + 1)$ posee algún estado predecesor que satisface $p(n)$.
Repitiendo este razonamiento,
puede reconstruirse una traza finita hasta un estado inicial descrito por $p(0)$.
De esta forma, la regla permite obtener:
$$\vdash_{\mathrm{IL}} \itriple{p(0)}{C^\star}{ok : \exists n. p(n) \land \operatorname{nat}(n)}.$$

La denominación ``variante hacia atrás'' se refiere a esta reconstrucción de los estados predecesores
exigida por la interpretación subaproximada de las tripletas.
Esta caracterización también puede combinarse con la regla de \textsc{Kleene non-zero} para describir una ejecución
que realiza una cantidad arbitraria de iteraciones normales y luego termina con un error.

\subsection*{Teorema de Soundness de la IL}
La relación entre las reglas deductivas y la semántica relacional se establece mediante el Teorema de Soundness.
Este resultado garantiza que toda tripleta derivable en la IL es semánticamente válida.

\begin{teorema}[Teorema de Soundness de la IL]
\label{teo:sound-il}
Para todo comando $C$, aserciones $p, q \subseteq \Sigma$ y modo de terminación $\varepsilon$:
$$\vdash_{\mathrm{IL}} \itriple{p}{C}{\varepsilon : q} \implies \models_{\mathrm{IL}} \itriple{p}{C}{\varepsilon : q}.$$
\end{teorema}

El teorema garantiza que toda tripleta construida mediante las reglas del sistema
describe resultados respaldados por ejecuciones reales.
La demostración mecanizada de este resultado para el lenguaje imperativo
utilizado en esta tesina se presenta en el \autoref{cap:formalizacion-il}.

\begin{corolario}
\label{cor:errores}
Si
$\vdash_{\mathrm{IL}} \itriple{p}{C}{er : q}$
y $q \neq \varnothing$,
entonces existen estados $\sigma_p \in p$ y $\sigma_q \in q$ tales que:
$$(\sigma_p, \sigma_q) \in [\![C]\!]_{er}.$$
\end{corolario}
\emph{Demostración.} Por el Teorema de Soundness~\ref{teo:sound-il},
la tripleta $\itriple{p}{C}{er : q}$ es semánticamente válida.
Como $q \neq \varnothing$,
existe un estado $\sigma_q \in q$.
Aplicando el Lema~\ref{lem:caracterizacion},
existe un estado $\sigma_p \in p$ tal que
$$(\sigma_p, \sigma_q) \in [\![C]\!]_{er}.$$
Por lo tanto,
existe una ejecución de $C$ que comienza en un estado que satisface $p$ y termina con el modo $er$.
$\ensuremath{\blacksquare}$

El corolario explicita,
para el caso de terminación con error,
la garantía de ausencia de falsos positivos que se desprende del Teorema de Soundness~\ref{teo:sound-il}:
todo estado descrito por la postcondición de una tripleta de error derivable
está respaldado por una ejecución real.
Además, si $q \neq \varnothing$,
existe al menos una ejecución que comienza en un estado que satisface $p$ y termina con el modo $er$.

Esta garantía es relativa a la presunción $p$.
Si el comando analizado forma parte de un programa mayor,
el corolario no establece por sí solo que ese programa alcance alguno de los estados iniciales desde los cuales se produce el error.

\subsection*{Relevancia práctica}
Una aplicación del razonamiento sobre incorrectitud es Pulse-X~\citep{bugs},
un analizador automático de errores de memoria basado en la ISL.
En su evaluación sobre OpenSSL,
un proyecto de código abierto que proporciona bibliotecas y herramientas criptográficas para comunicaciones seguras,
encontró quince errores previamente desconocidos, que fueron reportados a sus mantenedores y corregidos.
Este resultado muestra que la IL no solo proporciona un marco teórico para razonar sobre errores alcanzables,
sino que también puede servir como fundamento para diseñar herramientas capaces de detectar errores reales.

Además,
los autores evaluaron la escalabilidad de Pulse-X sobre diez programas de gran tamaño,
midiendo el tiempo de CPU y la cantidad de problemas reportados.
Los resultados mostraron un rendimiento comparable al de Infer~\citep{infer},
un analizador utilizado en la industria.
Esta evaluación aporta evidencia de la viabilidad del enfoque para analizar bases de código extensas.

Sin embargo,
para aplicar este tipo de razonamiento a programas con memoria dinámica y estructuras enlazadas,
Pulse-X incorpora mecanismos de razonamiento local.
Estos mecanismos son provistos por la Lógica de Separación,
que se presenta a continuación.

\section{Lógica de Separación (SL)}
\label{sec:logsep}
La HL es efectiva para razonar sobre variables simples,
pero su aplicación resulta más compleja cuando los programas manipulan memoria dinámica mediante punteros.
En estos casos,
no basta con describir los valores de las variables;
también es necesario representar las celdas de memoria asignadas y la información almacenada en ellas.

La SL, introducida por \citet{sl},
extiende la HL para
facilitar el razonamiento composicional sobre este tipo de programas.
Aunque la HL permite verificar programas con memoria dinámica,
exige expresar explícitamente condiciones que indiquen,
por ejemplo,
que dos punteros se refieren a celdas distintas.
Estas condiciones se vuelven difíciles de manejar cuando se trabaja con estructuras de datos enlazadas.
La SL permite describir
el fragmento de memoria relevante para cada comando
y combinar ese razonamiento con información sobre otros fragmentos independientes,
sin tener que describir explícitamente toda la memoria del programa~\citep{local}.

Para formalizar este modelo,
sean $\mathsf{Var}$ el conjunto de variables del programa,
$\mathsf{Val}$ el conjunto de valores y $\mathsf{Loc}$ el conjunto de direcciones de memoria.
Estas direcciones también son valores,
por lo que $\mathsf{Loc} \subseteq \mathsf{Val}$.
Por su parte,
un valor distinguido \texttt{null} representa el puntero nulo y no pertenece a $\mathsf{Loc}$.

Luego,
un \emph{store} es una función total que asigna un valor a cada variable:
$$\mathsf{Store} = \mathsf{Var} \rightarrow \mathsf{Val}.$$
Por su parte, un \emph{heap} se representa mediante una función parcial finita:
$$\mathsf{Heap} = \mathsf{Loc} \rightharpoonup \mathsf{Val}.$$
La parcialidad expresa que el \emph{heap} solo asigna valores a las direcciones
que forman parte de la memoria actualmente asignada.
Denotamos por $\operatorname{dom}(h)$ el conjunto de direcciones para las cuales el \emph{heap} $h$ está definido.

Un estado espacial es un par $(s, h)$,
donde $s \in \mathsf{Store}$ y $h \in \mathsf{Heap}$.
Las aserciones de la SL se interpretan sobre estados espaciales.
Formalmente, una aserción espacial es un predicado
$$P : \mathsf{Store} \times \mathsf{Heap} \rightarrow \mathsf{Prop}.$$
Escribimos $P(s, h)$
cuando el estado formado por el \emph{store} $s$ y el \emph{heap} $h$ satisface la aserción $P$.
Las aserciones que solo dependen del \emph{store} se denominan \emph{puras}.
Por ejemplo,
una igualdad entre expresiones exige que sus valores en el \emph{store} coincidan,
sin imponer restricciones sobre el \emph{heap}.

En presencia de memoria dinámica y punteros,
una dificultad central es determinar cuándo dos expresiones del programa
se refieren a la misma dirección de memoria.
A esta situación se la denomina \emph{aliasing}.
En este caso,
una modificación realizada a través de una de las expresiones puede afectar el valor observado a través de la otra.

\begin{ejemplo}
Consideremos el comando $C \equiv [x] := 2$,
que escribe el valor 2 en la celda cuya dirección está almacenada en $x$.
Saber que las celdas referenciadas por $x$ e $y$ contienen inicialmente el valor 1
no basta para garantizar que la segunda conserve su contenido.
En particular,
si $s \in \mathsf{Store}$ y $h \in \mathsf{Heap}$,
la siguiente tripleta no es válida:
$$\triple{h(s(x)) = 1 \land h(s(y)) = 1}{[x] := 2}{h(s(x)) = 2 \land h(s(y)) = 1}.$$

Para observar por qué falla,
supongamos que ambas variables apuntan a una misma dirección $\ell$:
$$s(x) = s(y) = \ell, \qquad h(\ell) = 1.$$
Este estado satisface la precondición de la tripleta.
Sin embargo,
después de la escritura se obtiene un \emph{heap} $h'$ tal que
$$h'(\ell) = h'(s(x)) = h'(s(y)) = 2,$$
por lo que la postcondición no se cumple.
Aunque la escritura se realiza mediante $x$,
también cambia el valor que se obtiene al leer mediante $y$.

Si se agrega a la precondición que ambas direcciones son distintas,
se obtiene una tripleta válida:
$$\triple{h(s(x)) = 1 \land h(s(y)) = 1 \land s(x) \neq s(y)}{[x] := 2}{h(s(x)) = 2 \land h(s(y)) = 1}.$$
En este caso,
la escritura modifica únicamente la celda referenciada por $x$,
por lo que la celda referenciada por $y$ conserva su contenido.
\end{ejemplo}

El ejemplo muestra la necesidad de establecer qué referencias corresponden a celdas distintas.
En programas que manipulan estructuras de datos enlazadas,
estas condiciones pueden ser difíciles de construir y de demostrar automáticamente,
lo que también puede limitar la escalabilidad del análisis~\citep{bttf}.

\subsection*{La conjunción separadora y aserciones espaciales}
Una contribución central de la SL es la conjunción separadora,
un conectivo que permite combinar aserciones sobre fragmentos disjuntos de memoria.
De este modo,
la separación se expresa mediante la estructura de las aserciones,
sin enumerar individualmente las desigualdades entre las direcciones involucradas.
Para definir este conectivo,
primero introducimos
cuándo dos \emph{heaps} son disjuntos y cómo pueden combinarse.

\begin{definicion}[\emph{Heaps} disjuntos]
\label{def:disj_heaps}
Sean $h_1, h_2 \in \mathsf{Heap}$.
Decimos que $h_1$ y $h_2$ son disjuntos,
lo que denotamos mediante $h_1 \perp h_2$,
si sus dominios no comparten ninguna dirección:
$$\operatorname{dom}(h_1) \cap \operatorname{dom}(h_2) = \varnothing.$$
\end{definicion}

Que los \emph{heaps} sean disjuntos garantiza que cada dirección de memoria pertenece, a lo sumo,
a uno de los dos \emph{heaps}.
Por lo tanto,
sus contenidos pueden combinarse en un único \emph{heap} sin generar conflictos.
Esta operación se denomina \emph{unión disjunta}.

\begin{definicion}[Unión disjunta]
Sean $h_1, h_2 \in \mathsf{Heap}$ dos \emph{heaps} disjuntos.
Su unión disjunta, denotada por $h_1 \uplus h_2$,
es el \emph{heap} cuyo dominio es
$$\operatorname{dom}(h_1 \uplus h_2) = \operatorname{dom}(h_1) \cup \operatorname{dom}(h_2)$$
y cuyo valor en cada dirección de memoria
$\ell \in \operatorname{dom}(h_1 \uplus h_2)$ se define por casos:
$$(h_1 \uplus h_2)(\ell) =
\begin{cases}
h_1(\ell), & \text{si } \ell \in \operatorname{dom}(h_1), \\
h_2(\ell), & \text{si } \ell \in \operatorname{dom}(h_2).
\end{cases}$$
\end{definicion}

Sobre esta estructura de \emph{heaps},
la SL incorpora aserciones específicas para describir fragmentos de memoria.
Entre las más importantes se encuentran $\mathsf{emp}$, \emph{points-to} y la conjunción separadora.

\begin{definicion}[\emph{Heap} vacío]
\label{def:heap_vacio}
Sean $s \in \mathsf{Store}$ y $h \in \mathsf{Heap}$.
La aserción $\mathsf{emp}$ se satisface cuando $h$ no contiene ninguna celda:
$$\mathsf{emp}(s,h) \iff \operatorname{dom}(h) = \varnothing.$$
\end{definicion}

La aserción $\mathsf{emp}$ no afirma que toda la memoria física del programa se encuentre vacía,
sino que el fragmento del \emph{heap} descrito por la aserción no contiene ninguna celda.
A continuación,
se introduce la aserción \emph{points-to},
que permite describir celdas individuales del \emph{heap}.

\begin{definicion}[Points-to]
\label{def:points_to}
Sea $s \in \mathsf{Store}$ y sean $e_1$ y $e_2$ expresiones del programa que se evalúan en $s$.
La aserción
$$e_1 \mapsto e_2$$
se satisface cuando $h \in \mathsf{Heap}$ contiene exactamente una celda,
cuya dirección es el resultado de evaluar $e_1$ en $s$ y cuyo contenido es el resultado de evaluar $e_2$ en $s$.
Formalmente:
$$(e_1 \mapsto e_2)(s, h) \iff \exists \ell \in \mathsf{Loc}. \
\ell = [\![e_1]\!]_s \land \operatorname{dom}(h) = \{\ell\} \land h(\ell) = [\![e_2]\!]_s.$$
\end{definicion}

Por ejemplo,
para una variable $x$ y un \emph{store} $s$,
la aserción $x \mapsto 3$ describe un \emph{heap} formado exactamente por una celda ubicada en la dirección $s(x)$,
cuyo contenido es 3.
Esta aserción no describe un \emph{heap} arbitrariamente grande que contiene esa celda,
sino un fragmento compuesto únicamente por ella,
a veces llamado \emph{heaplet}.
Para describir otras celdas independientes pueden agregarse nuevas aserciones mediante la conjunción separadora.

\begin{definicion}[Conjunción separadora]
\label{def:sep-conj}
Sean $P$ y $Q$ dos aserciones espaciales, $s \in \mathsf{Store}$ y $h \in \mathsf{Heap}$.
El estado $(s, h)$ satisface la conjunción separadora
$$P * Q$$
si $h$ puede dividirse en dos fragmentos disjuntos $h_1$ y $h_2$, de modo que $P$ se satisface sobre $h_1$ y $Q$ se satisface sobre $h_2$.
Formalmente:
$$(P * Q)(s, h) \iff \exists h_1, h_2 \in \mathsf{Heap}. \; h_1 \perp h_2 \land h = h_1 \uplus h_2 \land P(s, h_1) \land Q(s, h_2).$$
\end{definicion}

En la fórmula anterior,
tanto $P$ como $Q$ se comprueban usando el mismo \emph{store}.
En cambio, cada aserción se comprueba sobre un fragmento diferente del \emph{heap}:
$P$ sobre $h_1$ y $Q$ sobre $h_2$.
Esto se debe a que el \emph{store} contiene la asociación entre las variables del programa y sus valores,
por lo que funciona como un entorno común para evaluar las expresiones de ambas aserciones.
Por lo tanto, la separación se aplica solo al \emph{heap},
cuyas celdas se distribuyen entre fragmentos disjuntos.

La conjunción separadora satisface las siguientes propiedades algebraicas,
que permiten reorganizar las aserciones sin cambiar su significado.

\begin{lema}[Propiedades de la conjunción separadora]
Para todas las aserciones espaciales $P$, $Q$ y $R$:
\begin{enumerate}
  \item $P * Q \iff Q * P$.
  \item $(P * Q) * R \iff P * (Q * R)$.
  \item $P * \mathsf{emp} \iff P$.
\end{enumerate}
\end{lema}

La conjunción separadora también permite definir una variante de \emph{points-to} que admite memoria adicional.
Para ello,
utilizamos la aserción $\mathsf{true}$,
que se satisface en cualquier estado espacial,
sin imponer restricciones sobre el \emph{store} ni sobre el \emph{heap}.

A partir de esta aserción y la conjunción separadora,
definimos la siguiente variante de \emph{points-to}:
$$e_1 \hookrightarrow e_2 \equiv (e_1 \mapsto e_2) * \mathsf{true}.$$
Esta aserción indica que el \emph{heap} contiene una celda cuya dirección es
el resultado de evaluar $e_1$ y cuyo contenido es el resultado de evaluar $e_2$,
pero permite que existan otras celdas.

Las dos variantes de \emph{points-to} permiten ilustrar cómo la conjunción clásica ($\land$) y la separadora ($*$)
expresan distintas formas de compartir o separar memoria.
Para presentar los ejemplos,
utilizamos la abreviatura 
$$x \mapsto [e_1, e_2] \equiv (x \mapsto e_1) * (x + 1 \mapsto e_2),$$
que describe dos celdas consecutivas cuyos contenidos son los valores de las expresiones $e_1$ y $e_2$,
respectivamente.

En los ejemplos siguientes,
consideramos las aserciones $x \mapsto [3, y]$ e $y \mapsto [3, x]$.
Ambas aserciones se representan individualmente en la \autoref{fig:aserciones}.

\begin{figure}[h]
  \centering
  \begin{minipage}{0.45\textwidth}
    \centering
    \begin{tikzpicture}[
      cell/.style={rectangle, draw=black, thick, minimum width=0.8cm, minimum height=0.8cm, inner sep=0pt, outer sep=0pt},
      ptr/.style={->, thick, >=stealth}
    ]
      \node (x) at (-1.2, 0) {$x$};
      \node[cell] (c1) at (0, 0) {3};
      \node[cell, below=-\pgflinewidth of c1] (c2) {$y$};
      \draw[ptr] (x) -- (c1.west);
    \end{tikzpicture}
  \end{minipage}
  \hfill
  \begin{minipage}{0.45\textwidth}
    \centering
    \begin{tikzpicture}[
      cell/.style={rectangle, draw=black, thick, minimum width=0.8cm, minimum height=0.8cm, inner sep=0pt, outer sep=0pt},
      ptr/.style={->, thick, >=stealth}
    ]
      \node (y) at (-1.2, 0) {$y$};
      \node[cell] (c1) at (0, 0) {3};
      \node[cell, below=-\pgflinewidth of c1] (c2) {$x$};
      \draw[ptr] (y) -- (c1.west);
    \end{tikzpicture}
  \end{minipage}
\caption{Representación individual de las aserciones espaciales $x \mapsto [3, y]$ e $y \mapsto [3, x]$.}
\label{fig:aserciones}
\end{figure}

Consideremos primero la aserción $(x \mapsto [3, y]) \land (y \mapsto [3, x])$.
La conjunción clásica exige comprobar ambas aserciones sobre el mismo \emph{heap},
sin dividirlo.
La primera exige que ese \emph{heap} contenga exactamente el bloque que comienza en la dirección almacenada en $x$;
la segunda exige que contenga exactamente el bloque que comienza en la dirección almacenada en $y$.
Por ello,
ambas solo pueden cumplirse si las dos variables apuntan al mismo bloque,
como se muestra en la \autoref{fig:conjuncion-clasica}.

\begin{figure}[h]
  \centering
  \begin{tikzpicture}[
    cell/.style={rectangle, draw=black, thick, minimum width=0.8cm, minimum height=0.8cm, inner sep=0pt, outer sep=0pt},
    ptr/.style={->, thick, >=stealth}
  ]
    \node (x) at (-1.2, 0.3) {$x$};
    \node (y) at (-1.2, -0.3) {$y$};
    \node[cell] (c1) at (0, 0) {3};
    \node[cell, below=-\pgflinewidth of c1] (c2) {$\circ$};
    \draw[ptr] (x) -- ([yshift=0.03cm]c1.west);
    \draw[ptr] (y) -- ([yshift=-0.03cm]c1.west);
    \draw[ptr] ([xshift=0.1cm]c2.center) to[out=0, in=0, looseness=2] (c1.east);
  \end{tikzpicture}
\caption{Interpretación con aliasing: $x$ e $y$ apuntan al mismo bloque.}
\label{fig:conjuncion-clasica}
\end{figure}

En cambio,
la aserción $(x \mapsto [3, y]) * (y \mapsto [3, x])$ exige dividir el \emph{heap} en dos fragmentos disjuntos.
Sobre uno se comprueba $x \mapsto [3, y]$ y sobre el otro, $y \mapsto [3, x]$.
Cada fragmento contiene exactamente las dos celdas del bloque correspondiente.
Así se obtienen los dos bloques separados de la \autoref{fig:conjuncion-separadora},
cada uno de los cuales almacena un puntero hacia el otro.

\begin{figure}[h]
  \centering
  \begin{tikzpicture}[
    cell/.style={rectangle, draw=black, thick, minimum width=0.8cm, minimum height=0.8cm, inner sep=0pt, outer sep=0pt},
    ptr/.style={->, thick, >=stealth}
  ]
    \node (x) at (-2.0, 0) {$x$};
    \node (y) at (2.0, 0) {$y$};
    \node[cell] (c1) at (-1.0, 0) {3};
    \node[cell, below=-\pgflinewidth of c1] (c2) {$\circ$};
    \node[cell] (c3) at (1.0, 0) {3};
    \node[cell, below=-\pgflinewidth of c3] (c4) {$\circ$};
    \draw[ptr] (x) -- (c1.west);
    \draw[ptr] (y) -- (c3.east);
    \draw[ptr] ([xshift=0.1cm]c2.center) to[out=0, in=180, looseness=1.2] (c3.west);
    \draw[ptr] ([xshift=-0.1cm]c4.center) to[out=180, in=0, looseness=1.2] (c1.east);
  \end{tikzpicture}
\caption{Interpretación con bloques disjuntos: $x$ e $y$ apuntan a bloques separados.}
\label{fig:conjuncion-separadora}
\end{figure}

Por último,
consideremos qué ocurre al utilizar la variante que permite memoria adicional.
Extendemos la notación de la siguiente manera:
$$x \hookrightarrow [e_1, e_2] \equiv (x \mapsto [e_1, e_2]) * \mathsf{true}.$$
Con esta notación,
la aserción $(x \hookrightarrow [3, y]) \land (y \hookrightarrow [3, x])$
también comprueba ambos predicados sobre el mismo \emph{heap}.
Sin embargo,
ahora cada uno exige únicamente que el \emph{heap} contenga el bloque correspondiente.
Si ambos bloques coinciden,
se obtiene la configuración de la \autoref{fig:conjuncion-clasica}.
Si son disjuntos,
se obtiene la de la \autoref{fig:conjuncion-separadora}:
al comprobar cada predicado,
el otro bloque forma parte de la memoria adicional admitida por $\mathsf{true}$.
En ambos casos,
el \emph{heap} puede contener además otras celdas.

Si combinamos estas aserciones mediante la conjunción separadora,
$(x \hookrightarrow [3, y]) * (y \hookrightarrow [3, x])$,
los bloques deben pertenecer a fragmentos disjuntos del \emph{heap}.
Por lo tanto,
se obtiene la configuración de la \autoref{fig:conjuncion-separadora},
aunque pueden existir otras celdas además de las representadas.
La \autoref{tab:comparacion-sl} resume las diferencias entre las cuatro combinaciones.

\begin{table}[ht]
  \centering
  \small
  \renewcommand{\arraystretch}{1.3}
  \begin{tabular}{@{}cccc@{}}
    \toprule
    Aserción
    & Relación entre bloques
    & \begin{tabular}[c]{@{}c@{}}
        ¿Admite memoria\\
        adicional?
      \end{tabular}
    & \begin{tabular}[c]{@{}c@{}}
        Figuras\\
        ilustrativas
      \end{tabular} \\
    \midrule

    $(x \mapsto [3, y]) \land (y \mapsto [3, x])$
    & Deben coincidir
    & \xmark
    & \ref{fig:conjuncion-clasica} \\

    $(x \mapsto [3, y]) * (y \mapsto [3, x])$
    & Deben ser disjuntos
    & \xmark
    & \ref{fig:conjuncion-separadora} \\

    $(x \hookrightarrow [3, y]) \land (y \hookrightarrow [3, x])$
    & Pueden compartir celdas o ser disjuntos
    & \cmark
    & \ref{fig:conjuncion-clasica} y \ref{fig:conjuncion-separadora} \\

    $(x \hookrightarrow [3, y]) * (y \hookrightarrow [3, x])$
    & Deben ser disjuntos
    & \cmark
    & \ref{fig:conjuncion-separadora} \\

    \bottomrule
  \end{tabular}
  \caption{Comparación de las variantes de \emph{points-to} y las conjunciones.}
  \label{tab:comparacion-sl}
\end{table}

En la tercera fila de la \autoref{tab:comparacion-sl},
las aserciones pueden compartir celdas de distintas maneras;
las Figuras \ref{fig:conjuncion-clasica} y \ref{fig:conjuncion-separadora} ilustran algunos casos,
pero no todos.
Por ejemplo,
con $x = 2$, $y = 3$ y $h = \{2 \mapsto 3, 3 \mapsto 3, 4 \mapsto 2\}$,
ambas aserciones se satisfacen y comparten solo la celda en la posición 3.

La separación de los fragmentos del \emph{heap} impuesta por la conjunción separadora
permite razonar sobre las modificaciones de uno de ellos
y conservar la información sobre los demás.
La Regla del Marco, presentada más adelante,
formaliza este razonamiento bajo las condiciones necesarias para que la información conservada no resulte afectada por el comando.

\subsection*{El lenguaje imperativo extendido}
Las aserciones anteriores permiten describir fragmentos del \emph{heap},
pero todavía es necesario extender el lenguaje de comandos para poder operar sobre esa memoria.
Para ello, la SL incorpora cuatro comandos fundamentales para la manipulación de estructuras de datos mutables.
En estos comandos,
$x$ es una variable del programa,
mientras que $e$, $e_1$ y $e_2$ son expresiones.
Las expresiones utilizadas para acceder al \emph{heap} se evalúan en el \emph{store} para obtener la dirección de memoria correspondiente.
\begin{enumerate}
  \item \textbf{Alocación o reserva de memoria (allocation):}
  $$x := \texttt{alloc()}$$
  reserva una nueva celda y almacena su dirección en $x$.
  \item \textbf{Lectura (lookup):}
  $$x := [e]$$
  evalúa la expresión $e$ para obtener una dirección de memoria,
  lee el contenido de la celda ubicada en esa dirección y lo asigna a $x$.
  \item \textbf{Escritura o mutación (mutation):}
  $$[e_1] := e_2$$
  evalúa la expresión $e_1$ para obtener una dirección de memoria y
  sobrescribe el contenido de esa celda con el resultado de evaluar $e_2$.
  \item \textbf{Liberación (deallocation):}
  $$\texttt{free(}e\texttt{)}$$
  libera la celda ubicada en la dirección que se obtiene de evaluar $e$.
\end{enumerate}

\subsection*{Semántica relacional del lenguaje extendido}
La semántica del lenguaje extendido se interpreta sobre estados espaciales de la forma $(s,h)$,
con $s \in \mathsf{Store}$ y $h \in \mathsf{Heap}$.
Los comandos básicos presentados para la IL conservan su comportamiento y dejan el \emph{heap} sin modificaciones.
Asimismo,
las definiciones de composición secuencial,
elección no determinista e iteración mantienen la estructura vista en la \autoref{fig:semantica-il},
interpretadas ahora sobre estos estados.
La \autoref{fig:semantica-sl} presenta la semántica de los comandos que manipulan la memoria.

La notación $h[\ell \mapsto v]$ representa la actualización del \emph{heap} en $\ell$,
incorporando esa dirección si no pertenecía al dominio.
Por su parte,
$h \setminus \{\ell\}$ elimina $\ell$ del dominio y conserva las demás celdas.

\begin{figure}[htbp]
\centering
\begin{minipage}{0.98\textwidth}
\small
$$
\begin{aligned}
\relax[\![x := \texttt{alloc()}]\!]_{ok} &= \{((s, h), (s[x \mapsto \ell], h[\ell \mapsto v]))
  \mid v \in \mathsf{Val} \land \ell \in \mathsf{Loc} \setminus \operatorname{dom}(h)\}, \\[4pt]
[\![x := \texttt{alloc()}]\!]_{er} &= \varnothing, \\[4pt]
[\![\texttt{free(}e\texttt{)}]\!]_{ok} &= \{((s, h), (s, h \setminus \{\ell\}))
  \mid \ell = [\![e]\!]_s \land \ell \in \operatorname{dom}(h)\}, \\[4pt]
[\![\texttt{free(}e\texttt{)}]\!]_{er} &= \{(\sigma, \sigma)
  \mid \sigma = (s, h) \land [\![e]\!]_s \notin \operatorname{dom}(h)\}, \\[4pt]
[\![x := [e]]\!]_{ok} &= \{((s, h), (s[x \mapsto h(\ell)], h))
  \mid \ell = [\![e]\!]_s \land \ell \in \operatorname{dom}(h)\}, \\[4pt]
[\![x := [e]]\!]_{er} &= \{(\sigma, \sigma) \mid \sigma = (s, h) \land [\![e]\!]_s \notin \operatorname{dom}(h)\}, \\[4pt]
[\![[e_1] := e_2]\!]_{ok} &= \{((s, h), (s, h[\ell \mapsto [\![e_2]\!]_s]))
  \mid \ell = [\![e_1]\!]_s \land \ell \in \operatorname{dom}(h)\}, \\[4pt]
[\![[e_1] := e_2]\!]_{er} &= \{(\sigma, \sigma)
  \mid \sigma = (s, h) \land [\![e_1]\!]_s \notin \operatorname{dom}(h)\}.
\end{aligned}
$$
\vspace{0.1em}
\end{minipage}
\rule{\linewidth}{0.4pt}
\caption{Semántica relacional del lenguaje extendido.}
\label{fig:semantica-sl}
\end{figure}

Como vimos antes,
en el \emph{heap},
el dominio contiene las direcciones de las celdas actualmente asignadas.
Por ello,
las operaciones de lectura,
escritura y liberación requieren que la dirección obtenida al evaluar la expresión pertenezca a dicho dominio;
en caso contrario,
terminan con un error.
Una liberación exitosa elimina la celda del \emph{heap},
por lo que este modelo no distingue una dirección previamente liberada de otra que tampoco pertenece al dominio.

\subsection*{Reglas de inferencia de la SL}
Una vez extendido el lenguaje con comandos para manipular la memoria,
es necesario introducir las reglas de inferencia que describen sus efectos sobre las aserciones espaciales.
Las tripletas de la SL conservan la forma $\triple{P}{C}{Q}$,
pero ahora $P$ y $Q$ son aserciones espaciales interpretadas sobre estados compuestos por un \emph{store} y un \emph{heap}.

Antes de presentar las reglas,
definimos la noción de derivabilidad en la SL.

\begin{definicion}[Derivabilidad en la SL]
\label{def:derivabilidad-sl}
Escribimos
$$\vdash_{\mathrm{SL}} \triple{P}{C}{Q}$$
cuando la tripleta puede obtenerse mediante un número finito de aplicaciones de las reglas de inferencia de la SL.
\end{definicion}

Las reglas correspondientes a comandos como \texttt{skip}, asignación y composición secuencial se heredan de la HL.
La extensión propia de la SL está dada por las reglas asociadas a las operaciones de alocación,
lectura, escritura y liberación de memoria en la \autoref{fig:reglas-sl}.
En estas reglas se utiliza la abreviatura $e \mapsto -$,
equivalente a $\exists v. \ e \mapsto v$,
que representa una celda ubicada en la dirección obtenida al evaluar $e$,
cuyo contenido no es relevante para el razonamiento.

\begin{figure}[htbp]
\centering
\begin{mathpar}
\inferrule [Alloc]
  {}
  {\vdash_{\mathrm{SL}} \triple{\mathsf{emp}}{x := \texttt{alloc()}}{x \mapsto -}}

\inferrule [Free]
  {}
  {\vdash_{\mathrm{SL}} \triple{e \mapsto -}{\texttt{free(}e\texttt{)}}{\mathsf{emp}}}

\inferrule [Load]
  {}
  {\vdash_{\mathrm{SL}} \triple{x = x' \land e \mapsto v}{x := [e]}{x = v[x'/x] \land e[x'/x] \mapsto v[x'/x]}}

\inferrule [Store]
  {}
  {\vdash_{\mathrm{SL}} \triple{e_1 \mapsto -}{[e_1] := e_2}{e_1 \mapsto e_2}}

\end{mathpar}
\caption{Reglas de inferencia de la SL. Adaptadas de~\citep{sl}.}
\label{fig:reglas-sl}
\end{figure}

La regla \textsc{Alloc} expresa que la alocación reserva una nueva celda y almacena su dirección en $x$.
Por su parte,
\textsc{Free} elimina del fragmento de \emph{heap} descrito la celda ubicada en la dirección obtenida al evaluar $e$
y produce la aserción $\mathsf{emp}$.

La regla \textsc{Load}
lee el contenido de la celda indicada por $e$ y lo asigna a $x$,
sin modificar el \emph{heap}.
La variable auxiliar $x'$,
elegida distinta de $x$,
conserva el valor inicial de $x$,
mientras que $v$ representa el contenido leído.
Como la asignación puede cambiar el valor de la expresión $e$ si esta contiene a $x$,
la postcondición utiliza $e[x'/x]$ para seguir describiendo la dirección original.
Por el mismo motivo el contenido leído se describe mediante $v[x'/x]$:
si $v$ contiene a $x$,
su valor debe evaluarse con el valor de $x$ previo a la lectura.

La regla \textsc{Store} modifica el contenido de la celda
indicada por $e_1$ con el resultado de evaluar $e_2$.
Esta operación modifica el \emph{heap},
pero conserva el \emph{store}.

Estas reglas describen el comportamiento de cada operación sobre el fragmento de memoria que aparece explícitamente en la precondición.
Para extender este razonamiento a estados que contienen memoria adicional e independiente,
se introduce el principio de localidad,
formalizado mediante la Regla del Marco.

\subsection*{Razonamiento local y la Regla del Marco}
La posibilidad de describir fragmentos disjuntos del \emph{heap}
permite formular el principio de razonamiento local.
Este principio establece que,
para razonar sobre un comando,
basta considerar las celdas de memoria que el comando puede utilizar.
Las celdas adicionales, siempre que sean disjuntas del fragmento considerado,
no necesitan describirse.
Si el comando no las utiliza, permanecen inalteradas durante la ejecución.

\begin{ejemplo}
La tripleta
$$\triple{x \mapsto v}{[x] := w}{x \mapsto w}$$
describe localmente una escritura en memoria.
La precondición contiene solo la celda a la que se accede,
y la postcondición describe su nuevo valor.

Supongamos ahora que el programa se ejecuta en un \emph{heap} más grande
que contiene además el fragmento independiente descrito por $y \mapsto u$.
La escritura mediante $x$ no debería modificar esta segunda celda.
Por lo tanto, debería ser posible obtener:
$$\triple{x \mapsto v * y \mapsto u}{[x] := w}{x \mapsto w * y \mapsto u}$$
\end{ejemplo}

Este principio se materializa mediante la Regla del Marco (\textit{Frame Rule}),
que permite razonar sobre la parte de la memoria que interviene en la ejecución,
sin necesidad de describir el resto del \emph{heap}.

Sean $P$, $Q$ y $R$ aserciones espaciales y sea $C$ un comando.
Definimos la Regla del Marco como:
\begin{mathpar}
\inferrule []
  {\vdash_{\mathrm{SL}} \triple{P}{C}{Q} \\ \operatorname{mod}(C) \cap \operatorname{fv}(R) = \varnothing}
  {\vdash_{\mathrm{SL}} \triple{P * R}{C}{Q * R}}
\end{mathpar}

En la condición lateral tenemos $\operatorname{mod}(C)$,
el conjunto de variables del \emph{store} que pueden ser modificadas por $C$,
y $\operatorname{fv}(R)$, el conjunto de variables libres de la aserción $R$.
Esta condición garantiza que $C$ no modifique ninguna variable de las que depende $R$,
de modo que $R$ siga describiendo el mismo fragmento de memoria después de ejecutar $C$.
Por ejemplo,
sean $R \equiv x \mapsto 1$ y $C \equiv x := 3$,
donde $x$ es una variable del \emph{store}.
La aserción $R$ describe una celda del \emph{heap} ubicada en la dirección almacenada en $x$,
mientras que el comando modifica el valor de la propia variable.
Por lo tanto,
aunque el \emph{heap} permanezca intacto,
$R$ puede dejar de cumplirse porque $x$ ya no tiene la dirección de la celda descrita.
La condición lateral impide esta aplicación de la Regla del Marco,
ya que $x \in \operatorname{fv}(R)$ y $x \in \operatorname{mod}(C)$.

La SL proporciona las herramientas necesarias para razonar localmente sobre memoria mutable.
A continuación,
consideramos cómo combinar este razonamiento con la interpretación subaproximada de la IL.

\section{Lógica de Incorrectitud de Separación (ISL)}
\label{sec:logincsep}
La ISL,
introducida por \citet{isl}, combina los dos enfoques presentados anteriormente.
De la IL toma la interpretación subaproximada,
que permite respaldar los estados finales mediante ejecuciones reales.
De la SL toma la representación explícita del \emph{heap} y el razonamiento local mediante la conjunción separadora y la Regla del Marco.

Las tripletas de la ISL conservan la forma presentada para la IL:
$$\itriple{p}{C}{\varepsilon : q},$$
pero ahora la presunción $p$ y la postcondición $q$ son aserciones espaciales que describen conjuntos de estados compuestos por un \emph{store} y un \emph{heap}.
Su interpretación semántica sigue siendo subaproximada:
todo estado descrito por $q$ debe estar respaldado por una ejecución real de $C$ desde algún estado que satisface $p$,
con el modo de terminación $\varepsilon$.

Al igual que en las lógicas anteriores,
distinguimos esta interpretación semántica de la derivabilidad de una tripleta mediante reglas de inferencia.

\begin{definicion}[Derivabilidad en la ISL]
\label{def:derivabilidad-isl}
Escribimos
$$\vdash_{\mathrm{ISL}} \itriple{p}{C}{\varepsilon : q}$$
cuando la tripleta puede obtenerse mediante un número finito de aplicaciones de las reglas de inferencia de la ISL.
\end{definicion}

La construcción de este sistema deductivo requiere revisar cómo se combinan las reglas de memoria con la interpretación subaproximada.
En particular,
utilizar directamente la regla de liberación de la SL y aplicar la Regla del Marco puede producir una tripleta inválida.
El siguiente ejemplo muestra por qué esta combinación no es válida.

\subsection*{Heaps negativos}
En la semántica de la SL,
liberar una celda elimina del \emph{heap} la dirección que ocupaba.
De este modo,
el \emph{heap} deja de tener información sobre el hecho de que esa dirección fue liberada.
Esta pérdida de información resulta problemática para la interpretación subaproximada,
porque puede permitir derivar una postcondición satisfacible a partir de una presunción que ningún estado satisface.

Consideremos la siguiente tripleta sobre el modelo clásico de memoria:
$$\itriple{x \mapsto -}{\texttt{free(}x\texttt{)}}{ok : \mathsf{emp} \land \mathsf{loc}(x)}.$$
En este caso,
$\mathsf{loc}(x)$ expresa que $x$ contiene una dirección de memoria,
donde $x$ no puede ser \texttt{null}.
Esta aserción solo depende del \emph{store} y no exige que la dirección pertenezca al dominio del \emph{heap}.

Supongamos que aplicamos la Regla del Marco utilizando como marco otra celda descrita por $x \mapsto -$:
\begin{mathpar}
\inferrule []
  {\itriple{x \mapsto -}{\texttt{free(}x\texttt{)}}{ok : \mathsf{emp} \land \mathsf{loc}(x)}}
  {\itriple{x \mapsto - * x \mapsto -}{\texttt{free(}x\texttt{)}}{ok : (\mathsf{emp} \land \mathsf{loc}(x)) * x \mapsto -}}
\end{mathpar}

La presunción $x \mapsto - * x \mapsto -$ es insatisfacible porque exige que dos fragmentos disjuntos posean la misma dirección.
No existe ningún estado que la satisfaga.
Sin embargo, la postcondición $(\mathsf{emp} \land \mathsf{loc}(x)) * x \mapsto -$ es satisfacible.
Esto produciría una tripleta con una presunción inconsistente y una postcondición satisfacible.
Recordando la semántica de una tripleta de subaproximación,
y dado que la postcondición no es vacía,
esto indicaría que existe un estado que satisface la presunción,
lo cual es absurdo.

El problema no pertenece a la IL en sí misma,
sino a la combinación directa de la subaproximación,
la semántica clásica de liberación y la Regla del Marco.
El modelo clásico no permite distinguir entre una dirección sobre la cual no se posee información y una dirección que se sabe que fue liberada.
Esta distinción es la que introducen los \emph{heaps negativos}.

\begin{definicion}[\emph{Heap} extendido]
\label{def:heap_ext}
Un \emph{heap} de la ISL,
que denotamos $\mathsf{HeapExt}$,
es una función parcial finita
%\guido{por qué? hace falta? las otras versiones también son finitas?}:
$$\mathsf{HeapExt} = \mathsf{Loc} \rightharpoonup (\mathsf{Val} \cup \{\bot\}),$$
donde si $h \in \mathsf{HeapExt}$:
\begin{itemize}
  \item $h(\ell) = v$ indica que $\ell$ se encuentra asignada y contiene $v$;
  \item $h(\ell) = \bot$ indica que $\ell$ fue liberada y todavía no fue reasignada;
  \item $\ell \notin \operatorname{dom}(h)$ indica que el \emph{heap} no posee información sobre esa dirección.
\end{itemize}
\end{definicion}
Al igual que en el modelo de \emph{heaps} de la SL,
la finitud se refiere al dominio de cada \emph{heap}.
En este caso,
dicho dominio incluye tanto las celdas activas como las marcadas como liberadas.

A partir de este punto,
los estados de la ISL se interpretan sobre el \emph{heap} extendido
$$\Sigma_{\mathrm{ISL}} = \mathsf{Store} \times \mathsf{HeapExt}.$$
La relación de \emph{heaps} disjuntos y la unión disjunta se extienden de manera natural al \emph{heap} extendido,
incluyendo tanto celdas activas como celdas liberadas.

\begin{definicion}[\emph{Heap} negativo]
\label{def:heap_negativo}
La aserción:
$$e \nmapsto$$
se satisface cuando el \emph{heap} está formado por una celda liberada ubicada en la dirección indicada por $e$.
Formalmente, sean $s \in \mathsf{Store}$ y $h \in \mathsf{HeapExt}$:
$$(e \nmapsto)(s,h) \iff \exists \ell \in \mathsf{Loc}. \;
[\![e]\!]_s = \ell \land \operatorname{dom}(h) = \{\ell\} \land h(\ell) = \bot$$
\end{definicion}

La dirección obtenida al evaluar $e$ permanece en el dominio del \emph{heap},
aunque la celda esté liberada.
Por lo tanto,
la conjunción separadora impide combinar $e \nmapsto$ con otro fragmento que contenga esa misma dirección.
Esta propiedad permite conservar en la postcondición la información espacial que se perdía en el modelo clásico.
Con esta representación,
la regla local de liberación se expresa como:
$$\vdash_{\mathrm{ISL}} \itriple{x \mapsto -}{\texttt{free(}x\texttt{)}}{ok : x \nmapsto}$$

Al aplicar otra vez el marco $x \mapsto -$ obtenemos:
$$\vdash_{\mathrm{ISL}} \itriple{x \mapsto - * x \mapsto -}{\texttt{free(}x\texttt{)}}{ok : x \nmapsto * \; x \mapsto -}$$
La presunción sigue siendo inconsistente,
pero ahora la postcondición también lo es,
ya que exige que una misma dirección pertenezca a dos \emph{heaps} disjuntos,
estando liberada en uno y activa en el otro.
Por lo tanto, la conclusión es semánticamente válida:
$$\varnothing \subseteq post_{ok}(\texttt{free(}x\texttt{)})(\varnothing).$$
De este modo,
los \emph{heaps} negativos impiden que la Regla del Marco transforme una presunción inconsistente en una postcondición satisfacible.

El ejemplo anterior muestra cómo los \emph{heaps} negativos permiten conservar la inconsistencia al aplicar la Regla del Marco.
Para continuar,
es necesario describir formalmente la semántica relacional sobre el espacio de estados extendido.

\subsection*{Semántica relacional de la ISL}
La semántica presentada en la SL se adapta al espacio de estados extendido
$\Sigma_{\mathrm{ISL}} = \mathsf{Store} \times \mathsf{HeapExt}$.
Los comandos que no manipulan directamente el \emph{heap} conservan la semántica presentada para la IL,
interpretada ahora sobre estos estados.
En cambio,
las operaciones de memoria deben distinguir entre celdas activas,
celdas liberadas y las direcciones sobre las cuales el \emph{heap} no tiene información.

La \autoref{fig:semantica-isl} presenta las definiciones correspondientes a las operaciones de memoria sobre el \emph{heap} extendido.
En ellas,
la condición $h(\ell) \in \mathrm{Val}$ indica que la celda está activa,
mientras que $h(\ell) = \bot$ indica que fue liberada y todavía no fue reasignada.
Las operaciones de lectura, escritura y liberación sobre el puntero nulo \texttt{null} terminan con un error.

\begin{figure}[htbp]
\centering
\begin{minipage}{0.98\textwidth}
\small
$$
\begin{aligned}
\relax[\![x := \texttt{alloc()}]\!]_{ok} &= \{((s, h), (s[x \mapsto \ell], h[\ell \mapsto v]))
  \mid v \in \mathsf{Val} \land \ell \in \mathsf{Loc} \land (\ell \notin \operatorname{dom}(h) \lor h(\ell) = \bot)\}, \\[4pt]
[\![\texttt{free(}e\texttt{)}]\!]_{ok} &= \{((s,h), (s, h[\ell \mapsto \bot]))
  \mid \ell = [\![e]\!]_s \land h(\ell) \in \mathsf{Val}\}, \\[4pt]
[\![\texttt{free(}e\texttt{)}]\!]_{er} &= \{(\sigma, \sigma)
  \mid \sigma = (s, h) \land ([\![e]\!]_s = \texttt{null} \lor h([\![e]\!]_s) = \bot)\}, \\[4pt]
[\![x := [e]]\!]_{ok} &= \{((s, h), (s[x \mapsto v], h))
  \mid h([\![e]\!]_s) = v \in \mathsf{Val}\}, \\[4pt]
[\![x := [e]]\!]_{er} &= \{(\sigma, \sigma) \mid \sigma = (s, h) \land ([\![e]\!]_s = \texttt{null} \lor h([\![e]\!]_s) = \bot)\}, \\[4pt]
[\![[e_1] := e_2]\!]_{ok} &= \{((s, h), (s, h[\ell \mapsto [\![e_2]\!]_s]))
  \mid \ell = [\![e_1]\!]_s \land h(\ell) \in \mathsf{Val}\}, \\[4pt]
[\![[e_1] := e_2]\!]_{er} &= \{(\sigma, \sigma)
  \mid \sigma = (s, h) \land ([\![e_1]\!]_s = \texttt{null} \lor h([\![e_1]\!]_s) = \bot)\}.
\end{aligned}
$$
\vspace{0.1em}
\end{minipage}
\rule{\linewidth}{0.4pt}
\caption{Semántica de las operaciones de memoria sobre el \emph{heap} extendido.}
\label{fig:semantica-isl}
\end{figure}

A diferencia del modelo clásico,
una liberación exitosa conserva la dirección en el dominio del \emph{heap} y la marca con $\bot$.
La alocación puede reutilizar esa dirección o reservar una que no pertenezca al dominio.
Por su parte,
las terminaciones erróneas de lectura,
escritura y liberación se justifican mediante un puntero nulo o una celda marcada como liberada.

En el razonamiento subaproximado,
los resultados erróneos descritos por una postcondición deben estar respaldados por ejecuciones reales.
Para justificar estos errores localmente,
la ausencia de una dirección en el dominio del \emph{heap} no es suficiente,
ya que esa dirección podría corresponder a una celda activa de un fragmento de memoria adicional.
En cambio,
un puntero nulo no representa una dirección válida,
y una celda marcada con $\bot$ no puede combinarse mediante la conjunción separadora con otra celda en la misma dirección.
Por ello,
ambas condiciones permiten justificar errores que se conservan al agregar un fragmento disjunto de memoria.

Esta semántica considera programas secuenciales.
En un contexto concurrente,
donde otro hilo podría alocar una celda en una dirección liberada,
se debe considerar la interferencia entre hilos.
Esta extensión queda fuera del alcance del trabajo.

Sobre esta semántica secuencial,
la validez de las tripletas de la ISL se define mediante la interpretación subaproximada:
$$\models_\mathrm{ISL} \itriple{p}{C}{\varepsilon : q} \iff
\forall \sigma_q \in q. \exists \sigma_p \in p. (\sigma_p, \sigma_q) \in [\![C]\!]_\varepsilon.$$
Con esta definición,
se exige que cada estado espacial descrito por la postcondición esté respaldado por una ejecución real del comando
desde algún estado espacial que satisface la presunción.
La condición se aplica tanto a las terminaciones normales como a las terminaciones erróneas,
según el modo $\varepsilon$ considerado.

A partir de estas definiciones se obtiene una propiedad fundamental del \emph{heap} extendido.
Una ejecución puede modificar el estado asociado a una dirección,
por ejemplo,
transformando una celda activa en una celda liberada,
pero no puede eliminar del dominio una dirección que ya era conocida.
Del mismo modo,
una alocación puede transformar una dirección liberada en una celda activa sin modificar el dominio conocido.
Esta propiedad se formaliza en el siguiente lema.

\begin{lema}[Monotonía del \emph{heap}]
\label{lem:monotonia}
Para todo comando $C$,
modo de terminación $\varepsilon$ y transición $((s, h), (s', h')) \in [\![C]\!]_\varepsilon$,
se cumple:
$$\operatorname{dom}(h) \subseteq \operatorname{dom}(h').$$
\end{lema}

Esto quiere decir que una dirección conocida al comienzo de una ejecución sigue siendo conocida cuando esta finaliza.
La monotonía puede observarse en las operaciones básicas:
\begin{itemize}
  \item una alocación sobre una dirección nueva amplía el dominio;
  \item una lectura o una escritura conserva el dominio;
  \item una liberación transforma un valor $v$ en $\bot$,
  pero no elimina la dirección;
  \item una nueva alocación sobre una dirección previamente liberada transforma $\bot$ en un valor,
  pero tampoco modifica el dominio.
\end{itemize}
Por lo tanto, el dominio puede crecer o permanecer igual,
pero nunca disminuir.

Sobre esta semántica,
una ejecución puede extenderse con un fragmento de memoria disjunto de sus \emph{heaps} inicial y final,
conservando el modo de terminación.
Esta propiedad se expresa mediante el siguiente lema.

\begin{lema}[Localidad semántica]
\label{lem:localidad}
Sean $C$ un comando, $\varepsilon$ un modo de terminación,
$s$ y $s'$ los \emph{stores} inicial y final,
$h, h', h_f \in \mathsf{HeapExt}$ los \emph{heaps} inicial, final y adicional,
respectivamente.
Si
$$((s, h), (s', h')) \in [\![C]\!]_\varepsilon \land h' \perp h_f,$$
entonces
$$((s, h \uplus h_f), (s', h' \uplus h_f)) \in [\![C]\!]_\varepsilon.$$
\end{lema}
%\guido{creo que esto no es cierto. Considera
%estado inicial y final vacio, frame $x \leadsto 1$ y el comando $Read~x$.
%El programa falla y tiene modo de terminacion con error. Si agrego el frame
%ahora no falla.}
%\guido{update: si bien en la mecanizacion esto está bien,
%la presentación es distinta: para que el Load falle,
%pide que la dirección no esté en el dominio,
%y no que esté en Empty. Así que está mal.}
%\agustina{Entiendo, esto vale en la ISL pero no en la SL porque en la ISL tenemos el heap extendido con $\bot$
%y cambia la semantica para este caso. Mañana lo cambio.}

El lema establece que,
si una ejecución transforma el \emph{heap} $h$ en el \emph{heap} $h'$
y el fragmento adicional $h_f$ es disjunto de $h'$,
la ejecución puede extenderse al \emph{heap} inicial $h \uplus h_f$.
La hipótesis no exige por separado que $h_f$ sea disjunto de $h$,
ya que por la monotonicidad del dominio,
toda ubicación conocida en $h$ también está presente en $h'$.
Por lo tanto,
si $h'$ es disjunto de $h_f$,
también lo es $h$.
En ese caso,
el comando transforma $h \uplus h_f$ en $h' \uplus h_f$,
conservando el modo de terminación y dejando $h_f$ sin modificaciones.

\subsection*{Reglas de inferencia de la ISL}
Las reglas de la IL que no manipulan la memoria se mantienen en la ISL.
Las diferencias se encuentran en las reglas de alocación, lectura, escritura y liberación,
y en la aplicación de la Regla del Marco dentro de la interpretación subaproximada.
También se incluye la regla \textsc{Exist},
que permite cuantificar existencialmente una variable que no aparece en el comando.
La noción de derivabilidad utilizada en esta sección fue introducida en la Definición~\ref{def:derivabilidad-isl}.
A continuación se presentan las nuevas reglas de la ISL en la \autoref{fig:reglas-isl}.

%\guido{Alloc2 no necesita $x \notin FV(y)$? En la formalización toma una location,
%no una expresión como acá.}
\begin{figure}[htbp]
\centering
\begin{mathpar}
\inferrule [Exist]
  {\vdash_{\mathrm{ISL}} \itriple{p}{C}{\varepsilon : q} \quad x \notin \operatorname{fv}(C)}
  {\vdash_{\mathrm{ISL}} \itriple{\exists x. \ p}{C}{\varepsilon : \exists x. \ q}}

\inferrule [Alloc1]
  {}
  {\vdash_{\mathrm{ISL}} \itriple{x = x'}{x := \texttt{alloc()}}{ok : x \mapsto -}}

\inferrule [Alloc2]
  {}
  {\vdash_{\mathrm{ISL}} \itriple{x = x' * e \nmapsto}{x := \texttt{alloc()}}{ok : x = e[x'/x] * e \mapsto -}}

\inferrule [Free]
  {}
  {\vdash_{\mathrm{ISL}} \itriple{e \mapsto e'}{\texttt{free(}e\texttt{)}}{ok : e \nmapsto}}

\inferrule [FreeEr]
  {}
  {\vdash_{\mathrm{ISL}} \itriple{e \nmapsto}{\texttt{free(}e\texttt{)}}{er : e \nmapsto}}

\inferrule [FreeNull]
  {}
  {\vdash_{\mathrm{ISL}} \itriple{e = \texttt{null}}{\texttt{free(}e\texttt{)}}{er : e = \texttt{null}}}

\inferrule [Load]
  {}
  {\vdash_{\mathrm{ISL}} \itriple{x = x' * e \mapsto v}{x := [e]}{ok : x = v[x'/x] * e[x' / x] \mapsto v[x'/x]}}

\inferrule [LoadEr]
  {}
  {\vdash_{\mathrm{ISL}} \itriple{e \nmapsto}{x := [e]}{er : e \nmapsto}}

\inferrule [LoadNull]
  {}
  {\vdash_{\mathrm{ISL}} \itriple{e = \texttt{null}}{x := [e]}{er : e = \texttt{null}}}

\inferrule [Store]
  {}
  {\vdash_{\mathrm{ISL}} \itriple{e_1 \mapsto e}{[e_1] := e_2}{ok : e_1 \mapsto e_2}}

\inferrule [StoreEr]
  {}
  {\vdash_{\mathrm{ISL}} \itriple{e_1 \nmapsto}{[e_1] := e_2}{er : e_1 \nmapsto}}

\inferrule [StoreNull]
  {}
  {\vdash_{\mathrm{ISL}} \itriple{e_1 = \texttt{null}}{[e_1] := e_2}{er : e_1 = \texttt{null}}}

\end{mathpar}
\caption{Reglas de inferencia de la ISL. Adaptadas de~\citep{isl}.}
\label{fig:reglas-isl}
\end{figure}

La regla \textsc{Exist} permite cuantificar existencialmente una variable en la presunción y la postcondición,
siempre que no aparezca en el comando.
En esta regla denotamos por $\operatorname{fv}(C)$ las variables libres de $C$.
Luego,
la condición $x \notin \operatorname{fv}(C)$ garantiza que el comando no lee ni modifica dicha variable,
por lo que podemos omitir su valor concreto sin alterar la ejecución.

Las reglas \textsc{Alloc1} y \textsc{Alloc2} reflejan los dos casos identificados en el Lema~\ref{lem:monotonia}.
La primera reserva una dirección sobre la cual el \emph{heap} local no tenía información,
por lo que incorpora una nueva celda al fragmento descrito.
La segunda permite reutilizar una dirección que estaba representada mediante un \emph{heap} negativo.
En este caso,
la dirección ya pertenecía al dominio conocido y solo cambia su estado,
de liberado a activo.

En ambas reglas,
$x$ y $x'$ son variables distintas,
y la presunción $x = x'$ establece que almacenan el mismo valor antes de ejecutar el comando.
La alocación modifica solo $x$,
que pasa a contener la dirección reservada,
mientras que $x'$ conserva el valor inicial de $x$.
Aunque $x'$ no aparece en la postcondición de estas reglas,
puede aparecer en las aserciones incorporadas mediante la Regla del Marco,
permitiendo conservar información sobre ese valor inicial.
En \textsc{Alloc2} se exige además que $x$ no aparezca libre en $e$,
de modo que la asignación a $x$ no modifique la dirección que denota $e$.

La regla \textsc{Free} representa la liberación exitosa.
Requiere que la expresión $e$ evalúe a la dirección de una celda activa cuyo contenido está descrito por $e'$,
y transforma esa celda en un \emph{heap} negativo $e \nmapsto$.
Por otro lado,
\textsc{FreeEr} y \textsc{FreeNull} representan los casos de error.
La primera contempla el intento de liberar una celda ya marcada como liberada,
como ocurre en un error de tipo \emph{double free}.
La segunda contempla la liberación mediante un puntero nulo.

La regla \textsc{Load} representa una lectura exitosa.
Requiere que $e$ evalúe a la dirección de una celda activa cuyo contenido es $v$.
Este valor se asigna a $x$,
sin modificar el \emph{heap}.
Como en las reglas de alocación,
$x'$ conserva el valor inicial de $x$.
Por esta razón,
la postcondición utiliza $e[x' / x]$ para seguir describiendo la dirección original de la celda cuando $e$ contiene a $x$.
Las reglas \textsc{LoadEr} y \textsc{LoadNull} representan,
respectivamente, la lectura de una celda liberada y la lectura mediante un puntero nulo.

Finalmente, las reglas \textsc{Store} describen las escrituras sobre el \emph{heap}.
En el caso de éxito,
la celda cuya dirección se obtiene al evaluar $e_1$ contiene inicialmente el valor de $e$.
La escritura reemplaza ese contenido por el resultado de evaluar $e_2$,
sin modificar el \emph{store}.
Las reglas \textsc{StoreEr} y \textsc{StoreNull} representan,
respectivamente, una escritura sobre una celda liberada y una escritura mediante un puntero nulo.

La igualdad $x = x'$ es una aserción pura:
depende del \emph{store},
pero no impone condiciones sobre el \emph{heap}.
Por eso,
en estas reglas,
combinarla con una aserción espacial mediante $\land$ o mediante $*$ expresa la misma condición.
En el segundo caso,
la aserción pura puede satisfacerse sobre el \emph{heap} vacío y no consume recursos espaciales.
Las Figuras \ref{fig:reglas-sl} y \ref{fig:reglas-isl} utilizan,
respectivamente,
ambas notaciones.

\textbf{Observación.}
A diferencia de las reglas de la SL presentadas anteriormente,
las reglas \textsc{Free} y \textsc{Store} permiten especificar el contenido inicial de la celda mediante una expresión.
Reemplazar ese contenido por - produce tripletas válidas,
pero reglas más débiles para la deducción en ISL.
Por ejemplo,
la regla
$$\itriple{e_1 \mapsto -}{[e_1] := 5}{ok : e_1 \mapsto 5}$$
no permite obtener mediante la Regla de Consecuencia la tripleta
$$\itriple{e_1 \mapsto 7}{[e_1] := 5}{ok : e_1 \mapsto 5}.$$
Esto se debe a que,
en ISL,
dicha regla permite debilitar la presunción,
pero no fortalecerla.
En cambio,
una regla que admite una expresión arbitraria para el contenido inicial permite obtener esta última tripleta directamente,
eligiendo $e = 7$.

Las reglas de la \autoref{fig:reglas-isl} describen el efecto de cada operación sobre el fragmento local de memoria que utiliza.
Para extender estos razonamientos a estados que también contienen memoria adicional,
la ISL utiliza la Regla del Marco heredada de la SL:
\begin{equation}
  \inferrule [Frame]
  {\vdash_{\mathrm{ISL}} \itriple{p}{C}{\varepsilon : q} \ \text{mod}(C) \cap \text{fv}(r) = \varnothing}
  {\vdash_{\mathrm{ISL}} \itriple{p * r}{C}{\varepsilon : q * r}}
\label{eq:frame-isl}
\end{equation}

La regla muestra que el mismo marco $r$ se agrega tanto a la presunción como a la postcondición,
mientras que el comando $C$ y el modo de terminación $\varepsilon$ se mantienen sin cambios.
La regla conserva así la estructura de la Regla del Marco de la SL,
pero ahora se aplica dentro de la interpretación subaproximada de la ISL.

En particular, puede utilizarse tanto para tripletas con $\varepsilon = ok$
como para tripletas con $\varepsilon = er$.
En este último caso,
los \emph{heaps} negativos permiten conservar en la postcondición la información espacial que justifica el error.
De este modo,
una derivación local puede extenderse a un estado que contiene memoria adicional sin describir explícitamente todo el \emph{heap}.
Esta conservación de la información sobre las direcciones conocidas se apoya en la monotonía del \emph{heap} establecida en el Lema~\ref{lem:monotonia}.

Con estas reglas queda definido el sistema deductivo de la ISL.
Su relación con la semántica relacional se establece en la siguiente subsección.

\subsection*{Teorema de Soundness de la ISL}
El siguiente teorema relaciona la derivabilidad de una tripleta de la ISL con su validez respecto de la semántica relacional.

\begin{teorema}[Teorema de Soundness de la ISL]
\label{teo:sound-isl}
Para todo comando $C$, aserciones espaciales $p$ y $q$,
y un modo de terminación $\varepsilon \in \{ok, er\}$:
$$\vdash_{\mathrm{ISL}} \itriple{p}{C}{\varepsilon : q} \implies \models_{\mathrm{ISL}} \itriple{p}{C}{\varepsilon : q}.$$
\end{teorema}

El teorema establece que toda tripleta derivable mediante las reglas de la ISL
es válida respecto de la semántica relacional.
En particular,
si $\vdash_{\mathrm{ISL}} \itriple{p}{C}{er : q}$,
entonces cada estado que satisface $q$ puede alcanzarse mediante una ejecución de $C$
que comienza en un estado que satisface $p$ y termina con el modo $er$.

Para concluir además que existe al menos una ejecución errónea,
es necesario que la postcondición $q$ sea satisfacible.
Si $q = \varnothing$,
la tripleta puede ser válida de manera trivial,
pero no garantiza la existencia de una ejecución de error.

\begin{corolario}
\label{cor:errores-isl}
Si
$$\vdash_{\mathrm{ISL}} \itriple{p}{C}{er : q}$$
y $q \neq \varnothing$, entonces existen estados espaciales $\sigma_p \in p$ y $\sigma_q \in q$ tales que:
$$(\sigma_p, \sigma_q) \in [\![C]\!]_{er}.$$
\end{corolario}
\emph{Demostración.} Análoga a la demostración del Corolario~\ref{cor:errores},
utilizando el Teorema de Soundness de la ISL~\ref{teo:sound-isl} y su definición de validez semántica.
$\ensuremath{\blacksquare}$

Además de la solidez del sistema deductivo,
interesa caracterizar qué fragmento de memoria basta para justificar una ejecución.
Esto se presenta en el Teorema de la Huella.

\subsection*{El Teorema de la Huella (\textit{Footprint Theorem})}
El Lema~\ref{lem:localidad} de localidad semántica establece que una ejecución sobre un \emph{heap} extendido
puede ampliarse con un marco disjunto de sus \emph{heaps} inicial y final,
sin modificar dicho marco ni el modo de terminación.

El Teorema de la Huella completa esta propiedad en la dirección opuesta.
Mientras que la localidad permite extender una ejecución local con memoria adicional,
el teorema muestra que toda ejecución global puede descomponerse en una transición minimal y un marco adicional.

Para formalizar esta descomposición,
primero debemos definir cómo se incorpora un fragmento adicional de memoria a un estado.
Esta operación debe agregar el fragmento al \emph{heap} extendido sin modificar el \emph{store}
y solo puede realizarse cuando ambos fragmentos son disjuntos.

\begin{definicion}[Extensión de un estado por un marco]
\label{def:state-ext}
Sea $\sigma = (s, h)$ un estado espacial y sea $h_f \in \mathsf{HeapExt}$ tal que $h \perp h_f$.
Definimos:
$$\sigma \bullet h_f = (s, h \uplus h_f).$$
\end{definicion}

A partir de esta operación,
es posible definir la extensión por marcos de una relación entre estados.
Para expresarla,
si $\sigma = (s_\sigma,h_\sigma)$ es un estado espacial,
denotamos por $h(\sigma)$ su componente del \emph{heap},
es decir, $h(\sigma) = h_\sigma$.

\begin{definicion}[Extensión de una relación por marcos]
\label{def:frame-ext}
Sea $R \subseteq \Sigma_\mathrm{ISL} \times \Sigma_\mathrm{ISL}$ una relación entre estados.
Su extensión por marcos se define como:
$$\operatorname{frame}(R) = \{(\sigma_p \bullet h_f, \sigma_q \bullet h_f) \mid (\sigma_p, \sigma_q) \in R
\land h_f \in \mathsf{HeapExt} \land h(\sigma_p) \perp h_f \land h(\sigma_q) \perp h_f\}.$$
\end{definicion}

La extensión por marcos toma una transición entre dos estados
y agrega el mismo fragmento $h_f$ tanto al estado inicial como al final.
Las condiciones que exigen \emph{heaps} disjuntos garantizan que el marco no se superponga con los \emph{heaps} de esos estados.

Decimos que una transición $(\sigma_p, \sigma_q)$ es menor o igual que $(\sigma_p', \sigma_q')$ respecto de la extensión por marcos,
y escribimos $(\sigma_p, \sigma_q) \preceq (\sigma_p', \sigma_q')$,
si existe $h_f \in \mathsf{HeapExt}$ tal que:
$$\sigma_p' = \sigma_p \bullet h_f \land \sigma_q' = \sigma_q \bullet h_f.$$
Este orden no compara comandos ni modos de terminación,
sino la cantidad de memoria adicional que contienen dos transiciones.

Con esto podemos definir las huellas minimales de un comando.

\begin{definicion}[Huellas minimales]
\label{def:huella-minima}
La relación
$$\operatorname{foot}(C)_\varepsilon$$
está formada por las transiciones de $[\![C]\!]_\varepsilon$ que son minimales respecto del orden $\preceq$;
las llamamos huellas minimales.
\end{definicion}
%\guido{deberia decir ``minimal''?}

Una transición es minimal cuando no puede quitarse de sus estados inicial y final un mismo fragmento de memoria no vacío
y seguir obteniendo una transición válida de $[\![C]\!]_\varepsilon$.
Es decir,
no existe otra transición de $[\![C]\!]_\varepsilon$ estrictamente menor que ella respecto de $\preceq$.
Se puede observar que no se exige que la transición sea menor o igual que todas las demás transiciones del comando,
sino que ninguna otra se encuentre estrictamente por debajo de ella.

Una transición perteneciente a $\operatorname{foot}(C)_\varepsilon$
describe los recursos de memoria necesarios para una ejecución particular del comando.
Por ejemplo,
una ejecución exitosa de la instrucción $[x] := y$ requiere localmente una celda activa en la dirección almacenada en $x$,
pero no necesita describir las demás celdas del \emph{heap}.
El resto de la memoria puede incorporarse con un marco.

Además,
no necesariamente existe una única huella minimal para cada comando.
Un comando con distintas ramas puede tener varias ejecuciones minimales con requerimientos de memoria distintos,
y ninguna de ellas es menor que la otra.

\begin{teorema}[Teorema de la Huella]
\label{teo:huella}
Para todo comando $C$ y todo modo de terminación $\varepsilon \in \{ok, er\}$,
$$[\![C]\!]_{\varepsilon} = \operatorname{frame}(\operatorname{foot}(C)_\varepsilon).$$
\end{teorema}

La igualdad del Teorema~\ref{teo:huella} expresa dos direcciones complementarias.
La primera inclusión,
$\operatorname{frame}(\operatorname{foot}(C)_\varepsilon) \subseteq [\![C]\!]_\varepsilon$,
expresa la dirección local a global.
Toda transición minimal sigue siendo una ejecución válida cuando se le agrega un marco compatible.
Esta dirección se obtiene a partir de la propiedad de localidad presentada en el Lema~\ref{lem:localidad}.

La segunda inclusión,
$[\![C]\!]_\varepsilon \subseteq \operatorname{frame}(\operatorname{foot}(C)_\varepsilon)$,
expresa la dirección global a local.
Establece que toda ejecución global puede descomponerse en una transición minimal y un marco agregado a sus estados inicial y final.

\section{El lenguaje \fstar{}}
\label{sec:fstar}
Los resultados semánticos y deductivos presentados en este capítulo se mecanizan en los capítulos siguientes utilizando \fstar{},
un lenguaje de programación funcional con tipos dependientes que también funciona como asistente de pruebas
y verificador de programas~\citep{fstar}.
Este lenguaje está integrado con el solucionador SMT Z3
y adopta un enfoque de programación orientada a pruebas (\emph{proof-oriented programming}),
en el que los programas, sus especificaciones y sus demostraciones se desarrollan dentro de un mismo lenguaje.

Si bien \fstar{} permite desarrollar programas ejecutables,
en esta tesina se utiliza como demostrador de teoremas.
Su función es mecanizar la semántica de los lenguajes considerados,
representar los sistemas deductivos de la IL y la ISL,
y verificar sus propiedades metateóricas.

A continuación se introducen las características del lenguaje necesarias para comprender la formalización desarrollada en los siguientes capítulos.

\subsection*{Tipos dependientes y refinados}
En \fstar{}, los tipos pueden depender de valores.
La manifestación más práctica de esto son los tipos refinados.
Un tipo refinado permite expresar propiedades de los datos directamente en sus tipos,
de modo que su cumplimiento sea comprobado estáticamente.

Dado un tipo \lstinline|t| y un predicado \lstinline|P|,
la expresión \lstinline|x : t { P x }| representa los valores \lstinline|x| de tipo \lstinline|t| que además satisfacen la propiedad \lstinline|P x|.
Por ejemplo, en la biblioteca estándar de \fstar{},
el tipo de los números naturales no es un tipo primitivo aislado,
sino el tipo de los enteros refinado por una condición booleana:

\begin{lstlisting}
let nat = x : int { x >= 0 }
\end{lstlisting}

Cualquier función que requiera un argumento del tipo \lstinline|nat|
fuerza al verificador estático a demostrar matemáticamente que el valor provisto satisface la restricción $x \geq 0$ antes de compilar.
De este modo,
propiedades que en otros lenguajes se verificarían durante la ejecución se convierten en obligaciones estáticas.

Esta dependencia entre valores y tipos permite que la signatura de una función actúe como su propia especificación matemática.
Por ejemplo,
la función de incremento puede definirse garantizando estáticamente
que el valor de retorno \lstinline|y| será exactamente igual al argumento \lstinline|x| más uno:

\begin{lstlisting}
val incr : x : int -> y : int { y = x + 1 }
let incr x = x + 1
\end{lstlisting}

\subsection*{Tipos inductivos indexados y Curry-Howard}
\fstar{} hace un uso extensivo de los tipos inductivos para definir estructuras de datos y proposiciones lógicas.
Una de las características más potentes del lenguaje es la capacidad de crear tipos inductivos indexados,
los cuales permiten capturar invariantes complejos directamente en la forma de la estructura.

La correspondencia Curry-Howard \citep{howard} permite interpretar
las proposiciones como tipos y las demostraciones como valores que habitan esos tipos.
Cuando una proposición se representa mediante un tipo inductivo,
sus constructores describen las formas de construir una demostración.
A su vez,
las pruebas por inducción pueden expresarse mediante funciones recursivas que analizan esos constructores
y cuya terminación es verificada por \fstar{}.

Un ejemplo clásico en la programación de tipos dependientes es el tipo de los vectores,
donde la longitud forma parte del propio tipo:
\begin{lstlisting}
type vec (a : Type) : nat -> Type =
  | Nil : vec a 0
  | Cons : #n : nat -> hd : a -> tl : vec a n -> vec a (n + 1)
\end{lstlisting}
El constructor \lstinline|Nil| solo produce vectores de longitud cero,
mientras que \lstinline|Cons| recibe un vector de longitud \lstinline|n| y produce uno de longitud \lstinline|n + 1|.
Por lo tanto, la longitud no es una propiedad externa que deba verificarse luego,
sino un invariante impuesto por el sistema de tipos.
Esta información permite definir funciones de acceso cuyos índices estén refinados por la longitud del vector,
excluyendo estáticamente los accesos fuera de rango.

En el constructor \lstinline|Cons|,
el símbolo \lstinline|#| indica que \lstinline|n| es un parámetro implícito:
puede omitirse al aplicar el constructor,
ya que \fstar{} intenta inferirlo a partir de los demás argumentos y del contexto.
En este caso,
la longitud se obtiene del tipo del vector recibido.

Los tipos inductivos indexados también permiten representar relaciones.
Por ejemplo,
podemos definir una relación entre dos listas de enteros que tienen la misma longitud:
\begin{lstlisting}
type same_length : list int -> list int -> Type =
  | SLNil : same_length [] []

  | SLCons : x : int -> y : int ->
    xs : list int -> ys : list int ->
    same_length xs ys ->
    same_length (x :: xs) (y :: ys)
\end{lstlisting}
El constructor \lstinline|SLNil| establece que dos listas vacías tienen la misma longitud.
El constructor \lstinline|SLCons| establece que,
si dos listas tienen la misma longitud,
agregar un elemento al comienzo de cada una conserva esa propiedad.
Por lo tanto,
un valor de tipo \lstinline|same_length xs ys| representa una demostración
de que ambas listas tienen la misma longitud.
Este mecanismo permite codificar relaciones mediante tipos y sus reglas de construcción mediante constructores.

\subsection*{Código fantasma y automatización}
\fstar{} permite distinguir el código utilizado para demostrar propiedades mediante cómputos fantasma,
que son comprobados por el verificador pero eliminados durante la extracción de código.
El efecto \lstinline|Tot| describe cómputos puros y totales,
mientras que \lstinline|GTot| describe cómputos totales y fantasma.

Para enunciar y demostrar propiedades,
\fstar{} proporciona el efecto \lstinline|Lemma|.
Un lema toma ciertas premisas como argumentos mediante una cláusula \lstinline|requires|
y garantiza lógicamente una conclusión mediante la cláusula \lstinline|ensures|.
Su resultado computacional es irrelevante;
lo importante es que su cuerpo constituya una demostración aceptada por el comprobador.
Por ejemplo, consideremos la siguiente signatura:
\begin{lstlisting}
val lemma :
  x : t ->
  Lemma
    (requires P x)
    (ensures Q x)
\end{lstlisting}
Esta establece que, siempre que se cumpla la premisa \lstinline|P x|,
el lema garantiza la conclusión \lstinline|Q x|.

Retomando la relación \lstinline|same_length|,
un valor de tipo \lstinline|same_length xs ys| constituye un certificado construido mediante las reglas de la relación.
Podemos recibir este certificado como argumento de un lema que demuestre que las longitudes de ambas listas,
calculadas mediante la función \lstinline|length|,
son iguales.
En este caso,
la hipótesis está representada por el tipo del argumento \lstinline|pf|,
por lo que no necesitamos una cláusula \lstinline|requires|.
\begin{lstlisting}
let rec same_length_sound
  (#xs : list int) (#ys : list int)
  (pf : same_length xs ys) :
  Lemma
    (ensures length xs == length ys)
    (decreases pf) =
  match pf with
  | SLNil -> ()
  | SLCons _ _ _ _ tl ->
    same_length_sound tl
\end{lstlisting}

La demostración se hace por inducción sobre el certificado \lstinline|pf|.
En el caso \lstinline|SLNil|,
ambas listas son vacías y sus longitudes son cero.
En el caso \lstinline|SLCons|,
la llamada recursiva establece que las colas tienen la misma longitud.
Como ambas listas agregan un elemento,
sus longitudes también coinciden.
La cláusula \lstinline|decreases pf| permite justificar la terminación,
ya que la llamada recursiva recibe el certificado \lstinline|tl|,
que es una parte estructuralmente menor de \lstinline|pf|.

De este modo,
el tipo inductivo proporciona una estructura explícita que la demostración puede analizar,
mientras que la postcondición del lema expresa la propiedad obtenida a partir de ella.

En otras partes de la formalización basta con establecer que una proposición es válida,
sin analizar la estructura de su demostración.
Para estos casos se utiliza \lstinline|squash|:
dada una proposición \lstinline|P|,
el tipo \lstinline|squash P| conserva únicamente la información de que \lstinline|P| se cumple,
haciendo irrelevante el contenido computacional de su prueba.
Esta construcción resulta adecuada para condiciones laterales como:
\begin{lstlisting}
squash (forall s. pre s ==> pre' s)
\end{lstlisting}
que es una condición lateral de la Regla de Consecuencia.

En este caso, la formalización solo necesita establecer que la implicación es válida.
No necesita inspeccionar posteriormente cómo fue demostrada.
El uso de \lstinline|squash| también facilita que estas obligaciones sean tratadas por el \emph{solver} SMT Z3.
El \emph{solver} puede resolver automáticamente muchas obligaciones de primer orden,
como implicaciones, igualdades, restricciones aritméticas e inclusiones entre predicados.

Sin embargo, esta automatización no elimina la necesidad de construir las demostraciones.
El usuario debe proporcionar la estructura principal de los argumentos,
especialmente cuando requieren inducción o la construcción de testigos existenciales.
Cuando el \emph{solver} no puede resolver una obligación automáticamente,
es necesario agregar lemas auxiliares,
revelar definiciones o proporcionar pasos intermedios que orienten la prueba.

Estas características determinan la estrategia de mecanización utilizada en los capítulos siguientes.
Las aserciones se representan como predicados sobre estados;
las semánticas operacionales concretas y los sistemas deductivos,
como tipos inductivos indexados;
y sus propiedades metateóricas, como funciones recursivas y lemas fantasma.
La integración con Z3 permite automatizar las obligaciones lógicas locales,
mientras que la estructura global de cada demostración permanece explícita en el código de \fstar{}.

%%%%%%%%%%%%%%%%%%%%%%%%%%%%%%%%%%%%%%%%%%%%%%%%%%%%%%%%%%%%%%%%%%%%%%%%%%%%%%%%%%%%
%%%%%%%%%%%%%%%%%%%%%%%%%%%%%%%%%%%% CAPÍTULO 3 %%%%%%%%%%%%%%%%%%%%%%%%%%%%%%%%%%%%
%%%%%%%%%%%%%%%%%%%%%%%%%%%%%%%%%%%%%%%%%%%%%%%%%%%%%%%%%%%%%%%%%%%%%%%%%%%%%%%%%%%%

\chapter{Formalización de IL}
\label{cap:formalizacion-il}
\begin{quote}
``When you have eliminated the impossible, whatever remains, however improbable, must be the truth.'' \\
\raggedleft --- Arthur Conan Doyle
\end{quote}

En el capítulo anterior se establecieron los fundamentos teóricos de la IL.
En este capítulo se presenta la mecanización de este sistema deductivo utilizando el asistente de pruebas \fstar{}.

Para lograrlo, se aplica el paradigma de programación orientada a pruebas.
En primer lugar, se modela la sintaxis y semántica del lenguaje imperativo subyacente.
Luego, se codifican las reglas de inferencia de la IL a través de la correspondencia de Curry-Howard.
Por último, se demuestra su solidez (\emph{soundness}) respecto a la semántica operacional.

\section{Lenguaje y semántica operacional}
\label{sec:il-sintaxis}
Para poder razonar sobre el comportamiento de un programa,
el primer paso es definir formalmente su sintaxis y semántica.
En esta formalización inicial de la IL trabajaremos sobre un lenguaje imperativo simplificado.

Los estados de la \autoref{sec:loginc} se representan en esta formalización mediante un \emph{store} junto con un modo de terminación.
Las variables se representan mediante cadenas y sus valores mediante números naturales:

\begin{lstlisting}
type var = string
type value = nat
type store = var ^-> value
\end{lstlisting}

Un \lstinline|store| es una función total que asigna un valor a cada variable.
En la formalización se utiliza la representación de funciones extensionales provista por
\lstinline|FStar.FunctionalExtensionality|,
mediante la notación \lstinline|^->|.
Esta representación permite demostrar la igualdad entre dos \emph{stores} a partir de su coincidencia en todas las variables.
Luego, la actualización de una variable se expresa mediante la siguiente función:

\begin{lstlisting}
let override (st : store) (x : var) (v : value) : store =
  on var (fun y -> if y = x then v else st y)
\end{lstlisting}
%\guido{y si fijamos a/b a var y value? si no, hay que explicar eqtype. En el codigo
%puede ser general, aca puede ser monomorfica}
La expresión
\lstinline|override st x v| devuelve un \emph{store} que asigna \lstinline|v| a \lstinline|x| y conserva el valor de las demás variables.
La operación \lstinline|on| permite construir el resultado con la representación extensional requerida.

El lenguaje considerado distingue entre ejecuciones que continúan normalmente y ejecuciones que han alcanzado un error.
Para representar esta información se definen los siguientes tipos:

\begin{lstlisting}
type term_mode =
  | Ok
  | Er
type state = store & term_mode
\end{lstlisting}

El símbolo \lstinline|&| denota el tipo producto,
por lo que cada estado es un par formado por un \emph{store} y un modo de terminación.
Para un estado \lstinline|s|,
las proyecciones \lstinline|s._1| y \lstinline|s._2| permiten acceder a estos componentes,
respectivamente.

Con respecto a \lstinline|term_mode|,
el modo \lstinline|Ok| indica que la ejecución puede continuar,
mientras que \lstinline|Er| indica que se ha alcanzado un estado de error.
Esta distinción mecaniza los dos modos de terminación \emph{ok} y \emph{er}.
En la presentación abstracta del \autoref{cap:conceptos-previos},
el modo de terminación $\varepsilon \in \{ok, er\}$ aparece como un índice de la tripleta.
En esta formalización se incorpora directamente en el estado.
De esta manera,
una aserción puede distinguir los estados normales de los erróneos inspeccionando su segundo componente.
%\guido{y la semántica no permite que un estado de error avance en su ejecucion?}.

Las expresiones del lenguaje se definen mediante un tipo inductivo:

\begin{lstlisting}
type expr =
  | Var : var -> expr
  | Const : nat -> expr
  | Plus : expr -> expr -> expr
  | Minus : expr -> expr -> expr
  | Times : expr -> expr -> expr
  | Eq : expr -> expr -> expr
  | Lt : expr -> expr -> expr
  | Gt : expr -> expr -> expr
\end{lstlisting}

Las expresiones en nuestro lenguaje son puras,
es decir,
calcular $2 + 2$ no modifica la memoria ni puede fallar o arrojar múltiples resultados distintos.
Por lo tanto,
su semántica se modela mediante una función recursiva pura,
y una función adaptadora que toma el \emph{store} de la tupla del estado actual y la expresión y devuelve un natural:

\begin{lstlisting}
let rec eval_expr' (s : store) (e : expr) : GTot nat =
  match e with
  | Var x -> s x
  | Const n -> n
  | Plus e1 e2 -> eval_expr' s e1 + eval_expr' s e2
  | Minus e1 e2 ->
    let res = eval_expr' s e1 - eval_expr' s e2 in
    if res >= 0 then res else 0
  | Times e1 e2 -> eval_expr' s e1 * eval_expr' s e2
  | Eq e1 e2 -> if eval_expr' s e1 = eval_expr' s e2
                then 1 else 0
  | Lt e1 e2 -> if eval_expr' s e1 < eval_expr' s e2
                then 1 else 0
  | Gt e1 e2 -> if eval_expr' s e1 > eval_expr' s e2
                then 1 else 0

let eval_expr (s : state) (e : expr) : GTot nat =
  eval_expr' s._1 e
\end{lstlisting}

Como los valores son naturales,
la resta se trunca en cero cuando el segundo operando supera el primero.
A su vez,
las comparaciones devuelven uno cuando se cumplen y cero en caso contrario.
No se introduce un tipo separado de expresiones booleanas:
en los comandos de filtrado,
cero representa falso y cualquier valor distinto de cero representa verdadero.

La definición de comandos es similar,
pero incluye instrucciones que modifican el estado o controlan el flujo:

\begin{lstlisting}
type stmt =
  | Assign : var -> expr -> stmt
  | Nondet : var -> stmt
  | Skip : stmt
  | Error : stmt
  | Assume : expr -> stmt
  | Seq : stmt -> stmt -> stmt
  | Choice : stmt -> stmt -> stmt
  | Kleene : stmt -> stmt
\end{lstlisting}

\lstinline|Assign x e| evalúa la expresión \lstinline|e| y almacena su resultado en \lstinline|x|.
\lstinline|Nondet x| asigna un valor arbitrario a la variable.
\lstinline|Skip| no modifica el estado, mientras que \lstinline|Error| fuerza la terminación errónea.
\lstinline|Assume e| conserva solo los caminos donde \lstinline|e| es verdadera.
\lstinline|Seq p q| representa composición secuencial,
\lstinline|Choice p q| elección no determinista y \lstinline|Kleene p| una cantidad finita y no determinista de ejecuciones de \lstinline|p|.
Las construcciones derivadas \lstinline|if|, \lstinline|while| y \lstinline|assert| pueden expresarse mediante estos constructores,
siguiendo las expansiones de la \autoref{sec:loginc},
por lo que no requieren constructores adicionales.

La semántica relacional presentada en la \autoref{sec:loginc} se mecaniza mediante un tipo inductivo indexado
por el comando y sus estados inicial y final.
A diferencia de aquella presentación,
el modo de terminación se incorpora al estado final,
en lugar de aparecer como un índice separado de la relación:
\begin{lstlisting}
type runsto : stmt -> state -> state -> Type0
\end{lstlisting}

Un valor de tipo \lstinline|runsto c s0 s1| representa una derivación
de la ejecución del comando \lstinline|c| desde \lstinline|s0| hasta \lstinline|s1|.
Sus constructores codifican las reglas de la semántica operacional,
siguiendo la representación de relaciones mediante tipos inductivos indexados.
Las ejecuciones comienzan en un estado con modo \lstinline|Ok| y pueden finalizar con modo \lstinline|Ok| o \lstinline|Er|,
según terminen normalmente o con error.

Como los \emph{stores} son funciones,
dos expresiones funcionales diferentes pueden representar el mismo \emph{store} si coinciden en todas las variables.
Para tratar estas diferencias conviene definir la equivalencia de estados:

\begin{lstlisting}
let state_equiv (s1 s2 : state) : prop =
  (forall x. s1._1 x == s2._1 x) /\ s1._2 == s2._2
\end{lstlisting}

Con la representación extensional de los \emph{stores},
esta equivalencia implica la igualdad de los estados.
El siguiente lema establece esta propiedad:

\begin{lstlisting}
let state_equiv_eq (s1 s2 : state)
  : Lemma (requires (state_equiv s1 s2))
          (ensures s1 == s2) =
  extensionality var (fun _ -> value) s1._1 s2._1
\end{lstlisting}
%\guido{veo que \texttt{R\_Ext} ya no esta en el codigo, bien}

La llamada a \lstinline|extensionality| demuestra la igualdad de los \emph{stores} a partir de su coincidencia punto a punto.
Junto con la igualdad de los modos de terminación,
permite concluir la igualdad de los estados en su totalidad.

A partir de esta demostración de igualdad entre dos estados,
una derivación de \lstinline|runsto| que comienza o termina en uno de los estados también sirve para el otro.
Este reemplazo se justifica mediante la igualdad demostrada,
sin necesidad de agregar una regla a la semántica operacional.

A continuación,
se presentan los constructores de \lstinline|runsto| que describen la ejecución de los comandos.
Las reglas de \lstinline|runsto| exigen que el estado inicial tenga modo \lstinline|Ok|,
por lo que no permiten ejecutar un comando desde un estado con modo \lstinline|Er|.
Para simplificar sus signaturas,
se definen dos abreviaturas para los estados según su modo de terminación:

\begin{lstlisting}
type ok_state = s : state { s._2 == Ok }
type er_state = s : state { s._2 == Er }  
\end{lstlisting}

En el caso de la asignación,
se evalúa la expresión en el estado inicial y se actualiza el valor de la variable indicada:

\begin{lstlisting}
| R_Assign : s : ok_state -> #x : var -> #e : expr ->
  runsto (Assign x e) s (override s._1 x (eval_expr s e), Ok)
\end{lstlisting}
%\guido{ayudaria definir un \lstinline{ok_state} para no repetir el refinamiento?}
Toma un estado inicial \lstinline|s| con modo \lstinline|Ok|,
evalúa matemáticamente la expresión utilizando la función determinista \lstinline|eval_expr s e|,
y garantiza que el estado final es idéntico a \lstinline|s|,
pero con el valor de la variable \lstinline|x| actualizado mediante \lstinline|override|.
Por último, como \lstinline|Assign| es una operación segura en nuestro lenguaje,
el modo de terminación siempre es \lstinline|Ok|.

La asignación no determinista se representa mediante:

\begin{lstlisting}
| R_Nondet : s : ok_state -> #x : var -> v : value ->
  runsto (Nondet x) s (override s._1 x v, Ok)
\end{lstlisting}
El valor \lstinline|v| es un parámetro del constructor.
Para cada valor posible,
\lstinline|R_Nondet| permite construir una derivación que actualiza \lstinline|x| con ese valor.
De este modo,
la relación representa todas las elecciones permitidas por el comando.
Al construir una derivación concreta,
se proporciona el valor correspondiente a la ejecución que se desea justificar.

Representamos la terminación normal,
la terminación con error y el filtrado de caminos de la siguiente manera:

\begin{lstlisting}
| R_Skip : s : ok_state -> runsto Skip s s

| R_Error : s : ok_state -> runsto Error s (s._1, Er)

| R_Assume : s : ok_state -> #e : expr ->
  squash (eval_expr s e =!= 0) ->
  runsto (Assume e) s s
\end{lstlisting}

La regla \lstinline|R_Skip| conserva el estado inicial.
Por su parte,
\lstinline|R_Error| representa el comando \lstinline|error|:
conserva el \emph{store} y cambia el modo de \lstinline|Ok| a \lstinline|Er|.
Finalmente,
\lstinline|R_Assume| permite una ejecución que conserva el estado únicamente
cuando la expresión se evalúa a un valor distinto de cero.
Cuando la condición es falsa,
no existe una derivación de ejecución para ese comando desde ese estado;
esto no constituye una terminación con error.

La composición secuencial se divide según el modo de terminación del primer comando:

\begin{lstlisting}
| R_Seq : #p : stmt -> #q : stmt -> #s : ok_state -> #t : ok_state -> #u : state ->
  runsto p s t -> runsto q t u ->
  runsto (Seq p q) s u

| R_SeqEr : #p : stmt -> #q : stmt -> #s : ok_state -> #t : er_state ->
  runsto p s t -> runsto (Seq p q) s t
\end{lstlisting}


La regla \lstinline|R_Seq| describe el caso en que \lstinline|p| termina normalmente en un estado intermedio \lstinline|t|
y la ejecución continúa con \lstinline|q| desde ese estado.
El estado final \lstinline|u| es el resultado de esta segunda ejecución,
que puede terminar normalmente o con error.
En cambio,
\lstinline|R_SeqEr| establece que,
si \lstinline|p| termina con error en \lstinline|t|,
la secuencia completa termina en ese mismo estado.
No se requiere una derivación de la ejecución de \lstinline|q|,
ya que el segundo comando no llega a ejecutarse.

La elección no determinista posee una regla para cada rama:

\begin{lstlisting}
| R_ChoiceL : #p : stmt -> #q : stmt -> #s : ok_state -> #t : state ->
  runsto p s t -> runsto (Choice p q) s t

| R_ChoiceR : #p : stmt -> #q : stmt -> #s: ok_state -> #t : state ->
  runsto q s t -> runsto (Choice p q) s t
\end{lstlisting}
Para justificar una transición de \lstinline|Choice p q| alcanza con proporcionar una derivación de una de sus ramas.

Por último,
la estrella de Kleene representa una cantidad finita y no determinista de ejecuciones.
El caso base realiza cero iteraciones,
mientras que el caso inductivo agrega una nueva ejecución de \lstinline|p| a una repetición previa:

\begin{lstlisting}
| R_Kleene0 : #p : stmt -> #s : ok_state ->
  runsto (Kleene p) s s

| R_KleeneS : #p : stmt -> #s : ok_state -> #t : state ->
  runsto (Seq (Kleene p) p) s t ->
  runsto (Kleene p) s t
\end{lstlisting}
%\guido{alguna razon aca para anidar el Kleene p a la izquierda en vez de a la derecha?
%son lo mismo, pero a veces uno es mas facil de trabajar que el otro. Aca posiblemente
%esta forma sea mas comoda}

\section{La IL como tipo inductivo}
\label{sec:il-tipoind}
La Definición~\ref{def:derivabilidad} introdujo la derivabilidad como la existencia de un árbol finito
construido mediante las reglas de inferencia de la IL.
Utilizando la correspondencia Curry-Howard y los tipos inductivos indexados presentados en la \autoref{sec:fstar},
esta noción se mecaniza mediante el tipo \lstinline|il_triple|.

Una aserción puede interpretarse como un predicado sobre el conjunto de estados.
Dado que en la \autoref{sec:il-sintaxis} dicho conjunto fue instanciado mediante el tipo \lstinline|state|,
las aserciones se representan en \fstar{} como:

\begin{lstlisting}
type cond = state -> prop
\end{lstlisting}

Esta representación corresponde a un \emph{shallow embedding} de las aserciones:
en lugar de definir una sintaxis para el lenguaje de aserciones y su interpretación,
cada aserción es directamente un predicado de \fstar{} sobre estados.
Las consecuencias de esta decisión se discuten en la \autoref{sec:isl-tipoind}.

A partir de esta definición, la derivabilidad de una tripleta se representa mediante:

\begin{lstlisting}
type il_triple : (pre : cond) -> (p : stmt) -> (post : cond) -> Type
\end{lstlisting}

Un valor de tipo \lstinline|il_triple pre p post| es un árbol de derivación construido mediante las reglas del sistema.
Por lo tanto,
\lstinline|il_triple| mecaniza la noción sintáctica de derivabilidad de la Definición~\ref{def:derivabilidad},
no la validez semántica.
La relación entre ambas nociones será establecida en la \autoref{sec:il-soundness} mediante la función \lstinline|soundness|.

Cada constructor de \lstinline|il_triple| corresponde a una adaptación de las reglas de inferencia presentadas en la \autoref{fig:reglas-il}.
Del mismo modo que un término de \lstinline|runsto p s0 s1| constituye un testigo de ejecución,
un término de \lstinline|il_triple pre p post| constituye un testigo de derivabilidad.

Las presunciones de las reglas describen estados desde los cuales la ejecución todavía puede continuar,
es decir, tienen modo \lstinline|Ok|.
Para expresarlo se utiliza:

\begin{lstlisting}
let is_ok (c : cond) : cond = fun s -> c s /\ s._2 == Ok
\end{lstlisting}

Así, \lstinline|is_ok pre| exige que el estado satisfaga \lstinline|pre| y también que su modo de terminación sea \lstinline|Ok|.
A continuación se analiza cómo se codifican estructuralmente las reglas de la teoría.

En primer lugar, observamos las instrucciones más simples de nuestro lenguaje:

\begin{lstlisting}
| I_Skip : pre : cond ->
  il_triple (is_ok pre) Skip (is_ok pre)

| I_Error : pre : cond ->
  il_triple (is_ok pre) Error
    (fun s -> s._2 == Er /\ pre (s._1, Ok))

| I_Assume : pre : cond -> #e : expr ->
  il_triple (is_ok pre) (Assume e)
    (is_ok (fun s -> pre s /\ eval_expr s e =!= 0))
\end{lstlisting}
%\guido{Mostrá la regla de Consequence}
La regla \lstinline|I_Skip| conserva la aserción.
Por el contrario,
la regla \lstinline|I_Error| modela la interrupción del flujo normal.
Conserva el \emph{store} descrito por la presunción,
pero cambia el modo de terminación.
\lstinline|I_Assume| incorpora la condición evaluada al resultado.
La postcondición conserva solo los estados que satisfacen la presunción y para los cuales la expresión es verdadera.

Como se explicó en la \autoref{sec:loginc},
la regla de asignación de la IL conserva información sobre el valor anterior de la variable para justificar la alcanzabilidad del estado final.
En la formalización,
esta información se representa mediante un cuantificador existencial:

\begin{lstlisting}
| I_Assign : #pre : cond -> #x : var -> #e : expr ->
  il_triple (is_ok pre) (Assign x e)
    (is_ok (fun s -> exists x_init.
      pre ((override s._1 x x_init), s._2) /\
      (s._1 x = eval_expr ((override s._1 x x_init), s._2) e)))
\end{lstlisting}

El valor existencial \lstinline|x_init| representa el valor que tenía \lstinline|x| antes de la asignación.
A partir del \emph{store} final se reconstruye un posible \emph{store} inicial reemplazando el valor actual de \lstinline|x| por \lstinline|x_init|.
El primer miembro de la conjunción exige que este estado reconstruido satisfaga \lstinline|pre|.
El segundo establece que el valor actual de \lstinline|x| coincide con la evaluación de \lstinline|e| en dicho estado inicial.

La regla \lstinline|I_Nondet| usa una construcción similar.
En este caso,
para cada estado final se exige la existencia de algún valor anterior de \lstinline|x| 
que permita reconstruir un estado inicial que satisfaga la presunción.
Se representa mediante:

\begin{lstlisting}
| I_Nondet : #x : var -> #pre : cond ->
  il_triple (is_ok pre) (Nondet x)
    (is_ok (fun s -> exists v. pre ((override s._1 x v), s._2)))
\end{lstlisting}

El constructor de composición secuencial recibe las derivaciones de ambos comandos
y expresa los posibles resultados mediante una disyunción en la postcondición:

\begin{lstlisting}
| I_Seq : #p : stmt -> #q : stmt ->
  #pre : cond -> #mid : cond -> #post : cond ->
  il_triple (is_ok pre) p mid ->
  il_triple (is_ok mid) q post ->
  il_triple (is_ok pre) (Seq p q)
    (fun s -> post s \/ (s._2 == Er /\ mid s))
\end{lstlisting}

La postcondición reúne dos clases de resultados.
\lstinline|post s| contiene los estados obtenidos después de ejecutar tanto \lstinline|p| como \lstinline|q|.
El segundo disjunto 
%\guido{el segundo disjunto, la disjunción es la expresión entera}
contiene los estados de error producidos durante \lstinline|p|,
en los cuales \lstinline|q| no se ejecuta.
Cuando \lstinline|p| finaliza normalmente,
el estado intermedio satisface \lstinline|is_ok mid|,
que es la presunción utilizada para \lstinline|q|.
Cuando \lstinline|p| finaliza en \lstinline|Er|,
el error se propaga directamente al resultado de la secuencia.

La elección no determinista se codifica mediante dos constructores,
uno por cada rama:

\begin{lstlisting}
| I_ChoiceL : #p : stmt -> #q : stmt ->
  #pre : cond -> #post : cond ->
  il_triple (is_ok pre) p post ->
  il_triple (is_ok pre) (Choice p q) post

| I_ChoiceR : #p : stmt -> #q : stmt ->
  #pre : cond -> #post : cond ->
  il_triple (is_ok pre) q post ->
  il_triple (is_ok pre) (Choice p q) post
\end{lstlisting}
El constructor \lstinline|I_ChoiceL| utiliza una derivación para \lstinline|p|,
mientras que \lstinline|I_ChoiceR| utiliza una para \lstinline|q|.
En ambos casos,
la presunción y la postcondición se conservan al pasar a \lstinline|Choice p q|.

El caso de las iteraciones finitas se representa mediante:

\begin{lstlisting}
| I_Kleene0 :
  #p : stmt -> #pre : cond ->
  il_triple (is_ok pre) (Kleene p) (is_ok pre)

| I_KleeneS : #p : stmt -> #pre : cond -> #post : cond ->
  il_triple (is_ok pre) (Seq (Kleene p) p) post ->
  il_triple (is_ok pre) (Kleene p) post
\end{lstlisting}

Además,
la formalización incluye una regla paramétrica basada en una familia de aserciones:

\begin{lstlisting}
let kleene_pre (variant : nat -> cond) : cond =
  is_ok (variant 0)

let kleene_post (variant : nat -> cond) : cond =
  is_ok (fun s -> exists (n : nat). variant n s)

| I_KleeneVariant :
  #variant : (nat -> cond) -> #p : stmt ->
  (n : nat ->
    il_triple (is_ok (variant n)) p (is_ok (variant (n + 1)))) ->
  il_triple (kleene_pre variant) (Kleene p) (kleene_post variant)
\end{lstlisting}

La familia de aserciones se representa mediante el parámetro \lstinline|variant : nat -> cond|,
mientras que la premisa inductiva se codifica como una función que,
para cada natural \lstinline|n|,
proporciona una derivación de la tripleta correspondiente.
A partir de esta familia,
las definiciones auxiliares \lstinline|kleene_pre| y \lstinline|kleene_post| expresan las aserciones de la conclusión:
la primera corresponde a cero iteraciones y la segunda utiliza un cuantificador existencial sobre el número de iteraciones.
En ambas,
el uso de \lstinline|is_ok| restringe los estados a aquellos con terminación normal.
%\guido{Acá parece que la estas explicando desde cero, pero ya lo explicamos (mas o menos)
%en 2.3 cuando se introduce IL. Lo ideal acá es que expliques como se escriben esas reglas
%como tipo inductivo y si tomaste decisiones interesantes (e.g. usar el xinit en asignacion).
%No hay que explicar las reglas en sí mismas, eso ya lo hicimos.}

Por otra parte,
como se observó en el Lema~\ref{lem:prop-pi},
la postimagen preserva uniones.
Esta propiedad se refleja mediante:

\begin{lstlisting}
| I_Disjunction : #pre1 : cond -> #pre2 : cond ->
  #p : stmt -> #post1 : cond -> #post2 : cond ->
  il_triple (is_ok pre1) p post1 ->
  il_triple (is_ok pre2) p post2 ->
  il_triple (is_ok (fun s -> pre1 s \/ pre2 s)) p
    (fun s -> post1 s \/ post2 s)
\end{lstlisting}
La regla permite combinar derivaciones obtenidas a partir de conjuntos diferentes de estados iniciales.

%\section{La Regla de Consecuencia}
%\guido{merece realmente otra sección? es porque hacés referencia
%a esto más adelante?}
%\label{sec:il-consecuencia}
La Regla de Consecuencia se representa mediante un constructor que recibe una derivación
y las pruebas de las dos implicaciones requeridas:
%\guido{de nuevo, esto ya está explicado, acá deberías presentar
%la formalización nomás. Podés explicar que las implicaciones
%se modelan como una implicación squasheada y que generalmente
%el solver las maneja fácilmente}

\begin{lstlisting}
| I_Consequence : #pre : cond -> #p : stmt ->
  #post : cond -> pre' : cond -> post' : cond ->
  il_triple (is_ok pre) p post ->
  squash (forall x. pre x ==> pre' x) ->
  squash (forall x. post' x ==> post x) ->
  il_triple (is_ok pre') p post'
\end{lstlisting}

La clave en esta formalización está en haber envuelto estas proposiciones universales dentro del \lstinline|squash|.
Esto hace que la formalización solo necesite conocer la validez de las condiciones laterales,
sin necesidad de inspeccionar la estructura computacional de sus demostraciones.
\fstar{} transforma estas condiciones en obligaciones de verificación que,
en muchos casos,
pueden ser resueltas automáticamente por Z3.

\section{El Teorema de Soundness}
\label{sec:il-soundness}
El último paso para completar la formalización de la IL es demostrar que el sistema deductivo es sólido respecto a la semántica.
En el marco teórico,
esto se vio en el Teorema~\ref{teo:sound-il}.
Sin embargo,
enunciamos otra vez el teorema sobre la base de nuestra formalización.

\begin{teorema}[Soundness de la IL formalizada]
Para todo comando \lstinline|p|,
presunción \lstinline|pre| y postcondición \lstinline|post|,
si existe una derivación
$$\vdash_{\mathrm{IL}} \itriple{\text{\lstinline|pre|}}{\text{\lstinline|p|}}{\text{\lstinline|post|}}$$
entonces,
para todo estado final \lstinline|s1| que satisface \lstinline|post|,
existe un estado inicial \lstinline|s0| que satisface \lstinline|pre| y una ejecución operacional de \lstinline|p| desde \lstinline|s0| hasta \lstinline|s1|:
$$\forall \text{\lstinline|s1|}. \ \text{\lstinline|post s1|} \implies \exists \text{\lstinline|s0|}. \ \text{\lstinline|pre s0|} \land \text{\lstinline|runsto p s0 s1|}$$
\end{teorema}
En este enunciado,
el modo de terminación está incluido en los estados y puede restringirse mediante las aserciones,
en lugar de aparecer como un índice separado de la tripleta.

En \fstar{} codificamos este teorema mediante una función fantasma con tipos dependientes
que proporciona un estado inicial y una derivación operacional:

\begin{lstlisting}
let rec soundness
  (p : stmt) (pre : cond) (post : cond)
  (pf : il_triple pre p post) (s1 : state { post s1 })
  : GTot (s0 : state { pre s0 } & runsto p s0 s1) (decreases pf)
\end{lstlisting}
El resultado es un par dependiente:
su primera componente es un estado \lstinline|s0| que satisface \lstinline|pre|,
y la segunda es una derivación de \lstinline|runsto p s0 s1| cuyo tipo depende de ese estado.
La notación \lstinline{(|s0, r|)} utilizada en la prueba construye dicho par.

La signatura constituye una traducción directa del Lema~\ref{lem:caracterizacion}:
\begin{itemize}
  \item \lstinline|s1 : state { post s1 }| representa un estado final $\sigma_q \in q$;
  \item \lstinline|s0 : state { pre s0 }| representa un estado inicial testigo $\sigma_p \in p$;
  \item \lstinline|runsto p s0 s1| representa la transición $(\sigma_p, \sigma_q) \in [\![C]\!]_\varepsilon$.
\end{itemize}
La función conecta así los dos niveles mecanizados en las secciones anteriores.
Consume un habitante de \lstinline|il_triple|,
que certifica derivabilidad,
y produce un habitante de \lstinline|runsto|,
que certifica la existencia de una ejecución.

La demostración se realiza por inducción estructural sobre \lstinline|pf|.
A partir del estado final \lstinline|s1| que satisface \lstinline|post| y de la derivación recibida,
cada caso obtiene un estado inicial que satisface \lstinline|pre| y una ejecución que conecta ambos estados.
La cláusula \lstinline|decreases pf| indica que cada llamada recursiva se realiza sobre una subderivación menor.
A continuación,
detallamos algunos casos de la prueba.

Los casos de \lstinline|I_Skip|,
\lstinline|I_Error| e \lstinline|I_Assume| permiten reconstruir la ejecución de forma directa:

\begin{lstlisting}
| I_Skip _ ->
  let s0 = s1 in
  let r = R_Skip s0 in
  (|s0, r|)

| I_Error _ ->
  let (st1, _) = s1 in
  let s0 = (st1, Ok) in
  let r = R_Error s0 in
  (|s0, r|)

| I_Assume pre #e ->
  let s0 = s1 in
  let r = R_Assume s0 #e () in
  (|s0, r|)
\end{lstlisting}

En el caso \lstinline|I_Skip|,
el estado inicial coincide con el final y la ejecución se justifica mediante \lstinline|R_Skip|.
Para \lstinline|I_Error|,
el \emph{store} inicial coincide con el final,
pero se reemplaza su modo de terminación por \lstinline|Ok| para reconstruir el estado inicial.
Por último,
en \lstinline|I_Assume|,
recordamos que la regla tiene \lstinline|pre s| en su postcondición,
con lo cual \lstinline|s1| satisface \lstinline|pre|,
y podemos tomar \lstinline|s0 = s1|.
Esto tiene sentido intuitivo:
los estados finales están más restringidos que los estados iniciales,
y la operación \lstinline|Assume| no modifica el estado,
con lo cual es trivial demostrar la solidez de esta regla.

Las reglas de asignación y asignación no determinista contienen cuantificadores existenciales en sus postcondiciones.
Sus testigos permiten reconstruir un estado inicial desde el cual se alcanza el estado final recibido.

En el caso de asignación,
la postcondición garantiza la existencia de un valor histórico \lstinline|x_init| tal que,
al colocarlo en la variable \lstinline|x| del \emph{store} final,
se obtiene un estado que satisface la presunción.
Para seleccionar un valor anterior que satisfaga estas condiciones,
nos apoyamos en la operación fantasma \lstinline|indefinite_description_ghost|.
Esta operación de elección clásica,
disponible en el módulo \lstinline|FStar.IndefiniteDescription|,
permite obtener un testigo a partir de la existencia lógica de un valor.
Como la prueba tiene efecto \lstinline|GTot|,
el testigo puede utilizarse dentro de la demostración fantasma,
sin formar parte de un programa extraído ejecutable.

Con el valor elegido se reconstruye el \emph{store} inicial
mediante \lstinline|override st1 x x_init| y se define \lstinline|s0 = (st0, Ok)|:

\begin{lstlisting}
| I_Assign #pre #x #e ->
  let (st1, m) = s1 in
  let x_init = indefinite_description_ghost 
    _ (fun x_init -> pre ((override st1 x x_init), m) 
    /\ (st1 x = eval_expr ((override st1 x x_init), m) e))
  in
  let st0 = override st1 x x_init in
  let s0 = (st0, Ok) in
  let pf0 = R_Assign s0 #x #e in
  state_equiv_eq (override st0 x (eval_expr s0 e), Ok) s1;
  (|s0, pf0|)
\end{lstlisting}

Una vez reconstruido el \emph{store} inicial,
se aplica \lstinline|R_Assign| para obtener una derivación de la asignación.
Para que esta derivación alcance el estado final \lstinline|s1|,
se comprueba que el \emph{store} resultante tiene los mismos valores que el de \lstinline|s1| y que ambos estados tienen modo \lstinline|Ok|.
El lema \lstinline|state_equiv_eq| permite concluir que estos estados son iguales,
por lo que se puede utilizar directamente la derivación obtenida.

El caso \lstinline|I_Nondet| sigue el mismo patrón.
La postcondición proporciona un valor histórico que permite reconstruir un estado inicial,
mientras que la ejecución se instancia con el valor final de \lstinline|x|:

\begin{lstlisting}
| I_Nondet #x #pre ->
  let (st1, m) = s1 in
  let v = indefinite_description_ghost _ (fun v -> pre ((override st1 x v), m)) in
  let st0 = override st1 x v in
  let s0 = (st0, Ok) in
  let pf0 : runsto (Nondet x) s0 ((override st0 x (st1 x)), Ok) =
    R_Nondet s0 #x (st1 x)
  in
  state_equiv_eq (override st0 x (st1 x), Ok) s1;
  (|s0, pf0|)
\end{lstlisting}

En el caso \lstinline|I_Seq|,
el estado final puede estar justificado por una ejecución de ambos comandos o por una ejecución del primer comando \lstinline|p|
que termina con error antes de ejecutar el segundo \lstinline|q|.
La función \lstinline|soundness| distingue estas posibilidades para reconstruir una ejecución que alcance \lstinline|s1|:

\begin{lstlisting}
| I_Seq #p #q #pre #mid #post pf_p pf_q ->
  if t2b (post s1) then
    let (|s_mid, r_q|) = soundness q (is_ok mid) post pf_q s1 in
    let (|s0, r_p|) = soundness p (is_ok pre) mid pf_p s_mid in
    let r = R_Seq #p #q #s0 #s_mid #s1 r_p r_q in
    (|s0, r|)
  else
    let (|s0, r_p|) = soundness p (is_ok pre) mid pf_p s1 in
    let r = R_SeqEr #p #q #s0 #s1 r_p in
    (|s0, r|)
\end{lstlisting}

La comprobación mediante \lstinline|t2b| se realiza en el contexto fantasma de la prueba.
Si \lstinline|s1| satisface \lstinline|post|,
se reconstruye una ejecución de ambos comandos.
Primero se aplica la hipótesis inductiva a \lstinline|q|,
obteniendo el estado intermedio \lstinline|s_mid|,
y luego a \lstinline|p|,
obteniendo el estado inicial \lstinline|s0|.
Finalmente,
ambas derivaciones se componen mediante \lstinline|R_Seq|.

Si \lstinline|s1| no satisface \lstinline|post|,
la postcondición de \lstinline|I_Seq| garantiza que tiene modo \lstinline|Er| y satisface \lstinline|mid|.
En ese caso,
se reconstruye únicamente la ejecución de \lstinline|p| y se utiliza \lstinline|R_SeqEr| para propagar el fallo.
Si el estado final satisface tanto \lstinline|post| como \lstinline|mid| y tiene modo \lstinline|Er|,
ambas reconstrucciones son posibles.
Cuando esto ocurre,
la prueba utiliza la ejecución de ambos comandos,
pues basta con justificar una ejecución que llegue a ese estado.

Los casos de elección no determinista aplican la hipótesis inductiva solo en la rama seleccionada por la derivación:

\begin{lstlisting}
| I_ChoiceL #p #q #pre #post pf_p ->
  let (|s0, r_p|) = soundness p (is_ok pre) post pf_p s1 in
  let r = R_ChoiceL #p #q #s0 #s1 r_p in
  (|s0, r|)
\end{lstlisting}
El caso derecho es análogo.
No es necesario reconstruir una ejecución de la rama descartada.

Los casos del desenrollado de bucles son también simples.
El caso de \lstinline|I_Kleene0| es análogo al de \lstinline|I_Skip|:
se toma el estado final como estado inicial y se utiliza \lstinline|R_Kleene0| para justificar una ejecución sin iteraciones.
Para \lstinline|I_KleeneS|,
se aplica la hipótesis inductiva en la derivación de \lstinline|Seq (Kleene p) p| recibida como premisa.
Esto proporciona un estado inicial y una derivación operacional de la secuencia,
que \lstinline|R_KleeneS| convierte en una derivación de \lstinline|Kleene p| entre los mismos estados.

El tratamiento de la iteración mediante \lstinline|I_KleeneVariant| usa una función auxiliar \lstinline|aux|,
recursiva sobre el número de iteraciones.
Esta recibe un natural \lstinline|m| y un estado que satisface \lstinline|variant m|,
y obtiene un estado inicial que satisface \lstinline|variant 0|,
junto con una ejecución de \lstinline|Kleene p| hasta el estado recibido.
Ambos estados tienen modo \lstinline|Ok|.
Este caso se define de la siguiente manera:

\begin{lstlisting}
| I_KleeneVariant #variant #p pf_var ->
  let p_n (n : nat) : prop = variant n s1 /\ s1._2 == Ok in
  let n = indefinite_description_ghost _ p_n in
  let rec aux (m : nat) (t : state { variant m t /\ t._2 == Ok } )
    : GTot (s0 : state { variant 0 s0 /\ s0._2 == Ok } & runsto (Kleene p) s0 t) (decreases m) =
    if m = 0 then
      let s0 = t in
      let r = R_Kleene0 #p in
      (|s0, r|)
    else 
      let m' = m - 1 in
      let pf_p = pf_var m' in
      let (|s_mid, r_p|) =
        soundness p (is_ok (variant m')) (is_ok (variant (m' + 1))) pf_p t in
      let (|s0, r_kleene|) = aux m' s_mid in
      let r = R_KleeneS #p (R_Seq r_kleene r_p) in
      (|s0, r|)
  in
  aux n s1
\end{lstlisting}
Primero,
utiliza la descripción indefinida para seleccionar un número de iteraciones \lstinline|n| que satisfaga la postcondición.
Luego,
aplica \lstinline|aux n s1| para reconstruir una derivación operacional de esas iteraciones.
El caso base utiliza \lstinline|R_Kleene0|.
En el caso inductivo,
se reconstruye primero la última ejecución de \lstinline|p|,
luego la repetición anterior y finalmente se componen ambas mediante \lstinline|R_Seq| y \lstinline|R_KleeneS|.
%\guido{y si mostramos el tipo de este lema inductivo auxiliar?}

En \lstinline|I_Consequence|,
el estado final que satisface \lstinline|post'| también satisface \lstinline|post| por la condición lateral de fortalecimiento.
Se aplica entonces la hipótesis inductiva a la derivación original.
El estado inicial obtenido satisface \lstinline|pre| y,
por debilitamiento, 
también satisface \lstinline|pre'|,
por lo que puede devolverse junto con la misma derivación operacional.

Finalmente,
en \lstinline|I_Disjunction|,
el estado final satisface al menos una de las dos postcondiciones.
La prueba distingue los casos mediante \lstinline|t2b| y aplica la hipótesis inductiva a la derivación correspondiente.
Si ambas postcondiciones se cumplen,
cualquiera de las dos derivaciones permite justificar el resultado;
la prueba utiliza la primera.
El estado inicial obtenido satisface la presunción de esa derivación y,
por lo tanto,
también la disyunción de las presunciones,
por lo que puede devolverse con la misma derivación operacional.

La función \lstinline|soundness| instancia el Teorema~\ref{teo:sound-il} para el lenguaje concreto de este capítulo.
En particular,
demuestra,
en el contexto fantasma,
la existencia del estado inicial y de la derivación operacional
exigidos por la caracterización del Lema~\ref{lem:caracterizacion}:

$$\text{\lstinline|il_triple pre p post|} \implies \forall \text{\lstinline|s1|} \in \text{\lstinline|post|}. \
\exists \text{\lstinline|s0|} \in \text{\lstinline|pre|}. \ \text{\lstinline|runsto p s0 s1|}.$$
Por lo tanto,
el Corolario~\ref{cor:errores} también se aplica a esta formalización,
ya que si una derivación produce un resultado de error satisfacible,
existe un estado inicial y una ejecución operacional que alcanza dicho error.

Esta garantía debe entenderse respecto del lenguaje y de la semántica formalizados.
\fstar{} certifica que el sistema deductivo no introduce estados finales sin respaldo operacional dentro del modelo definido.
Sobre esta base,
el próximo capítulo extenderá el estado con memoria dinámica y trasladará la misma estrategia de mecanización a la ISL.

%%%%%%%%%%%%%%%%%%%%%%%%%%%%%%%%%%%%%%%%%%%%%%%%%%%%%%%%%%%%%%%%%%%%%%%%%%%%%%%%%%%%
%%%%%%%%%%%%%%%%%%%%%%%%%%%%%%%%%%%% CAPÍTULO 4 %%%%%%%%%%%%%%%%%%%%%%%%%%%%%%%%%%%%
%%%%%%%%%%%%%%%%%%%%%%%%%%%%%%%%%%%%%%%%%%%%%%%%%%%%%%%%%%%%%%%%%%%%%%%%%%%%%%%%%%%%

\chapter{Formalización de ISL}
\label{cap:formalizacion-isl}
\begin{quote}
``The second, to divide each of the difficulties I examined into as many parts as possible and as may be required in order to resolve them better.'' \\
\raggedleft --- René Descartes
\end{quote}

En el capítulo anterior se mecanizó la IL sobre un lenguaje imperativo sin memoria dinámica.
Sin embargo,
el modelo de memoria que se utilizó es insuficiente para analizar programas realistas.
Al introducir referencias de memoria dinámica surge el problema de \emph{aliasing}.

En este capítulo presentamos la formalización mecanizada de la ISL.
No se reexplicarán los fundamentos de la subaproximación,
los cuales se mantienen idénticos a los del capítulo anterior,
sino que detallaremos cómo la arquitectura de nuestra formalización en \fstar{} debió
evolucionar para incorporar el razonamiento local sobre porciones disjuntas de memoria,
la conjunción separadora y la detección de \emph{bugs} de memoria
como el \emph{use after free}.

\section{Extensión del lenguaje y del modelo de estado}
\label{sec:isl-estado}
Para soportar el razonamiento local propuesto por \citet{sl} en la SL,
y su posterior extensión en la ISL,
debemos abandonar la premisa de que un programa conoce y opera sobre la totalidad de la memoria.
En el \autoref{cap:formalizacion-il},
el estado estaba formado por un \emph{store} y un modo de terminación.
Ese modelo era suficiente para razonar sobre asignaciones,
no determinismo y control de flujo,
pero no permitía representar memoria dinámica ni \emph{aliasing}.

Como se vio en la \autoref{sec:logsep},
la memoria se representa mediante dos componentes:
el \emph{store},
que asocia variables con valores,
y el \emph{heap},
que describe el contenido de las ubicaciones de memoria.
En consecuencia,
el estado utilizado para la formalización de la ISL se redefine como una terna formada por un \emph{store},
un \emph{heap} y un modo de terminación.

En primer lugar,
se incorporan las ubicaciones de memoria y se amplía el tipo de los valores:

\begin{lstlisting}
type var = string
type loc = nat
type value =
  | Nat of nat
  | Loc of loc
type store = var ^-> value
\end{lstlisting}

A diferencia del \autoref{cap:formalizacion-il},
donde todos los valores eran números naturales,
ahora se distingue entre valores numéricos y direcciones de memoria.
El constructor \lstinline|Nat n| representa el valor natural \lstinline|n|,
mientras que \lstinline|Loc l| representa la ubicación \lstinline|l|.
La dirección 0 se reserva para modelar el puntero nulo;
los accesos exitosos a memoria exigen direcciones no nulas.

La Definición~\ref{def:heap_ext} presentó los \emph{heaps} de la ISL como funciones parciales
cuyo codominio incluye tanto valores como una marca negativa $\perp$.
Para modelar esta parcialidad en \fstar{} se decidió codificar el \emph{heap} como una función total hacia un nuevo tipo de dato \lstinline|cell|,
el cual explicita las tres situaciones posibles de una celda:

\begin{lstlisting}
type cell =
  | Full of value
  | Empty
  | Unknown
type heap = loc ^-> cell
\end{lstlisting}

Cada uno de los constructores de \lstinline|cell| define cuál es el estado actual de la celda de memoria:
\begin{enumerate}
  \item \lstinline|Full v|:
  indica que la ubicación es conocida,
  está asignada y contiene el valor $v$.
  \item \lstinline|Empty|:
  indica que la ubicación es conocida y no está asignada;
  en particular,
  permite representar una celda que fue liberada y todavía no fue reasignada.
  \item \lstinline|Unknown|:
  indica que el fragmento local del \emph{heap} no tiene información sobre esa dirección.
\end{enumerate}

El constructor \lstinline|Empty| constituye la representación mecanizada del elemento negativo $\perp$ introducido en la Definición~\ref{def:heap_ext}.
Por otro lado,
\lstinline|Unknown| permite representar la parcialidad del \emph{heap} dentro de una función total.
A diferencia del modelo matemático de la \autoref{sec:logsep},
el tipo \lstinline|heap| no impone que el conjunto de ubicaciones distintas de \lstinline|Unknown| sea finito.

Es importante distinguir \lstinline|Empty| de \lstinline|Unknown|. 
Una ubicación marcada como \lstinline|Empty| pertenece al fragmento conocido de memoria y registra que no está asignada.
En cambio,
una ubicación marcada como \lstinline|Unknown| queda fuera de ese fragmento
y puede pertenecer a otro \emph{heap} disjunto agregado mediante un marco.
Como se explicó en la \autoref{sec:logincsep},
esta distinción no aparece en el modelo clásico de \emph{heaps} de la SL,
pero resulta relevante para el razonamiento de la ISL.

En la semántica operacional,
una liberación exitosa transforma una celda \lstinline|Full v| en \lstinline|Empty|,
pero no la convierte en \lstinline|Unknown|.
De esta manera,
la formalización conserva la información de que la ubicación es conocida incluso después de haber sido liberada.
Esta decisión de representación se corresponde con la intuición de monotonía presentada en el Lema~\ref{lem:monotonia},
donde una ubicación conocida no deja de ser conocida como consecuencia de una liberación.

El tipo de estado queda definido como:

\begin{lstlisting}
type state = store & heap & term_mode
type cond = state -> prop
\end{lstlisting}

%La formalización no impone un único \emph{heap} inicial:
%el fragmento de memoria considerado depende de la presunción de cada tripleta.
%Un \emph{heap} que asigna \lstinline|Unknown| a todas las ubicaciones representa un fragmento vacío,
%mientras que \lstinline|Empty| expresa información conocida sobre una ubicación no asignada.
%\guido{el heap inicial está en Empty o en Unknown?}
%\guido{entiendo.. pero en general hay un estado inicial cuando
%uno define una semántica operacional. supongo que es todo empty,
%si fuera todo unknown no se puede alocar, no?}
%\guido{hmmm la regla no requiere Empty...}

Sobre este modelo de estado,
el lenguaje conserva los comandos del \autoref{cap:formalizacion-il} e incorpora cuatro operaciones de memoria:
\lstinline|Alloc| para reservar una celda,
\lstinline|Free| para liberarla,
\lstinline|Load| para leer su contenido y \lstinline|Store| para modificarlo.
La extensión se representa mediante:

\begin{lstlisting}
type stmt =
  ... // Se omiten los constructores ya presentados. 
  | Alloc : var -> stmt
  | Free : expr -> stmt
  | Load : var -> expr -> stmt
  | Store : expr -> expr -> stmt
\end{lstlisting}

Esta extensión también requiere adaptar la evaluación de expresiones a los nuevos valores.
La función \lstinline|eval_expr| sigue devolviendo un natural:
para una variable,
devuelve \lstinline|n| si contiene \lstinline|Nat n| y \lstinline|l| si contiene \lstinline|Loc l|.
Para conservar el valor de una variable al copiarlo mediante una asignación o una escritura en memoria,
se incorpora \lstinline|eval_value|:

\begin{lstlisting}
let eval_value (st : store) (e : expr) : GTot value =
  match e with
  | Var x -> st x
  | _ -> Nat (eval_expr' st e)
\end{lstlisting}
Cuando la expresión es una variable,
esta función conserva su constructor.
Para las demás expresiones devuelve el resultado numérico mediante \lstinline|Nat|.
Luego,
copiar una variable que contiene \lstinline|Loc l| conserva ese valor.
La evaluación de las direcciones de acceso y de las comparaciones sigue utilizando \lstinline|eval_expr'|,
por lo que esta adaptación no impone una separación de tipos entre números y punteros.

\section{Álgebra de \emph{heaps} y aserciones espaciales}
\label{sec:isl-aserciones}
Para representar la conjunción separadora de la Definición~\ref{def:sep-conj}
es necesario definir cuándo dos \emph{heaps} representan fragmentos independientes y cómo puede construirse su unión.
Para el primer caso, como los \emph{heaps} se representan como funciones totales,
la condición de que dos \emph{heaps} sean disjuntos se define punto a punto.
Dos celdas son disjuntas cuando al menos una de ellas no contiene información local:

\begin{lstlisting}
let cell_disjoint (c1 c2 : cell) : prop =
  c1 == Unknown \/ c2 == Unknown
\end{lstlisting}

Luego,
dos \emph{heaps} son disjuntos cuando sus celdas son disjuntas para toda ubicación:

\begin{lstlisting}
let heaps_disjoint (h1 h2 : heap) : prop =
  forall l. cell_disjoint (h1 l) (h2 l)
\end{lstlisting}

Esta definición mecaniza la relación $h_1 \perp h_2$ de la Definición~\ref{def:disj_heaps}.
En el modelo matemático del \autoref{cap:conceptos-previos},
la condición de que dos \emph{heaps} sean disjuntos se expresaba mediante la separación de los dominios de dos funciones parciales.
En la representación de \fstar{},
donde estas funciones se han totalizado,
se exige que no exista una ubicación conocida simultáneamente por ambos \emph{heaps}.

Definimos además el operador infijo \texttt{><} como abreviatura de \lstinline|heaps_disjoint|,
que en los fragmentos de código representamos mediante $\perp$.

Para combinar dos celdas se define \lstinline|cell_merge|,
que conserva la información de la primera cuando esta es conocida y utiliza la segunda cuando la primera es \lstinline|Unknown|.
La operación se extiende a los \emph{heaps} punto a punto:

\begin{lstlisting}
let cell_merge (c1 c2 : cell) : cell =
  match c1 with
  | Unknown -> c2
  | _ -> c1

let heap_merge (h1 h2 : heap) : heap =
  on loc (fun l -> cell_merge (h1 l) (h2 l))

let heap_union (h1 h2 : heap { h1 >< h2 }) : heap =
  heap_merge h1 h2
\end{lstlisting}

La operación \lstinline|heap_merge| no exige que sus argumentos sean disjuntos:
si ambos describen una ubicación,
conserva la información del primero.
En cambio,
el refinamiento de \lstinline|heap_union| exige que los \emph{heaps} sean disjuntos,
por lo que en ese caso la combinación representa la unión disjunta.
Bajo las condiciones de disjunción correspondientes,
esta unión es conmutativa y asociativa.

De manera análoga a lo hecho para \lstinline|heaps_disjoint|,
definimos el operador infijo \texttt{<+>} como abreviatura de \lstinline|heap_union| y lo representamos mediante $\uplus$ en los fragmentos de código.

Para construir las especificaciones de los programas,
formalizamos los tres predicados espaciales presentados en las Definiciones \ref{def:heap_vacio},
\ref{def:points_to} y \ref{def:heap_negativo}.
El más simple es \lstinline|emp|,
que describe un fragmento vacío,
donde todas las ubicaciones se representan mediante \lstinline|Unknown|:

\begin{lstlisting}
let emp : cond =
  fun (st, hp, m) -> forall l. hp l == Unknown
\end{lstlisting}

Por otro lado,
cuando un programa reserva memoria o accede a un puntero, necesitamos describir celdas específicas.
Para esto definimos los predicados \lstinline|points_to| y \lstinline|points_to_empty|:

\begin{lstlisting}
let points_to (l : loc) (v : value) : cond =
  fun (st, hp, m) -> l =!= 0 /\ hp l == Full v /\
                     forall l'. (l' <> l) ==> (hp l' == Unknown)

let points_to_empty (l : loc) : cond =
  fun (st, hp, m) -> l =!= 0 /\ hp l == Empty /\
                     forall l'. (l' <> l) ==> (hp l' == Unknown)
\end{lstlisting}
%\guido{serviría usar sintaxis de operadores? \lstinline`points_to_empty` podría
%llamarse \lstinline`is_empty` o algo mas claro?}

El predicado \lstinline|points_to l v| describe un \emph{heaplet} formado exclusivamente por una celda activa en la ubicación no nula \lstinline|l|,
con contenido \lstinline|v|.
La condición sobre las demás ubicaciones exige que todas sean \lstinline|Unknown|,
lo que corresponde a la interpretación exacta del \emph{points-to} de la Definición~\ref{def:points_to}.
Por su parte,
\lstinline|points_to_empty l| describe un \emph{heaplet} cuya única ubicación conocida es \lstinline|l|,
también no nula,
marcada como \lstinline|Empty|.
Este predicado representa la aserción de \emph{heap} negativo de la Definición~\ref{def:heap_negativo}.

A partir del álgebra anterior,
la conjunción separadora se representa mediante:

\begin{lstlisting}
let sep_conj (p q : cond) : cond =
  fun (st, hp, m) -> exists h1 h2.
    h1 >< h2 /\
    hp == (h1 <+> h2) /\
    p (st, h1, m) /\ q (st, h2, m)

let ( ** ) = sep_conj
\end{lstlisting}

Esta definición constituye la traducción directa de la Definición~\ref{def:sep-conj}.
Para que un estado satisfaga \lstinline|p ** q|,
deben construirse dos \emph{heaps} disjuntos cuya unión sea el \emph{heap} original,
de forma que \lstinline|p| se satisfaga sobre el primer fragmento y \lstinline|q| sobre el segundo.
Ambas aserciones se evalúan con el mismo \emph{store} y el mismo modo de terminación;
solo se divide el \emph{heap}.

A partir de la conjunción separadora,
se define \lstinline|points_to_flex|,
que permite describir una celda sin exigir que sea la única del fragmento considerado:

\begin{lstlisting}
let points_to_flex (l : loc) (v : value) : cond =
  points_to l v ** (fun _ -> True)
\end{lstlisting}
La aserción \lstinline|points_to l v| describe la celda,
mientras que \lstinline|True| permite cualquier contenido en el fragmento restante,
disjunto de ella.
De esta forma,
el \emph{heap} puede contener otras ubicaciones conocidas.

\section{Semántica operacional y metateoría}
\label{sec:isl-semantica}
La relación \texttt{runsto} conserva la estructura introducida en la \autoref{sec:il-sintaxis}.
La diferencia en este caso es que opera sobre estados formados por un \emph{store},
un \emph{heap} y un modo de terminación.
Los constructores correspondientes a asignación,
no determinismo, secuencia, elección y estrella de Kleene mantienen el mismo comportamiento general.
La extensión principal consiste en definir el comportamiento operacional de las operaciones \lstinline|Alloc|,
\lstinline|Free|, \lstinline|Load| y \lstinline|Store|,
incorporadas al lenguaje en la \autoref{sec:isl-estado}.

Como tanto el \emph{store} como el \emph{heap} se representan mediante funciones,
la equivalencia de estados del \autoref{cap:formalizacion-il} se amplía para comparar ambos componentes punto a punto:

\begin{lstlisting}
let state_equiv (s1 s2 : state) : prop =
  (forall x. s1._1 x == s2._1 x) /\
  (forall l. s1._2 l == s2._2 l) /\
  s1._3 == s2._3
\end{lstlisting}
Al igual que en el \autoref{cap:formalizacion-il},
la extensionalidad funcional permite obtener la igualdad de los estados a partir de esta equivalencia.
El lema \lstinline|state_equiv_eq| aplica este principio tanto al \emph{store} como al \emph{heap}.

En este capítulo se utiliza una versión generalizada de \lstinline|override| que permite actualizar tanto el \emph{store} como el \emph{heap},
ya que ambos se representan mediante funciones.

A diferencia de lo visto en la \autoref{sec:il-sintaxis},
\lstinline|R_Assign| utiliza ahora \lstinline|eval_value| para actualizar la variable destino.
Además de este cambio,
presentamos los constructores operacionales de las nuevas operaciones de memoria.

Al igual que antes,
se utiliza \lstinline|ok_state| para restringir los estados iniciales al modo \lstinline|Ok|.
Luego,
la regla operacional de reserva es:

\begin{lstlisting}
| R_Alloc : s : ok_state -> #x : var -> l : loc { l =!= 0 } -> v : value ->
  #(squash (s._2 l == Unknown \/ s._2 l == Empty)) ->
  runsto (Alloc x) s (override s._1 x (Loc l), override s._2 l (Full v), s._3)
\end{lstlisting}
La ubicación seleccionada debe ser no nula.
Además,
debe ser una dirección sobre la cual el \emph{heap} no poseía información o una dirección conocida que no estaba asignada.
Admitir \lstinline|Unknown| permite reservar una celda desde un fragmento vacío.
Exigir únicamente \lstinline|Empty| requeriría que la ubicación ya estuviera descrita como no asignada.

La liberación exitosa se formaliza mediante:

\begin{lstlisting}
| R_Free : s : ok_state -> e : expr -> l : loc { l =!= 0 } -> v : value ->
  #(squash (eval_expr s e == l)) ->
  #(squash (s._2 l == Full v)) ->
  runsto (Free e) s (s._1, override s._2 l Empty, s._3)
\end{lstlisting}
La expresión debe evaluarse a una dirección no nula que contenga una celda activa.
La ejecución no elimina la ubicación del \emph{heap},
sino que transforma \lstinline|Full v| en \lstinline|Empty|.
De esta manera,
se preserva el recurso negativo requerido por la ISL.

Si la dirección
está marcada como \lstinline|Empty|,
la liberación termina con error.
Este caso permite representar un \emph{double free} cuando la celda ya había sido liberada.
El \emph{heap} se conserva y el modo de terminación cambia a \texttt{Er}:

\begin{lstlisting}
| R_FreeEr : s : ok_state -> e : expr ->
  #(squash (eval_expr s e =!= 0)) ->
  #(squash (s._2 (eval_expr s e) == Empty)) ->
  runsto (Free e) s (s._1, s._2, Er)
\end{lstlisting}

El intento de liberar la dirección nula se trata mediante un constructor separado:

\begin{lstlisting}
| R_FreeNull : s : ok_state -> e : expr ->
  #(squash (eval_expr s e == 0)) ->
  runsto (Free e) s (s._1, s._2, Er)
\end{lstlisting}

Las reglas de lectura siguen el mismo patrón.
Una lectura exitosa requiere una celda activa y actualiza el \emph{store} con su contenido:

\begin{lstlisting}
| R_Load : s : ok_state -> x : var -> e : expr -> v : value ->
  #(squash (eval_expr s e =!= 0)) ->
  #(squash (s._2 (eval_expr s e) == Full v)) ->
  runsto (Load x e) s (override s._1 x v, s._2, s._3)
\end{lstlisting}

Los casos de error se representan mediante \lstinline|R_LoadEr|,
cuando la dirección no nula está marcada como \lstinline|Empty|,
y \lstinline|R_LoadNull|,
cuando la expresión evalúa a la dirección nula.
En ambos casos se conservan el \emph{store} y el \emph{heap},
y el modo de terminación cambia a \lstinline|Er|.
Se omiten sus definiciones por ser análogas a \lstinline|R_FreeEr| y \lstinline|R_FreeNull|,
respectivamente.

La escritura exitosa requiere que la ubicación se encuentre activa y reemplaza su contenido:

\begin{lstlisting}
| R_Store : s : ok_state -> e1 : expr -> e2 : expr ->
  l : loc { l =!= 0 } -> v : value ->
  #(squash (s._2 l == Full v)) ->
  #(squash (eval_expr s e1 == l)) ->
  runsto (Store e1 e2) s (s._1, override s._2 l (Full (eval_value s._1 e2)), Ok)
\end{lstlisting}

Los constructores \lstinline|R_StoreEr| y \lstinline|R_StoreNull| representan,
respectivamente, una escritura sobre una celda marcada como \lstinline|Empty| y una escritura sobre la dirección nula.
Los omitimos por ser análogos.

Estos constructores distinguen los accesos a celdas activas,
los accesos a ubicaciones marcadas como \lstinline|Empty| y los accesos a la dirección nula,
de acuerdo con la semántica de la \autoref{fig:semantica-isl}.
El tipo \lstinline|term_mode| solo distingue entre \lstinline|Ok| y \lstinline|Er|.
La causa específica del error queda registrada en el constructor utilizado para construir la derivación operacional.

La propiedad central del razonamiento local es que una ejecución sobre un fragmento de memoria
puede extenderse con un \emph{heap} disjunto que permanece inalterado.
Esta propiedad fue presentada de forma abstracta en el Lema~\ref{lem:localidad}.
En la formalización se mecaniza con la función \lstinline|r_frame|:

\begin{lstlisting}
let rec r_frame (#p : stmt) (#s0 : state) (#s1 : state)
  (r : runsto p s0 s1) (h_fr : heap)
  (#_ : squash (s0._2 >< h_fr)) (#_ : squash (s1._2 >< h_fr))
  : GTot (runsto p 
    (s0._1, s0._2 <+> h_fr, s0._3) (s1._1, s1._2 <+> h_fr, s1._3))
  (decreases r)
\end{lstlisting}

La función recibe la derivación operacional local y un \emph{heap} adicional \lstinline|h_fr|.
Las dos premisas que exigen \emph{heaps} disjuntos garantizan que el marco no se superponga con los \emph{heaps} inicial y final.
Como resultado se obtiene una nueva derivación:
\begin{center}
\lstinline|runsto p (s0| $\bullet$ \lstinline|h_fr) (s1| $\bullet$ \lstinline|h_fr)|,
\end{center}
donde $\bullet$ denota la extensión de un estado por un marco,
presentada en la Definición~\ref{def:state-ext}:
se conserva el \emph{store} y el modo de terminación,
y se agrega el marco al \emph{heap} mediante la unión disjunta $\uplus$.

La condición de que el marco sea disjunto del \emph{heap} final es especialmente relevante para la alocación.
Aunque una ubicación sea \lstinline|Unknown| en el \emph{heap} inicial,
al reservarla pasa a ser conocida y,
por lo tanto,
no puede pertenecer también al marco.
Si posteriormente se libera,
queda marcada como \lstinline|Empty| y continúa siendo conocida.
Por ello,
la condición del \emph{heap} final también impide que el marco contenga una ubicación reservada y luego liberada durante la ejecución.

La demostración se hace por inducción estructural sobre \lstinline|r|,
reconstruyendo en cada caso una derivación sobre los \emph{heaps} extendidos.
La mayoría de los casos se resuelven de manera directa.
El caso secuencial requiere,
además,
conocer que el marco también es disjunto del \emph{heap} intermedio.
Para obtener esta propiedad se utiliza el lema:

\begin{lstlisting}
let rec lemma_runsto_disjoint (#p : stmt) (#s0 : state) (#s1 : state)
  (h_fr : heap) (r : runsto p s0 s1) :
  Lemma (requires (s1._2 >< h_fr))
        (ensures (s0._2 >< h_fr)) (decreases r)
\end{lstlisting}

Este lema permite demostrar que el marco es disjunto del \emph{heap} inicial a partir de que lo es del final.
En el caso secuencial de \lstinline|r_frame|,
se aplica a la ejecución del segundo comando para demostrar que el marco también es disjunto del \emph{heap} intermedio.
Por lo tanto,
la condición de que el marco sea disjunto del \emph{heap} inicial es derivable,
aunque se mantiene explícita en la signatura de \lstinline|r_frame|.
%\guido{si tenes este lema no podes borrar una de las
%precondiciones de disjoint en r\_frame?}
%\guido{esta es la parte que me confunde de poder alocar
%en una direcciion que esta en Unknown, por cierto, porque el programa
%p podria alocar y liberarla, y en ese caso podria poner de marco
%algo que tambien contiene a esa direccion \textbf{pero} en realidad no
%porque queda en Empty... debe estar bien (la prueba anda) pero me suena raro}

La función \lstinline|r_frame| demuestra la dirección local a global de la propiedad de localidad.
Luego,
será reutilizada tanto en la prueba de solidez de la Regla del Marco como en la descomposición mediante huellas locales de la \autoref{sec:isl-huella}.

\section{La ISL como tipo inductivo}
\label{sec:isl-tipoind}
Estamos en condiciones de mecanizar el sistema deductivo de la ISL.
Al igual que en la \autoref{sec:il-tipoind},
representamos la relación de derivabilidad mediante un tipo inductivo en \fstar{},
al que denominaremos \lstinline|isl_triple|.
Su signatura mantiene la estructura presentada para la IL,
pero ahora las aserciones también pueden describir el \emph{heap}:

\begin{lstlisting}
type isl_triple : (pre : cond) -> (p : stmt) -> (post : cond) -> Type
\end{lstlisting}

Los constructores correspondientes a la asignación,
no determinismo, secuencia, elección, consecuencia, disyunción y estrella de Kleene mantienen la estructura presentada para la IL.
En particular,
la asignación adapta el valor de la variable destino mediante \lstinline|eval_value|,
de acuerdo con la semántica operacional.
Omitimos aquí los constructores compartidos con la IL y nos concentramos en las reglas de memoria dinámica y en la Regla del Marco.

Las reglas de memoria se presentan aquí con una presunción parametrizada \lstinline|pre|:
la postcondición exige que el estado anterior,
reconstruido a partir del estado final,
satisfaga esa presunción.
Esto explica su diferencia de forma respecto de las reglas de la \autoref{fig:reglas-isl}.

Comenzamos por la alocación,
cuyas dos reglas distinguen la información disponible sobre la ubicación reservada.
En \lstinline|ISL_Alloc1|,
la ubicación estaba marcada como \lstinline|Unknown|.
En este caso,
la postcondición describe una celda activa en una dirección no nula \lstinline|l|
y exige que la variable destino \lstinline|x| contenga \lstinline|Loc l|.
Para relacionar este resultado con la presunción \lstinline|pre|,
se recupera existencialmente el valor anterior de \lstinline|x| y se reconstruye el \emph{heap} inicial colocando \lstinline|Unknown| en \lstinline|l|:

\begin{lstlisting}
| ISL_Alloc1 : #pre : cond -> x : var ->
  isl_triple (is_ok pre) (Alloc x)
    (is_ok (fun s -> exists l v.
      s._1 x == Loc l /\
      l =!= 0 /\
      points_to l v s /\
      (exists x_old. pre (override s._1 x x_old, override s._2 l Unknown, Ok))))
\end{lstlisting}

La regla \lstinline|ISL_Alloc2| sigue el mismo esquema,
pero reconstruye el \emph{heap} inicial colocando \lstinline|Empty| en la dirección reservada.
De este modo,
describe la alocación sobre una ubicación previamente conocida como no asignada.
Al igual que en \lstinline|ISL_Alloc1|,
la dirección,
el contenido de la nueva celda y el valor anterior de la variable destino se cuantifican existencialmente en la postcondición:

\begin{lstlisting}
| ISL_Alloc2 : #pre : cond -> x : var ->
  isl_triple (is_ok pre) (Alloc x)
    (is_ok (fun s -> exists (l : loc) (v x_old : value).
      s._1 x == Loc l /\
      l =!= 0 /\
      points_to l v s /\
      pre (override s._1 x x_old, override s._2 l Empty, Ok)))
\end{lstlisting}

Para leer el contenido de una celda activa,
la regla \lstinline|ISL_Load| conserva el \emph{heap} y actualiza la variable destino \lstinline|x| con el contenido leído.
Su postcondición cuantifica la dirección \lstinline|l|,
el contenido \lstinline|v| y el valor anterior \lstinline|x_old|.
Este último permite reconstruir el \emph{store} inicial,
sobre el que se exige \lstinline|pre| y se evalúa la expresión \lstinline|e|.
Evaluar la dirección en ese \emph{store} es necesario porque \lstinline|e| puede depender de la propia variable \lstinline|x|:

\begin{lstlisting}
| ISL_Load : #pre : cond -> x : var -> e : expr ->
  isl_triple (is_ok pre) (Load x e)
    (is_ok (fun s -> exists l v x_old.
      eval_expr (override s._1 x x_old, s._2, Ok) e == l /\
      points_to l v s /\
      s._1 x == v /\
      pre (override s._1 x x_old, s._2, Ok)))
\end{lstlisting}

Si la dirección consultada corresponde a una celda negativa,
la lectura termina en error.
La regla \lstinline|ISL_LoadEr| expresa este caso mediante \lstinline|points_to_empty|
y conserva tanto el \emph{store} como el \emph{heap}.
Por ello,
para recuperar el estado inicial basta con sustituir el modo \lstinline|Er| por \lstinline|Ok|
y exigir que el estado resultante satisfaga \lstinline|pre|:

\begin{lstlisting}
| ISL_LoadEr : #pre : cond -> x : var -> e : expr ->
  isl_triple (is_ok pre) (Load x e)
    (fun s ->
      s._3 == Er /\
      points_to_empty (eval_expr s e) s /\
      pre (s._1, s._2, Ok))
\end{lstlisting}

La lectura también falla cuando la expresión \lstinline|e| evalúa a la dirección nula.
En este caso,
la regla \lstinline|ISL_LoadNull| conserva el mismo esquema de reconstrucción del estado inicial,
pero sustituye la condición sobre la celda negativa por la condición de nulidad:

\begin{lstlisting}
| ISL_LoadNull : #pre : cond -> x : var -> e : expr ->
  isl_triple (is_ok pre) (Load x e)
    (fun s -> s._3 == Er /\ eval_expr s e == 0 /\ 
      pre (s._1, s._2, Ok))
\end{lstlisting}

A diferencia de la lectura,
la escritura modifica el \emph{heap} y conserva el \emph{store}.
La postcondición de \lstinline|ISL_Store| exige que la celda indicada por \lstinline|e1|
contenga el valor obtenido al evaluar \lstinline|e2| mediante \lstinline|eval_value|.
Además,
cuantifica el contenido anterior \lstinline|v_old|,
de modo que al restaurarlo se obtenga un estado que satisface \lstinline|pre|:

\begin{lstlisting}
| ISL_Store : #pre : cond -> e1 : expr -> e2 : expr ->
  isl_triple (is_ok pre) (Store e1 e2)
    (is_ok (fun s ->
      points_to (eval_expr s e1) (eval_value s._1 e2) s /\
      (exists v_old. 
        pre (s._1, override s._2 (eval_expr s e1) (Full v_old), Ok))))
\end{lstlisting}

La liberación también modifica el \emph{heap},
pero transforma una celda activa en una celda negativa.
La postcondición de \lstinline|ISL_Free| exige que la ubicación liberada esté marcada como \lstinline|Empty|
y que exista un valor \lstinline|v| tal que,
al restaurar la celda como \lstinline|Full v|,
se obtenga un estado que satisface \lstinline|pre|.
El cuantificador existencial permite que el contenido anterior sea distinto para cada estado final:

\begin{lstlisting}
| ISL_Free : #pre : cond -> e : expr ->
    isl_triple (is_ok pre) (Free e)
      (is_ok (fun s -> exists (v : value).
        points_to_empty (eval_expr s e) s /\
        pre (s._1, override s._2 (eval_expr s e) (Full v), Ok)))
\end{lstlisting}

Las reglas de error de la escritura y liberación siguen el esquema presentado para la lectura:
conservan el \emph{store} y el \emph{heap},
y distinguen entre el acceso a una celda negativa y el acceso a la dirección nula.

Además de las reglas de memoria,
la regla \textsc{Exist} de la \autoref{fig:reglas-isl} se representa mediante \lstinline|ISL_Exists|.
Esta permite cuantificar existencialmente una variable del \emph{store} en la presunción y postcondición,
siempre que dicha variable no aparezca en el comando.
El predicado \lstinline|fv_stmt p x| indica que \lstinline|x| aparece libre en \lstinline|p|.
Para expresar la cuantificación,
\lstinline|exists_var x c| establece que existe un valor para \lstinline|x| con el cual se satisface \lstinline|c|,
manteniendo el resto del estado:

\begin{lstlisting}
let exists_var (x : var) (c : cond) : cond =
  fun s -> exists v. c (override s._1 x v, s._2, s._3)

| ISL_Exists : #p : stmt -> #pre : cond -> #post : cond -> x : var ->
  isl_triple (is_ok pre) p post ->
  squash (~(fv_stmt p x)) ->
  isl_triple (is_ok (exists_var x pre)) p (exists_var x post)
\end{lstlisting}
La condición sobre \lstinline|fv_stmt| garantiza que el comando no lee ni modifica la variable cuantificada.
Esto permite omitir su valor concreto en las aserciones sin alterar el comportamiento del comando.

Por otro lado, la Regla del Marco debe preservar no solo el \emph{heap} adicional,
sino también la validez de la aserción enmarcada respecto del \emph{store}.
Aunque ese \emph{heap} no cambie,
la aserción podría dejar de cumplirse si el comando modifica una variable de la que depende.
Para expresar esta condición se calcula el conjunto de variables que un comando puede modificar:

\begin{lstlisting}
let rec modifies (p : stmt) (x : var) : prop =
  match p with
  | Assign y _ -> x = y
  | Nondet y -> x = y
  | Seq q r -> modifies q x \/ modifies r x
  | Choice q r -> modifies q x \/ modifies r x
  | Kleene q -> modifies q x
  | Alloc y -> x = y
  | Load y _ -> x = y
  | _ -> False
\end{lstlisting}

A diferencia de lo que ocurre con \lstinline|ISL_Exists|,
cuya condición lateral se refiere solo al comando y puede expresarse sintácticamente mediante \lstinline|fv_stmt|,
la condición de la Regla del Marco~\ref{eq:frame-isl} involucra también las variables libres del marco.
Como las aserciones se representan mediante un \emph{shallow embedding},
es decir,
como funciones arbitrarias del tipo \lstinline|state -> prop|,
no existe una sintaxis sobre la cual calcular esas variables.
Por eso,
la condición $\mathrm{mod}(C) \cap \mathrm{fv}(r) = \emptyset$ se reemplaza por su contraparte semántica.
En lugar de exigir que el marco no mencione las variables modificadas,
se exige que su verdad no dependa de ellas.
Para expresarlo,
definimos cuándo una aserción es semánticamente inmune a la modificación de ciertas variables.

Dos \emph{stores} coinciden fuera de un conjunto de variables cuando asignan los mismos valores a todas las variables no incluidas en dicho conjunto.
Basándonos en esto,
una aserción es independiente de un conjunto de variables si su evaluación lógica devuelve el mismo resultado
sin importar cómo cambien los valores de dichas variables:

\begin{lstlisting}
let match_except_vars (vars : string -> prop) (st1 st2 : store) : prop =
  forall x. ~(vars x) ==> st1 x == st2 x

let independent_on_vars (vars : string -> prop) (c : cond) : prop =
  forall (st1 st2 : store) (hp : heap) (m1 m2 : term_mode).
    match_except_vars vars st1 st2 ==> (c (st1, hp, m1) <==> c (st2, hp, m2))
\end{lstlisting}

La definición exige que la verdad de \lstinline|c| no cambie cuando se modifican las variables incluidas en \lstinline|vars|.
Además,
exige que la aserción del marco no dependa del modo de terminación,
ya que compara su validez para modos arbitrarios \lstinline|m1| y \lstinline|m2|.
Esta restricción es necesaria porque la Regla del Marco también se aplica a tripletas de error:
el marco se cumple en la presunción,
cuyo modo es \lstinline|Ok|,
y debe seguir cumpliéndose en una postcondición cuyo modo puede ser \lstinline|Er|.
No supone una pérdida respecto de la presentación de \citet{isl},
donde el modo de terminación forma parte de la tripleta y no de las aserciones.
Por otro lado,
la condición semántica es más general que la sintáctica,
ya que admite marcos que mencionan una variable modificada sin depender realmente de su valor.

Con esto, la Regla del Marco se codifica como el siguiente constructor:

\begin{lstlisting}
| ISL_Frame : #pre : cond -> #p : stmt -> #post : cond -> fr : cond ->
  isl_triple (is_ok pre) p post ->
  squash (independent_on_vars (modifies p) fr) ->
  isl_triple (is_ok (pre ** fr)) p (post ** fr)
\end{lstlisting}

El constructor \lstinline|ISL_Frame| representa la Regla del Marco~\ref{eq:frame-isl}.
A partir de una derivación para \lstinline|pre| y \lstinline|post|,
permite incorporar la aserción \lstinline|fr| a ambas mediante la conjunción separadora,
que exige que los respectivos fragmentos del \emph{heap} sean disjuntos.
La condición \lstinline|independent_on_vars (modifies p) fr| garantiza que
los cambios en las variables modificables por \lstinline|p| y en el modo de terminación no alteren la verdad de \lstinline|fr|.

El \emph{shallow embedding} tiene ventajas y desventajas.
Las implicaciones entre aserciones,
como las que exige \lstinline|ISL_Consequence|,
son proposiciones de \fstar{} que Z3 descarga automáticamente en la mayoría de los casos,
y las aserciones pueden usar toda la expresividad de la lógica de \fstar{}.
Como desventaja,
el lenguaje de aserciones no existe como objeto dentro del desarrollo,
por lo que no es posible razonar sobre él.
Por ejemplo,
no puede demostrarse la completitud de la lógica respecto de un lenguaje de aserciones concreto.

\section{El Teorema de Soundness}
\label{sec:isl-soundness}
El Teorema de Soundness~\ref{teo:sound-isl} establece que toda tripleta derivable de la ISL es semánticamente válida.
Su mecanización conserva la misma forma general que la función \lstinline|soundness| desarrollada para IL en la \autoref{sec:il-soundness}.

\begin{teorema}[Soundness de la ISL formalizada]
\label{teo:soundness-isl}
Para todo comando \lstinline|p|,
presunción \lstinline|pre| y postcondición \lstinline|post|,
si existe una derivación
$$\vdash_{\mathrm{ISL}} \itriple{\text{\lstinline|pre|}}{\text{\lstinline|p|}}{\text{\lstinline|post|}}$$
entonces, para todo estado final \lstinline|s1| que satisface \lstinline|post|,
existe un estado inicial \lstinline|s0| que satisface \lstinline|pre| y una ejecución operacional de \lstinline|p| desde \lstinline|s0| hasta \lstinline|s1|:
$$\forall \text{\lstinline|s1|}. \ \text{\lstinline|post s1|} \implies
\exists \text{\lstinline|s0|}. \ \text{\lstinline|pre s0|} \land \text{\lstinline|runsto p s0 s1|}.$$
\end{teorema}

De manera análoga a la IL,
el teorema se codifica en \fstar{} mediante una función con tipos dependientes.
La signatura mantiene la misma forma,
reemplazando \lstinline|il_triple| por \lstinline|isl_triple| y utilizando los estados y la semántica operacional de la ISL:

\begin{lstlisting}
let rec soundness
  (p : stmt) (pre : cond) (post : cond)
  (pf : isl_triple pre p post) (s1 : state { post s1 })
  : GTot (s0 : state { pre s0 } & runsto p s0 s1) (decreases pf)
\end{lstlisting}

La signatura constituye la misma traducción de la caracterización subaproximada utilizada en el \autoref{cap:formalizacion-il}:
\begin{itemize}
  \item \lstinline|s1 : state { post s1 }| representa un estado final que satisface la postcondición;
  \item \lstinline|s0 : state { pre s0 }| representa el estado inicial testigo;
  \item \lstinline|runsto p s0 s1| es la ejecución que conecta ambos estados.
\end{itemize}

La demostración de la función se realiza por inducción estructural sobre el árbol de derivación \lstinline|pf|.
Los constructores heredados de la IL siguen el argumento desarrollado en la \autoref{sec:il-soundness}.
Los nuevos casos corresponden a las operaciones de memoria,
a la cuantificación existencial y a la Regla del Marco.

Se presentan los casos de memoria en el mismo orden que sus reglas deductivas.
En los casos de éxito
la prueba utiliza los testigos existenciales de la postcondición para reconstruir un estado inicial.
Estos testigos se obtienen mediante \lstinline|indefinite_description_ghost|.
Por su parte,
\lstinline|exists_intro| permite justificar una propiedad existencial proporcionando un testigo.
Luego se aplica la regla operacional correspondiente y,
cuando es necesario,
se utiliza \lstinline|state_equiv_eq| para identificar el estado final construido con el estado \lstinline|s1| recibido.

En \lstinline|ISL_Alloc1|,
se obtiene la dirección reservada \lstinline|l|,
su contenido \lstinline|v| y el valor anterior \lstinline|x_old| de la variable destino.
Con ellos se reconstruye \lstinline|s0|,
restaurando \lstinline|x_old| en el \emph{store} y \lstinline|Unknown| en la ubicación reservada.
Desde ese estado,
\lstinline|R_Alloc| permite reservar la misma dirección con el contenido observado en \lstinline|s1|:

\begin{lstlisting}
| ISL_Alloc1 #pre #x -> 
  let (st1, hp1, m1) = s1 in
  let unfold p_lv (l_i : loc) (v_i : value) : prop =
    st1 x == Loc l_i /\ l_i =!= 0 /\ points_to l_i v_i s1
  in
  let l = indefinite_description_ghost loc (fun l ->
    exists v. p_lv l v) in
  let v = indefinite_description_ghost value (fun v ->
    p_lv l v) in
  let p_xold (x_o : value) : prop =
    pre (override st1 x x_o, override hp1 l Unknown, Ok)
  in
  let x_old = indefinite_description_ghost value p_xold in
  let st0 = override st1 x x_old in
  let hp0 = override hp1 l Unknown in
  let s0 : state = (st0, hp0, Ok) in
  let r = R_Alloc s0 #x l v #() in
  state_equiv_eq (override st0 x (Loc l), override hp0 l (Full v), Ok) s1;
  (|s0, r|)
\end{lstlisting}

El caso \lstinline|ISL_Alloc2| sigue la misma construcción,
pero restaura \lstinline|Empty| en la dirección reservada.
La postcondición garantiza que el estado inicial así reconstruido satisface \lstinline|pre|.
Como la ubicación está marcada como no asignada,
se puede aplicar nuevamente \lstinline|R_Alloc|:

\begin{lstlisting}
| ISL_Alloc2 #pre #x ->
  let (st1, hp1, m1) = s1 in
  let unfold p_lv (l_i : loc) (v_i : value) : prop =
    st1 x == Loc l_i /\ l_i =!= 0 /\ points_to l_i v_i s1
  in
  let l = indefinite_description_ghost loc (fun l ->
    exists v. p_lv l v) in
  let v = indefinite_description_ghost value (fun v ->
    p_lv l v) in
  let p_xold (x_o : value) : prop =
    pre (override st1 x x_o, override hp1 l Empty, Ok)
  in
  let x_old = indefinite_description_ghost value p_xold in
  let st0 = override st1 x x_old in
  let hp0 = override hp1 l Empty in
  let s0 : state = (st0, hp0, Ok) in
  let r = R_Alloc s0 #x l v in
  state_equiv_eq (override st0 x (Loc l), override hp0 l (Full v), Ok) s1;
  (|s0, r|)
\end{lstlisting}

Para la lectura,
primero se obtiene \lstinline|x_old| junto con la garantía de que existen una dirección y un contenido compatibles con ese valor anterior.
Luego se eligen \lstinline|l| y \lstinline|v| y se reconstruye el \emph{store} inicial,
conservando el \emph{heap}.
La postcondición garantiza que la evaluación de \lstinline|e| en ese estado conduce a la celda que contiene \lstinline|v|.
Así se aplica \lstinline|R_Load| y se identifica el estado resultante con \lstinline|s1|:

\begin{lstlisting}
| ISL_Load #pre #x #e ->
  let (st1, hp1, m1) = s1 in
  let p_xold (x_o : value) : prop =
    exists (l_i : loc) (v_i : value).
      eval_expr (override st1 x x_o, hp1, Ok) e == l_i /\
      points_to l_i v_i s1 /\ st1 x = v_i /\
      pre (override st1 x x_o, hp1, Ok)
  in
  let x_old = indefinite_description_ghost value p_xold in
  let unfold p_lv (l_i : loc) (v_i : value) : prop =
    eval_expr (override st1 x x_old, hp1, Ok) e == l_i /\ 
    points_to l_i v_i s1 /\ st1 x == v_i /\
    pre (override st1 x x_old, hp1, Ok)
  in
  let l = indefinite_description_ghost loc (fun l ->
    exists v. p_lv l v) in
  let v = indefinite_description_ghost value (fun v ->
    p_lv l v) in
  let st0 = override st1 x x_old in
  let hp0 = hp1 in
  let s0 : state = (st0, hp0, Ok) in
  exists_intro (fun v_i -> points_to l v_i s0) v;
  let r = R_Load s0 x e v in
  let st1' = override st0 x v in
  state_equiv_eq (st1', hp0, Ok) s1;
  (|s0, r|)
\end{lstlisting}

En el caso \lstinline|ISL_LoadEr| no es necesario recuperar valores anteriores,
ya que la lectura fallida conserva el \emph{store} y el \emph{heap}.
Se construye \lstinline|s0| con esos mismos componentes y modo \lstinline|Ok|.
La condición \lstinline|points_to_empty| de la postcondición permite aplicar \lstinline|R_LoadEr|
y obtener el estado final con modo \lstinline|Er|:

\begin{lstlisting}
| ISL_LoadEr #pre #x #e ->
  let (st, hp, m1) = s1 in
  let s0 = (st, hp, Ok) in
  let l = eval_expr s0 e in
  let r = R_LoadEr s0 x e in
  (|s0, r|)
\end{lstlisting}

El caso \lstinline|ISL_LoadNull| reconstruye el estado inicial de la misma manera.
En lugar de utilizar una celda negativa,
la prueba aprovecha la condición de que \lstinline|e| evalúa a cero para construir la ejecución mediante \lstinline|R_LoadNull|:

\begin{lstlisting}
| ISL_LoadNull #pre #x #e ->
  let (st, hp, m1) = s1 in
  let s0 = (st, hp, Ok) in
  let r = R_LoadNull s0 x e in
  (|s0, r|)
\end{lstlisting}

Para demostrar el caso \lstinline|ISL_Store|,
se identifica la dirección escrita y se obtiene un contenido anterior \lstinline|v_old| que permita satisfacer \lstinline|pre|.
El estado inicial se reconstruye colocando \lstinline|Full v_old| en esa dirección,
sin modificar el \emph{store}.
A continuación,
\lstinline|R_Store| escribe el valor de \lstinline|e2| y \lstinline|state_equiv_eq| permite identificar el estado resultante con \lstinline|s1|:

\begin{lstlisting}
| ISL_Store #pre #e1 #e2 -> 
  let (st1, hp1, m1) = s1 in
  let p_l (l_i : loc) : prop =
    eval_expr s1 e1 == l_i /\ points_to l_i (eval_value s1._1 e2) s1
  in
  let l = indefinite_description_ghost loc p_l in
  let p_vold (v_o : value) : prop =
    pre (st1, override hp1 l (Full v_o), Ok)
  in
  let v_old = indefinite_description_ghost value p_vold in
  let st0 = st1 in
  let hp0 = override hp1 l (Full v_old) in
  let s0 = (st0, hp0, Ok) in
  exists_intro (fun v_i -> eval_expr s0 e1 == l /\ points_to l v_i s0) v_old;
  exists_intro (fun l_i -> exists v_i. eval_expr s0 e1 == l_i /\ points_to l_i v_i s0) l;
  let r = R_Store s0 e1 e2 l v_old in
  let hp1' = override hp0 l (Full (eval_value s0._1 e2)) in
  state_equiv_eq (st0, hp1', Ok) s1;
  (|s0, r|)
\end{lstlisting}

La liberación utiliza una reconstrucción similar a la escritura:
se obtiene el contenido anterior \lstinline|v| y se restaura la celda activa \lstinline|Full v| en el \emph{heap} inicial.
Al aplicar \lstinline|R_Free|,
esa celda pasa a estar marcada como \lstinline|Empty|,
tal como exige la postcondición.
Finalmente,
\lstinline|state_equiv_eq| identifica el estado obtenido con \lstinline|s1|:


\begin{lstlisting}
| ISL_Free #pre e -> 
  let (st1, hp1, m1) = s1 in
  let l = eval_expr s1 e in
  let p_v (v_i : value) : prop = 
    m1 == Ok /\ points_to_empty l s1 /\ pre (st1, override hp1 l (Full v_i), Ok) 
  in
  let v = indefinite_description_ghost value p_v in
  let st0 = st1 in
  let hp0 = override hp1 l (Full v) in
  let s0 = (st0, hp0, Ok) in
  let r = R_Free s0 e l v in
  state_equiv_eq (st0, override hp0 l Empty, Ok) s1;
  (|s0, r|)
\end{lstlisting}

Los casos de error de la escritura y liberación se resuelven como los de lectura:
se conserva el \emph{store} y el \emph{heap} del estado final,
se restablece el modo inicial \lstinline|Ok| y se aplica la regla operacional de error correspondiente.

En el caso \lstinline|ISL_Exists|,
se elige un valor testigo de la cuantificación existencial en la postcondición
y se aplica la hipótesis inductiva al estado final con ese valor en la variable cuantificada.
Como la variable no aparece en el comando,
la ejecución obtenida puede adaptarse para conservar el valor de esa variable en el estado final original.
El estado inicial resultante satisface la presunción existencialmente cuantificada.

Para esta construcción,
\lstinline|set_var| actualiza una variable del \emph{store} sin modificar el \emph{heap} ni el modo de terminación.
Por otro lado,
el lema \lstinline|r_fresh| permite reemplazar por un mismo valor esa variable en ambos extremos de una ejecución,
siempre que no aparezca en el comando.
El caso se demuestra de la siguiente manera:

\begin{lstlisting}
| ISL_Exists #p #pre #post x pf_p _ ->
  let v = indefinite_description_ghost value (fun v -> post (set_var s1 x v)) in
  let s1_v = set_var s1 x v in
  let (|s0_v, r|) = soundness p (is_ok pre) post pf_p s1_v in
  let orig = s1._1 x in
  let s0 = set_var s0_v x orig in
  let r' = r_fresh r x orig #() in
  state_equiv_eq (set_var s1_v x orig) s1;
  state_equiv_eq (set_var s0 x (s0_v._1 x)) s0_v;
  exists_intro (fun w -> pre (set_var s0 x w)) (s0_v._1 x);
  (|s0, r'|)
\end{lstlisting}

El caso de la Regla del Marco constituye el mayor desafío lógico de la formalización,
ya que requiere combinar la ejecución local obtenida mediante la hipótesis inductiva con la propiedad operacional del marco demostrada en \lstinline|r_frame|.
Además de extender la ejecución con el \emph{heap} adicional,
es necesario justificar que el estado inicial reconstruido satisface también la aserción del marco.

En este caso,
\lstinline|p_cmd| es el comando local,
\lstinline|p_pre| y \lstinline|p_post| son sus aserciones,
y \lstinline|pf_p| es la derivación a la que se aplica la hipótesis inductiva.
La aserción del marco se denota por \lstinline|fr|.

Cuando la prueba fue construida mediante el marco,
sabemos lógicamente que el estado final \lstinline|s1| satisface la conjunción separadora \lstinline|p_post ** fr|.
Por la Definición~\ref{def:sep-conj} de la conjunción separadora,
el \emph{heap} final puede dividirse en dos fragmentos disjuntos:
\begin{center}
\lstinline|s1._2 = h_local <+> h_fr|,
\end{center}
donde \lstinline|p_post| se satisface sobre \lstinline|h_local| y \lstinline|fr| sobre \lstinline|h_fr|.

En la prueba se define esta información mediante un predicado auxiliar:

\begin{lstlisting}
let p_two (h_local : heap) (h_fr : heap) : prop =
  h_local >< h_fr /\
  s1._2 == (h_local <+> h_fr) /\
  p_post (s1._1, h_local, s1._3) /\ fr (s1._1, h_fr, s1._3)
in
\end{lstlisting}
La definición de \lstinline|sep_conj| garantiza la existencia de dos \emph{heaps} que satisfacen \lstinline|p_two|.
Primero se elige \lstinline|h_local| junto con la garantía de que existe un marco compatible;
luego se elige \lstinline|h_fr| para ese fragmento.
Ambas elecciones se realizan en el contexto fantasma de la prueba.

Una vez aislado el fragmento local del \emph{heap},
aplicamos la hipótesis inductiva de nuestro teorema:

\begin{lstlisting}
  let (|s0_local, r_local|) =
    soundness p_cmd (is_ok p_pre) p_post pf_p (s1._1, h_local, s1._3) in
\end{lstlisting}
Esta llamada devuelve el estado inicial local \lstinline|s0_local| y una derivación \lstinline|r_local| de la ejecución hasta el estado \lstinline|(s1._1, h_local, s1._3)|.
El \emph{heap} \lstinline|h_local| es el fragmento que satisface la postcondición local proporcionada por la conjunción separadora.
En esta parte de la prueba no se afirma todavía que dicho fragmento sea una huella minimal en el sentido de la Definición~\ref{def:huella-minima}.

El último paso es elevar esa certeza local al contexto global,
es decir,
agregar \lstinline|h_fr| al estado inicial local.
La partición de \lstinline|s1| garantiza que el marco es disjunto del \emph{heap} final.
Para demostrar que también es disjunto del \emph{heap} inicial,
se utiliza:

\begin{lstlisting}
lemma_runsto_disjoint h_fr r_local;
\end{lstlisting}
Luego,
se usa la propiedad operacional del marco:

\begin{lstlisting}
let r_global = r_frame r_local h_fr #() #() in
\end{lstlisting}
Esto produce una ejecución sobre los \emph{heaps} globales.
Finalmente,
tenemos que la condición \lstinline|independent_on_vars (modifies p_cmd) fr|
y el lema que caracteriza las variables modificadas por \lstinline|runsto|
permiten demostrar que la aserción del marco también se satisfacía sobre el \emph{store} inicial reconstruido.
Como los \emph{heaps} iniciales son disjuntos,
el estado inicial global satisface \lstinline|is_ok (p_pre ** fr)|,
completando el caso:

\begin{lstlisting}
| ISL_Frame #p_pre #p_cmd #p_post fr pf_p _ ->
  let unfold p_two (h_local h_fr : heap) : prop = 
    h_local >< h_fr /\ 
    s1._2 == (h_local <+> h_fr) /\ 
    p_post (s1._1, h_local, s1._3) /\ 
    fr (s1._1, h_fr, s1._3)
  in
  let h_local = indefinite_description_ghost heap (fun h_local ->
    exists h_fr. p_two h_local h_fr) in
  let h_fr = indefinite_description_ghost heap (fun h_fr ->
    p_two h_local h_fr) in
  let (|s0_local, r_local|) = 
    soundness p_cmd (is_ok p_pre) p_post pf_p (s1._1, h_local, s1._3) in
  let (st0, hp0, m0) = s0_local in
  lemma_runsto_disjoint h_fr r_local;
  let s0 : state = (st0, hp0 <+> h_fr, m0) in
  let r_global = r_frame r_local h_fr #() #() in
  lemma_runsto_modifies r_local;
  (|s0, r_global|)
\end{lstlisting}

En conclusión,
la verificación de la función \lstinline|soundness| por parte de \fstar{} mecaniza el Teorema~\ref{teo:sound-isl}.
Luego,
$$\text{\lstinline|isl_triple pre p post|} \implies \forall \text{\lstinline|s1|} \in \text{\lstinline|post|}. \
\exists \text{\lstinline|s0|} \in \text{\lstinline|pre|}. \ \text{\lstinline|runsto p s0 s1|}.$$
Por el Corolario~\ref{cor:errores-isl},
cuando \lstinline|post| describe estados en modo \lstinline|Er| y es satisfacible,
existe una ejecución concreta del modelo formalizado que alcanza uno de esos estados de error.

\section{El Teorema de la Huella}
\label{sec:isl-huella}
En la \autoref{sec:logincsep} se presentó el Teorema de la Huella~\ref{teo:huella},
que caracteriza la semántica operacional de un comando a partir de sus transiciones minimales y de la extensión de estas con marcos disjuntos.
El resultado dice:
$$[\![C]\!]_\varepsilon = \operatorname{frame}(\operatorname{foot}(C)_\varepsilon).$$

La dirección local a global se obtiene mediante \lstinline|r_frame|,
presentada anteriormente.
Si \lstinline|runsto p s0 s1| y un \emph{heap} es disjunto de los \emph{heaps} de \lstinline|s0| y \lstinline|s1|,
entonces ese \emph{heap} puede agregarse a ambos estados y se obtiene una nueva derivación de \lstinline|runsto|.
Sin embargo,
esta propiedad por sí sola no identifica qué parte de una ejecución es local ni demuestra que toda ejecución global admita una descomposición parecida.

Siguiendo este esquema de descomposición,
se introducen dos relaciones.
La primera es \lstinline|footprint|,
que describe inductivamente ejecuciones sobre fragmentos locales del \emph{heap}.
La segunda,
\lstinline|framed_footprint|,
representa su clausura por marcos disjuntos.
Primero se demuestra que esta clausura coincide con \lstinline|runsto|.
A partir de la descomposición obtenida,
se establece después la existencia de una transición minimal respecto de la extensión por marcos,
completando la caracterización del Teorema de la Huella~\ref{teo:huella}.

\subsection*{La relación de huellas locales}
La relación \lstinline|runsto| permite que los estados inicial y final tengan memoria adicional que no participa de la ejecución.
Por ejemplo,
una ejecución de \lstinline|Load x e| puede ocurrir en un \emph{heap} con varias celdas,
aunque esa ejecución solo consulta una única ubicación.
Para aislar el fragmento local utilizado por cada comando se define el siguiente tipo inductivo:

\begin{lstlisting}
type footprint : stmt -> state -> state -> Type0
\end{lstlisting}
%\guido{Qué pasa con el programa $C = skip + x := Load(1)$, con estado inicial $x = 0$ y heap ${1 \mapsto 0}$?
%Acá tengo que $C$ deja ese estado sin cambiar, pero si saco ${1 \mapsto 0}$ y me quedo con el heap vacío,
%\emph{también} tengo eso, porque una rama no avanza pero la otra (skip) me deja igual, así que también tengo
%esta transición con el heap vacío.
%La huella según esta relación es ese heaplet. Entonces no es minimal?}

Llamamos \emph{huella local} a un habitante de \lstinline|footprint p s0 s1|,
que representa una ejecución local de \lstinline|p| desde \lstinline|s0| hasta \lstinline|s1|.
A diferencia de las huellas minimales de la Definición~\ref{def:huella-minima},
una huella local no es necesariamente minimal.

El tipo está indexado por el comando y por sus estados inicial y final.
La diferencia con \lstinline|runsto| es que sus constructores restringen los \emph{heaps} a los fragmentos requeridos por cada operación.

Para representar los comandos que no acceden a memoria dinámica se define el \emph{heap} vacío:

\begin{lstlisting}
let empty_heap : heap =
  on loc (fun _ -> Unknown)
\end{lstlisting}
Todas sus ubicaciones son \lstinline|Unknown|,
por lo que el fragmento no contiene ninguna celda activa ni negativa.
Este \emph{heap} se utiliza en las huellas locales de instrucciones como \lstinline|Skip|,
\lstinline|Error|, \lstinline|Assign|,
\lstinline|Nondet| y \lstinline|Assume|,
cuyos efectos solo involucran el \emph{store} o el modo de terminación.
Por ejemplo,
la huella de \lstinline|Skip| se codifica como:

\begin{lstlisting}
| F_Skip : st : store ->
  footprint Skip (st, empty_heap, Ok) (st, empty_heap, Ok)
\end{lstlisting}
La instrucción no requiere recursos espaciales y conserva todos los componentes del estado local.

Para las operaciones de memoria se definen \emph{heaps} formados por una única celda:

\begin{lstlisting}
let singleton_empty_heap (l : loc) : heap =
  on loc (fun l' -> if l' = l then Empty else Unknown)

let singleton_full_heap (l : loc) (v : value) : heap =
  on loc (fun l' -> if l' = l then Full v else Unknown)
\end{lstlisting}
La primera describe exactamente una celda negativa,
mientras que la segunda describe exactamente una celda activa en una dirección específica.
Todas las demás direcciones quedan fuera del fragmento local.

La alocación distingue dos situaciones iniciales.
El constructor \lstinline|F_AllocFresh| reserva una ubicación que no estaba descrita por el fragmento local,
por lo que parte de \lstinline|empty_heap|.
En cambio,
\lstinline|F_AllocReuse| reutiliza una ubicación conocida y marcada como \lstinline|Empty|.
Ambos casos producen una celda activa:

\begin{lstlisting}
| F_AllocFresh : st : store -> x : var -> l : loc { l =!= 0 } -> v : value ->
  footprint (Alloc x) (st, empty_heap, Ok) 
  (override st x (Loc l), singleton_full_heap l v, Ok)

| F_AllocReuse : st : store -> x : var -> l : loc { l =!= 0 } -> v : value ->
  footprint (Alloc x) (st, singleton_empty_heap l, Ok)
  (override st x (Loc l), singleton_full_heap l v, Ok)
\end{lstlisting}

Una lectura exitosa requiere únicamente la celda que será consultada:

\begin{lstlisting}
| F_Load : st : store -> x : var -> e : expr ->
  l : loc { l =!= 0 } -> v : value ->
  #(squash (eval_expr' st e == l)) ->
  footprint (Load x e) (st, singleton_full_heap l v, Ok)
  (override st x v, singleton_full_heap l v, Ok)
\end{lstlisting}
El \emph{heap} inicial contiene una única celda activa en \lstinline|l|.
La lectura conserva esa celda y actualiza en el \emph{store} el valor de \lstinline|x|.

En cambio,
una lectura sobre una ubicación marcada como \lstinline|Empty| requiere una celda negativa:

\begin{lstlisting}
| F_LoadEr : st : store -> x : var -> e : expr -> l : loc { l =!= 0 } ->
  #(squash (eval_expr' st e == l)) ->
  footprint (Load x e) (st, singleton_empty_heap l, Ok)
  (st, singleton_empty_heap l, Er)
\end{lstlisting}
En este caso, la operación no modifica el \emph{heap},
pero cambia el modo de terminación.
La celda negativa permanece en el estado final como justificación espacial del error.

Las reglas correspondientes a la escritura y a la liberación de memoria siguen el mismo patrón.
Una liberación exitosa transforma un \emph{heap} de una única celda activa en otro de una única celda negativa:

\begin{lstlisting}
| F_Free : st : store -> e : expr ->
  l : loc { l =!= 0 } -> v : value ->
  #(squash (eval_expr' st e == l)) ->
  footprint (Free e) (st, singleton_full_heap l v, Ok)
  (st, singleton_empty_heap l, Ok)
\end{lstlisting}
Una escritura exitosa conserva la dirección,
pero reemplaza su contenido:

\begin{lstlisting}
| F_Store : st : store -> e1 : expr -> e2 : expr ->
  l : loc { l =!= 0 } -> v : value ->
  #(squash (eval_expr' st e1 == l)) ->
  footprint (Store e1 e2) (st, singleton_full_heap l v, Ok)
    (st, singleton_full_heap l (eval_value st e2), Ok)
\end{lstlisting}

El caso más complejo de la definición es la composición secuencial.
Las huellas locales de dos comandos consecutivos no tienen por qué usar los mismos recursos en el estado intermedio.
El primer comando puede usar celdas que el segundo no necesita.
Análogamente,
el segundo puede requerir celdas que el primero no utilizó.
Por lo tanto,
no sería suficiente exigir que el \emph{heap} final de la huella local del primer comando coincida con el \emph{heap} inicial de la huella local del segundo.

Para representar esta situación,
el constructor \lstinline|F_Seq| divide la memoria del estado intermedio en tres componentes,
\lstinline|h_common|, \lstinline|h_p| y \lstinline|h_q|.
El \emph{heap} \lstinline|h_common| contiene los recursos utilizados por los dos comandos.
El \emph{heap} \lstinline|h_p| contiene los recursos exclusivos de \lstinline|p|,
mientras que \lstinline|h_q| contiene los recursos exclusivos de \lstinline|q|.

La estructura principal del constructor es:

%\guido{usa un simbolo para heaps-disjoint}
\begin{lstlisting}
| F_Seq : p : stmt -> q : stmt -> st0 : store -> stm : store ->
  st1 : store -> h_common : heap ->
  h_p : heap { h_common >< h_p } ->
  h_q : heap { h_common >< h_q /\ h_p >< h_q } ->
  h0 : heap { h0 >< h_q } ->
  h1 : heap { h1 >< h_p } ->
  m1 : term_mode ->
  footprint p (st0, h0, Ok) (stm, h_common <+> h_p, Ok) ->
  footprint q (stm, h_common <+> h_q, Ok) (st1, h1, m1) ->
  footprint (Seq p q) (st0, h0 <+> h_q, Ok)
                      (st1, h1 <+> h_p, m1)
\end{lstlisting}

En los esquemas siguientes se muestra únicamente la componente del \emph{heap} de cada estado.
La huella local de \lstinline|p| comienza con \lstinline|h0| y termina usando los recursos comunes y los exclusivos del primer comando,
\begin{center}
\lstinline|h0| $\xrightarrow{\texttt{p}}$ \lstinline|h_common <+> h_p|.
\end{center}
Sin embargo,
para construir la ejecución completa de la secuencia,
los recursos exclusivos de \lstinline|q| deben estar presentes durante la ejecución de \lstinline|p|,
aunque este no los utilice.
Por eso,
\lstinline|h_q| se agrega como marco,
\begin{center}
\lstinline|h0 <+> h_q| $\xrightarrow{\texttt{p}}$ \lstinline|(h_common <+> h_p) <+> h_q|.
\end{center}
De manera simétrica,
la huella local de \lstinline|q| usa los recursos comunes y sus recursos exclusivos,
\begin{center}
\lstinline|h_common <+> h_q| $\xrightarrow{\texttt{q}}$ \lstinline|h1|.
\end{center}
Durante esta segunda ejecución,
\lstinline|h_p| permanece como marco,
\begin{center}
\lstinline|(h_common <+> h_q) <+> h_p| $\xrightarrow{\texttt{q}}$ \lstinline|h1 <+> h_p|.
\end{center}
Las condiciones que exigen \emph{heaps} disjuntos permiten demostrar que los dos estados intermedios globales son equivalentes,
\begin{center}
\lstinline|(h_common <+> h_p) <+> h_q == (h_common <+> h_q) <+> h_p|.
\end{center}
Por lo tanto,
las dos ejecuciones pueden componerse mediante la regla operacional \lstinline|R_Seq|.

La propagación de un error producido por el primer comando se representa con un constructor separado:

\begin{lstlisting}
| F_SeqEr : p : stmt -> q : stmt ->
  s : state { s._3 == Ok } -> t : state { t._3 == Er } ->
  footprint p s t ->
  footprint (Seq p q) s t
\end{lstlisting}
Como el primer comando termina en error,
el segundo no se ejecuta y no incorpora nuevos requisitos espaciales.

Como en el caso de \lstinline|runsto|,
se utiliza \lstinline|state_equiv_eq| para reutilizar una derivación de \lstinline|footprint| sobre estados equivalentes.
%\guido{esta no es admisible tambien?}

\subsection*{Corrección operacional de las huellas locales}
Una vez definida la relación local,
el primer resultado consiste en demostrar que toda huella local representa una ejecución operacional válida:

\begin{lstlisting}
let rec footprint_sound (p : stmt) (s0 : state) (s1 : state)
  (f : footprint p s0 s1)
  : GTot (runsto p s0 s1) (decreases f)
\end{lstlisting}

La demostración se hace por inducción estructural sobre \lstinline|f|.
Los casos primitivos se traducen al constructor correspondiente de \lstinline|runsto|.
En algunos de ellos es necesario identificar estados cuyos \emph{heaps} se construyen mediante expresiones distintas.
Por ejemplo,
en el caso \lstinline|F_Free|,
la regla operacional actualiza la celda liberada mediante \lstinline|override|,
mientras que la huella local describe el \emph{heap} final mediante \lstinline|singleton_empty_heap|:

\begin{lstlisting}
| F_Free st e l v ->
  let s0 = (st, singleton_full_heap l v, Ok) in
  let s1 = (st, override (singleton_full_heap l v) l Empty, Ok) in
  let r = R_Free s0 e l v in
  state_equiv_eq s1 (st, singleton_empty_heap l, Ok);
  r
\end{lstlisting}
Ambas expresiones describen el mismo \emph{heap}:
la ubicación \lstinline|l| está marcada como \lstinline|Empty| y las demás permanecen como \lstinline|Unknown|.
El lema \lstinline|state_equiv_eq| establece la igualdad de los estados,
lo que permite reutilizar la derivación \lstinline|r| con el estado final requerido por la huella local.

En el caso \lstinline|F_Seq|,
las hipótesis inductivas proporcionan las ejecuciones de los dos comandos.
Se extienden mediante \lstinline|r_frame| con \lstinline|h_q| y \lstinline|h_p|,
respectivamente,
según el esquema anterior.
Luego,
\lstinline|state_equiv_eq| identifica los estados intermedios y \lstinline|R_Seq| compone las ejecuciones.

\subsection*{Clausura por marcos}
El resultado anterior permite convertir una huella local en una ejecución operacional.
Para caracterizar también las ejecuciones que contienen memoria adicional,
se define la clausura por marcos de \lstinline|footprint|:

\begin{lstlisting}
type framed_footprint : stmt -> state -> state -> Type0 =
  | FF_Frame : p : stmt -> s0 : state -> s1 : state ->
    h_fr : heap -> fp : footprint p s0 s1 ->
    d0 : squash (s0._2 >< h_fr) ->
    d1 : squash (s1._2 >< h_fr) ->
    framed_footprint p (s0._1, s0._2 <+> h_fr, s0._3)
                       (s1._1, s1._2 <+> h_fr, s1._3)
\end{lstlisting}
\lstinline|FF_Frame| recibe una huella local y agrega el mismo \emph{heap} \lstinline|h_fr| a sus estados inicial y final.
Las dos premisas que exigen \emph{heaps} disjuntos garantizan que el marco no se superponga con los recursos locales de ninguno de los estados.

La corrección operacional de esta clausura se obtiene al convertir la huella local en una ejecución mediante
\lstinline|footprint_sound| y extenderla con el marco mediante \lstinline|r_frame|:

\begin{lstlisting}
let framed_footprint_sound (p : stmt) (s0 : state) (s1 : state) 
  (ff : framed_footprint p s0 s1)
  : GTot (runsto p s0 s1) (decreases ff) =
  match ff with
  | FF_Frame p s0 s1 h_fr f_p _ _ ->
    let r_p = footprint_sound p s0 s1 f_p in
    r_frame r_p h_fr #() #()
\end{lstlisting}

\subsection*{Descomposición de ejecuciones globales}
Para completar la caracterización es necesario demostrar la dirección inversa.
Dada una derivación operacional,
debe ser posible separar los recursos utilizados localmente por el comando de aquellos que permanecen como marco.
Esta propiedad se mecaniza con:

\begin{lstlisting}
let rec runsto_decompose (#p : stmt) (#s0 #s1 : state) (r : runsto p s0 s1)
  : GTot (framed_footprint p s0 s1) (decreases r)
\end{lstlisting}
La demostración procede por inducción estructural sobre \lstinline|r|.

El auxiliar \lstinline|frame_local| construye una derivación de \lstinline|framed_footprint| a partir de una huella local y un marco disjunto,
y utiliza \lstinline|state_equiv_eq| para identificar los estados resultantes con los de la ejecución original:

\begin{lstlisting}
let frame_local (#p : stmt) (#t0 #t1 : state) (fp : footprint p t0 t1)
  (s0 s1 : state) (h_fr : heap)
  (#_ : squash (t0._2 >< h_fr /\ t1._2 >< h_fr /\
                state_equiv (t0._1, t0._2 <+> h_fr, t0._3) s0 /\
                state_equiv (t1._1, t1._2 <+> h_fr, t1._3) s1))
  : GTot (framed_footprint p s0 s1) =
  state_equiv_eq (t0._1, t0._2 <+> h_fr, t0._3) s0;
  state_equiv_eq (t1._1, t1._2 <+> h_fr, t1._3) s1;
  FF_Frame p t0 t1 h_fr fp () ()
\end{lstlisting}

Con este auxiliar,
el caso de \lstinline|Skip| usa \lstinline|F_Skip| como huella local y conserva todo el \emph{heap} original como marco:

\begin{lstlisting}
| R_Skip s -> frame_local (F_Skip s._1) s0 s1 s._2
\end{lstlisting}

Las operaciones de memoria requieren separar del \emph{heap} global la celda utilizada.
Para ello se define:

\begin{lstlisting}
let heap_without (h : heap) (l : loc) : heap =
  on loc (fun l' ->
    if l' = l then Unknown
    else h l')
\end{lstlisting}
En el caso de la liberación,
por ejemplo,
la transición local describe el cambio de la celda ubicada en \lstinline|l|,
que pasa de \lstinline|Full v| a \lstinline|Empty|.
El resto del \emph{heap} se conserva como marco mediante \lstinline|heap_without s._2 l|:

\begin{lstlisting}
| R_Free s e l v ->
  frame_local (F_Free s._1 e l v) s0 s1 (heap_without s._2 l)
\end{lstlisting}

Los casos de acceso nulo utilizan \lstinline|empty_heap| como \emph{heap} local y conservan todo el \emph{heap} global en el marco.

Esta elección del marco satisface las condiciones de \lstinline|frame_local|:
el fragmento local y el marco son disjuntos,
y su unión reconstruye el \emph{heap} correspondiente.
Por ejemplo,
si \lstinline|h l = Full v|, entonces:
\begin{center}
\lstinline|singleton_full_heap l v >< heap_without h l|, y

\lstinline|(singleton_full_heap l v <+> heap_without h l) == h|.
\end{center}
El caso de una operación sobre una celda negativa se realiza de manera análoga.

Las reglas de elección no determinista y el caso inductivo de Kleene reutilizan el marco obtenido por la hipótesis inductiva
y transforman únicamente la huella local.
Esta transformación se expresa mediante \lstinline|map_framed_footprint|:

\begin{lstlisting}
let map_framed_footprint (#p #q : stmt) (#s0 #s1 : state)
  (transform : (#sl0 : state -> #sl1 : state -> footprint p sl0 sl1 ->
    GTot (footprint q sl0 sl1)))
  (ff : framed_footprint p s0 s1)
  : GTot (framed_footprint q s0 s1) (decreases ff) =
  match ff with
  | FF_Frame _ s0_local s1_local h_fr fp d0 d1 ->
    let fp' = transform fp in
    FF_Frame q s0_local s1_local h_fr fp' d0 d1
\end{lstlisting}

Por ejemplo,
en \lstinline|R_ChoiceL| se descompone recursivamente la ejecución de la rama elegida.
Luego se aplica \lstinline|F_ChoiceL| a la derivación local obtenida,
conservando su marco:

\begin{lstlisting}
| R_ChoiceL #p #q r_p ->
  let ff_p = runsto_decompose r_p in
  map_framed_footprint (fun #sl0 #sl1 fp_p ->
    footprint_starts_ok fp_p;
    F_ChoiceL p q sl0 sl1 fp_p) ff_p
\end{lstlisting}

En el caso de \lstinline|R_SeqEr|,
el primer comando termina en error y el segundo no se ejecuta.
La hipótesis inductiva proporciona la descomposición de la ejecución de \lstinline|p|.
Para adaptarla a \lstinline|Seq p q|,
se utiliza el auxiliar \lstinline|framed_seq_er|,
que aplica \lstinline|F_SeqEr| a la derivación local y conserva el mismo marco:

\begin{lstlisting}
| R_SeqEr #p #q #s #t r_p ->
  let ff_p = runsto_decompose r_p in
  framed_seq_er #p #q #s #t ff_p
\end{lstlisting}

El caso de \lstinline|R_Seq| es la parte central de la prueba.
Supongamos una ejecución operacional \lstinline|s0| $\xrightarrow{\texttt{p}}$ \lstinline|t| $\xrightarrow{\texttt{q}}$ \lstinline|s1|.
Las llamadas recursivas producen descomposiciones de ambas subejecuciones \lstinline|framed_footprint p s0 t| y \lstinline|framed_footprint q t s1|.
Cada descomposición divide el \emph{heap} del estado intermedio \lstinline|t| en un fragmento local y un marco.

Para acceder a estos componentes,
\lstinline|flatten_framed| extrae los estados locales,
el marco,
la huella local y las propiedades que relacionan su reconstrucción con los estados globales.
En el caso de \lstinline|p|,
llamamos \lstinline|p1| al estado local final y \lstinline|fr_p| a su marco.
Para \lstinline|q|,
llamamos \lstinline|q0| al estado local inicial y \lstinline|fr_q| a su marco.
Así,
ambas descomposiciones reconstruyen el mismo \emph{heap} intermedio:

\begin{center}
\lstinline|t._2 == (p1._2 <+> fr_p)|,

\lstinline|t._2 == (q0._2 <+> fr_q)|.
\end{center}

Para construir \lstinline|F_Seq| es necesario reconciliar estas dos particiones.
Como ambas reconstruyen el \emph{heap} de \lstinline|t|,
se obtiene:
\begin{center}
\lstinline|(p1._2 <+> fr_p) == (q0._2 <+> fr_q)|.
\end{center}

Esta igualdad permite aplicar una partición cruzada,
denominada \emph{cross split}.
Esta construcción subdivide los fragmentos de ambas particiones en cuatro componentes compatibles.
En general,
dados \emph{heaps} \lstinline|h1|,
\lstinline|h2|, \lstinline|h3| y \lstinline|h4| tales que
\begin{center}
\lstinline|(h1 <+> h2) == (h3 <+> h4)|,
\end{center}
la partición cruzada construye cuatro \emph{heaps} \lstinline|h13|,
\lstinline|h14|, \lstinline|h23| y \lstinline|h24|,
disjuntos dos a dos y tales que
\begin{center}
\lstinline|h1 == (h13 <+> h14)|,

\lstinline|h2 == (h23 <+> h24)|,

\lstinline|h3 == (h13 <+> h23)|,

\lstinline|h4 == (h14 <+> h24)|.
\end{center}
El predicado \lstinline|heaps_pairwise_disjoint| exige que los cuatro \emph{heaps} sean disjuntos dos a dos.
Para expresar las igualdades se reutiliza \lstinline|heap_merge|,
definida en la \autoref{sec:isl-aserciones}.
Las condiciones del registro garantizan que los fragmentos son disjuntos,
de modo que esta operación coincide allí con la unión disjunta.
Luego,
la estructura que almacena estos componentes se define mediante:

\begin{lstlisting}
type cross_split_witness (h1 h2 h3 h4 : heap) = {
  h13 : heap; h14 : heap; h23 : heap; h24 : heap;
  pairwise_disjoint : squash (heaps_pairwise_disjoint h13 h14 h23 h24);
  split_h1 : squash (forall l. heap_merge h13 h14 l == h1 l);
  split_h2 : squash (forall l. heap_merge h23 h24 l == h2 l);
  split_h3 : squash (forall l. heap_merge h13 h23 l == h3 l);
  split_h4 : squash (forall l. heap_merge h14 h24 l == h4 l)
}
\end{lstlisting}

En la composición secuencial,
se tienen \lstinline|dp1| y \lstinline|dq0|,
que justifican, respectivamente,
\lstinline|p1._2 >< fr_p| y \lstinline|q0._2 >< fr_q|.
Con estas pruebas,
la partición cruzada se obtiene mediante:

\begin{lstlisting}
let cs = heap_cross_split p1._2 fr_p q0._2 fr_q #dp1 #dq0 #() in
let h_common = cs.h13 in
let h_p = cs.h14 in
let h_q = cs.h23 in
let h_ext = cs.h24 in
\end{lstlisting}
Esta elección recupera la distribución de los recursos utilizada por \lstinline|F_Seq|:
en el estado intermedio,
\lstinline|h_common| contiene los recursos compartidos por ambas huellas locales,
mientras que \lstinline|h_p| y \lstinline|h_q| contienen los exclusivos de cada una.
El cuarto componente,
\lstinline|h_ext|,
pertenece a ambos marcos y se conserva como marco exterior de toda la secuencia.

A partir de esta partición,
\lstinline|state_equiv_eq| permite identificar los estados de las huellas locales con los requeridos por \lstinline|F_Seq|.
Este constructor combina ambas huellas locales,
y \lstinline|FF_Frame| incorpora \lstinline|h_ext| como marco exterior de la secuencia.
Finalmente,
se vuelve a utilizar \lstinline|state_equiv_eq| para identificar los estados reconstruidos con los de la derivación operacional original.

\subsection*{Resultado de descomposición}
Las funciones \lstinline|framed_footprint_sound| y \lstinline|runsto_decompose| establecen las dos direcciones de la caracterización.
Por lo tanto,
se obtiene el siguiente resultado:

\begin{teorema}
Para todo comando \lstinline|p| y estados \lstinline|s0| y \lstinline|s1|,
\begin{center}
\lstinline|runsto p s0 s1| $\iff$ \lstinline|framed_footprint p s0 s1|.
\end{center}
\end{teorema}
En \fstar{},
ambas direcciones se expresan mediante funciones que transforman una derivación de una relación en una derivación de la otra:

\begin{lstlisting}
let framed_decomposition_sound (#p : stmt) (#s0 #s1 : state)
  (ff : framed_footprint p s0 s1) : GTot (runsto p s0 s1) =
  framed_footprint_sound p s0 s1 ff

let framed_decomposition_complete (#p : stmt) (#s0 #s1 : state)
  (r : runsto p s0 s1) : GTot (framed_footprint p s0 s1) =
  runsto_decompose r
\end{lstlisting}

Esta caracterización proporciona una descomposición local,
pero no exige que la transición obtenida sea minimal respecto de la extensión por marcos de la Definición~\ref{def:huella-minima}.
Por ejemplo,
para el comando \lstinline|skip| $+ \ x := [1]$,
con $x = 0$ y \emph{heap} $1 \mapsto 0$,
la rama de lectura produce una derivación de \lstinline|footprint| que usa la celda 1 y deja el estado sin cambios.
Sin embargo,
la rama \lstinline|skip| produce la misma transición sobre el \emph{heap} vacío,
por lo que la primera no es minimal.
Cada constructor elige los recursos necesarios para su instrucción,
pero no puede tener en cuenta que otra rama del programa justifique la misma transición con menos memoria.
Por eso,
\lstinline|footprint| se usa como paso intermedio para demostrar a continuación la existencia de una transición minimal por debajo de cada ejecución.

\subsection*{Orden por marcos y transiciones minimales}
Para expresar la minimalidad como una propiedad lógica,
\lstinline|runs p s0 s1| afirma que existe una derivación operacional entre los estados.
Además,
\lstinline|known| identifica las celdas distintas de \lstinline|Unknown|
y \lstinline|nonempty| establece que un \emph{heap} contiene al menos una ubicación conocida:

\begin{lstlisting}
let runs (p : stmt) (s0 s1 : state) : prop =
  exists (_ : runsto p s0 s1). True

let known (c : cell) : bool = not (Unknown? c)
let nonempty (h : heap) : prop = exists l. known (h l)
\end{lstlisting}

A partir de estas definiciones,
el predicado \lstinline|framed_by| expresa que los estados \lstinline|s0| y \lstinline|s1|
se obtienen agregando el mismo marco a \lstinline|t0| y \lstinline|t1|.
Esta relación permite definir \lstinline|minimal| como una transición válida de la que no puede retirarse un marco no vacío
conservando una ejecución del mismo comando.
Finalmente,
\lstinline|has_min_below| afirma que existe una transición minimal cuya extensión reconstruye la transición original:

\begin{lstlisting}
let framed_by (t0 t1 s0 s1 : state) (hf : heap) : prop =
  t0._2 >< hf /\ t1._2 >< hf /\
  s0 == (t0._1, t0._2 <+> hf, t0._3) /\
  s1 == (t1._1, t1._2 <+> hf, t1._3)

let minimal (p : stmt) (s0 s1 : state) : prop =
  runs p s0 s1 /\
  (forall (t0 t1 : state) (hf : heap).
    (runs p t0 t1 /\ nonempty hf /\ framed_by t0 t1 s0 s1 hf) ==> False)

let has_min_below (p : stmt) (s0 s1 : state) : prop =
  exists (t0 t1 : state) (hf : heap).
    minimal p t0 t1 /\ framed_by t0 t1 s0 s1 hf
\end{lstlisting}

\subsection*{Finitud del fragmento local}
Para demostrar que existe una transición minimal,
se retiran marcos no vacíos mientras sea posible.
La finitud del fragmento local permite justificar que este proceso termina.

Aunque el tipo \lstinline|heap| no impone finitud,
el \emph{heap} inicial de una derivación de \lstinline|footprint| tiene un conjunto finito de ubicaciones conocidas.
Como las ubicaciones son números naturales,
esta propiedad se expresa mediante una cota:
a partir de cierto índice,
todas las celdas son \lstinline|Unknown|.
La función \lstinline|count| cuenta las ubicaciones conocidas por debajo de esa cota:

\begin{lstlisting}
let bounded (h : heap) (n : nat) : prop =
  forall (l : loc). l >= n ==> h l == Unknown

let rec count (h : heap) (n : nat) : GTot nat (decreases n) =
  if n = 0 then 0
  else count h (n - 1) + (if known (h (n - 1)) then 1 else 0)
\end{lstlisting}

La prueba de que el fragmento local es finito tiene la siguiente signatura:

\begin{lstlisting}
let rec footprint_bounded (#p : stmt) (#s0 #s1 : state) (f : footprint p s0 s1)
  : Lemma (ensures exists (n : nat). bounded s0._2 n) (decreases f)
\end{lstlisting}
La prueba procede por inducción sobre la huella local.
Los casos primitivos parten de un \emph{heap} vacío o una única celda,
para los que basta una cota igual a cero o a la dirección más uno.
En la composición secuencial,
se combinan las cotas de los fragmentos mediante \lstinline|union_bounded|,
que toma el máximo de ambas.
Los constructores de elección y de iteración reutilizan el resultado inductivo.
La propiedad demostrada se refiere al \emph{heap} inicial de la huella local y permite utilizar \lstinline|count| como medida finita en la prueba de minimalidad.

Con esta cota,
el lema \lstinline|minimize| demuestra la existencia de una transición minimal.
Para ello utiliza dos propiedades de la extensión por marcos:
\lstinline|empty_frame| establece que toda transición es extensión de sí misma por el \emph{heap} vacío,
y \lstinline|compose_frames| permite reemplazar la extensión por un marco seguida de la extensión por otro marco por una única extensión con la unión de ambos.
La signatura de \lstinline|minimize| es la siguiente:

\begin{lstlisting}
let rec minimize (p : stmt) (s0 s1 : state) (n : nat)
  : Lemma (requires runs p s0 s1 /\ bounded s0._2 n)
          (ensures has_min_below p s0 s1) (decreases (count s0._2 n))
\end{lstlisting}

La demostración distingue si la transición recibida ya es minimal.
En ese caso,
basta tomarla como testigo y utilizar el marco vacío.
En caso contrario,
la negación de minimalidad permite elegir una transición menor y un marco no vacío cuya extensión produce la original.
Ese marco contiene al menos una ubicación conocida que,
por ser disjunto del fragmento restante,
no es conocida en dicho fragmento.
Por ello,
al retirarlo disminuye estrictamente la cantidad de ubicaciones conocidas del \emph{heap} inicial,
como demuestra \lstinline|count_strict|.
La llamada recursiva obtiene una transición minimal por debajo de la transición menor
y \lstinline|compose_frames| combina los marcos para reconstruir la transición original.

\subsection*{Resultado de minimalidad}
La descomposición anterior permite aplicar \lstinline|minimize| al fragmento local acotado y extender el resultado a una ejecución cualquiera.
El lema \lstinline|footprint_theorem_minimal| establece que toda ejecución admite una transición minimal del comando que,
al extenderse con un marco disjunto de sus \emph{heaps} inicial y final,
reconstruye los estados originales:

\begin{lstlisting}
let footprint_theorem_minimal (#p : stmt) (#s0 #s1 : state) (r : runsto p s0 s1)
  : Lemma (has_min_below p s0 s1)
\end{lstlisting}

Para demostrarlo,
\lstinline|runsto_decompose| descompone la ejecución en una huella local y un marco.
Luego,
\lstinline|footprint_sound| proporciona la derivación operacional de la huella local,
y \lstinline|minimize| obtiene una transición minimal por debajo de ella.
Finalmente,
\lstinline|compose_frames| combina el marco retirado durante la minimización con el marco de la descomposición inicial,
de modo que la extensión de la transición minimal reconstruye la ejecución original.

Para completar la caracterización del Teorema de la Huella,
resta considerar la dirección inversa:
toda transición minimal extendida con un marco compatible es una ejecución del comando.
Esta propiedad se obtiene mediante \lstinline|r_frame| y se expresa en el lema \lstinline|minimal_framed_runs|:

\begin{lstlisting}
let minimal_framed_runs (p : stmt) (r0 r1 : state) (hf : heap)
  : Lemma (requires minimal p r0 r1 /\ r0._2 >< hf /\ r1._2 >< hf)
          (ensures runs p (r0._1, r0._2 <+> hf, r0._3) 
            (r1._1, r1._2 <+> hf, r1._3)) =
  let r = indefinite_description_ghost (runsto p r0 r1) (fun _ -> True) in
  exists_intro (fun (_ : runsto p (r0._1, r0._2 <+> hf, r0._3) 
    (r1._1, r1._2 <+> hf, r1._3)) -> True)
    (r_frame r hf #() #())
\end{lstlisting}

Combinando ambas direcciones obtenemos la versión mecanizada del Teorema de la Huella~\ref{teo:huella}.

\begin{teorema}(Teorema de la Huella formalizado)
Para todo comando \lstinline|p| y estados \lstinline|s0| y \lstinline|s1|,
\begin{center}
\lstinline|runs p s0 s1| $\iff$
\lstinline|has_min_below p s0 s1|.
\end{center}
\end{teorema}

\section{Completado de \emph{heaps} y certificados de error}
Los resultados anteriores permiten razonar sobre fragmentos de memoria que dejan algunas ubicaciones como \lstinline|Unknown|.
Este valor no tiene un equivalente en una máquina real:
indica que el fragmento no describe esa ubicación,
y que podría pertenecer a otro fragmento agregado mediante un marco.
En una ejecución concreta,
en cambio,
cada ubicación está asignada o no lo está.
Para relacionar las ejecuciones sobre fragmentos con ejecuciones sobre una memoria de este tipo,
se define una operación de completado,
que reemplaza cada celda \lstinline|Unknown| por \lstinline|Empty| y conserva las celdas conocidas.
Un \emph{heap} completado no tiene ubicaciones desconocidas,
sino que corresponde a una máquina cuya memoria asignada es exactamente la del fragmento,
y en la que todas las demás ubicaciones están libres:

\begin{lstlisting}
let complete_cell (c : cell) : cell =
  match c with
  | Unknown -> Empty
  | _ -> c

let complete (h : heap) : heap =
  on loc (fun l -> complete_cell (h l))

let cstate (s : state) : state =
  (s._1, complete s._2, s._3)
\end{lstlisting}

El estado \lstinline|cstate| aplica el completado al \emph{heap} de un estado sin modificar su \emph{store} ni su modo de terminación.
Este estado puede describirse mediante la conjunción separadora.
Sea \lstinline|no_alloc| la aserción que describe un fragmento sin celdas activas,
en el que toda ubicación conocida está liberada:

\begin{lstlisting}
let no_alloc : cond =
  fun s -> forall l. s._2 l == Empty \/ s._2 l == Unknown

let complete_sat (pre : cond) (s : state)
  : Lemma (requires pre s)
          (ensures (pre ** no_alloc) (cstate s)) =
  complete_split s._2;
  introduce_frame pre no_alloc (cstate s) s._2 (rest s._2)
\end{lstlisting}

El lema \lstinline|complete_sat| establece que,
si un estado satisface \lstinline|pre|,
su completado satisface \lstinline|pre ** no_alloc|:
el \emph{heap} completado se descompone en el fragmento original y un fragmento disjunto
que marca como \lstinline|Empty| todas las ubicaciones que el original no conocía.
Para demostrarlo,
se demuestra que \lstinline|complete h| coincide con la unión disjunta de \lstinline|h| y ese fragmento,
y se aplica \lstinline|introduce_frame|.

Para demostrar que el completado conserva las ejecuciones,
\lstinline|run_complete| demuestra,
por inducción sobre la derivación operacional original,
que el mismo comando puede ejecutarse entre los estados obtenidos al completar su \emph{heap}:

\begin{lstlisting}
let rec run_complete (#p : stmt) (#s0 #s1 : state) (r : runsto p s0 s1)
  : GTot (runsto p (cstate s0) (cstate s1)) (decreases r)
\end{lstlisting}

En el caso de la alocación,
\lstinline|R_Alloc| permite reservar una ubicación marcada como \lstinline|Unknown| o \lstinline|Empty|.
Por lo tanto,
si el completado reemplaza \lstinline|Unknown| por \lstinline|Empty|,
la misma ubicación sigue disponible para la alocación.

Para los casos que modifican el \emph{heap},
se utiliza el lema \lstinline|complete_override|.
Este establece que actualizar una celda con un valor distinto de \lstinline|Unknown|
y luego completar el \emph{heap} produce el mismo resultado que completar primero el \emph{heap} y realizar después esa actualización:

\begin{lstlisting}
let complete_override (st : store) (h : heap) (m : term_mode)
  (l : loc) (c : cell { known c })
  : Lemma ((st, complete (override h l c), m) == 
    (st, override (complete h) l c, m)) =
  state_equiv_eq (st, complete (override h l c), m)
    (st, override (complete h) l c, m)
\end{lstlisting}

Con estos resultados,
la solidez puede combinarse con un estado final concreto para certificar la existencia de un error.
La derivación lógica por sí sola no garantiza esa existencia si su postcondición es vacía.
Por eso se proporciona además un estado que satisface la postcondición y tiene modo \lstinline|Er|.
El predicado \lstinline|error_reachable| reúne lo que se certifica,
y la función \lstinline|certify| lo demuestra a partir de la derivación y del estado final testigo:

\begin{lstlisting}
let error_reachable (pre : cond) (p : stmt) : prop =
  exists (s0 s1 : state). pre s0 /\ s1._3 == Er /\ runs p s0 s1 /\
    (pre ** no_alloc) (cstate s0) /\
    (cstate s1)._3 == Er /\ runs p (cstate s0) (cstate s1)

let certify (#pre #post : cond) (#p : stmt)
  (pf : isl_triple pre p post) (s1 : state)
  : Lemma (requires post s1 /\ s1._3 == Er)
          (ensures error_reachable pre p) =
  let (| s0, r |) = soundness p pre post pf s1 in
  complete_sat pre s0;
  exists_intro (fun (_ : runsto p s0 s1) -> True) r;
  exists_intro (fun (_ : runsto p (cstate s0) (cstate s1)) -> True)
    (run_complete r);
  assert (pre s0 /\ s1._3 == Er /\ runs p s0 s1 /\
    (pre ** no_alloc) (cstate s0) /\
    (cstate s1)._3 == Er /\ runs p (cstate s0) (cstate s1))
\end{lstlisting}

Cada parte de \lstinline|error_reachable| certifica una propiedad distinta.
La primera,
\lstinline|pre s0 \land s1._3 == Er \land runs p s0 s1|,
sale del Teorema de Soundness:
existe una ejecución errónea del programa desde un estado que satisface la presunción.
Esta ejecución ocurre sobre el fragmento local,
por lo que puede apoyarse en ubicaciones \lstinline|Unknown|;
por ejemplo,
una alocación puede elegir una ubicación sobre la que el fragmento no tenía información.
La segunda parte proviene de \lstinline|complete_sat| y \lstinline|run_complete|,
y certifica que el mismo programa alcanza el error en la máquina en la que toda ubicación fuera del fragmento está libre.
El estado inicial de esta ejecución no satisface en general la presunción,
ya que esta suele exigir que las ubicaciones fuera del fragmento sean \lstinline|Unknown|;
lo que satisface es \lstinline|pre ** no_alloc|,
es decir,
la presunción junto con el resto de la memoria liberada.
Por último,
la condición \lstinline|(cstate s1)._3 == Er| se deduce de \lstinline|s1._3 == Er|,
ya que el completado no modifica el modo de terminación,
y se incluye para que el enunciado sea explícito.

El certificado cubre esta completación particular,
en la que toda la memoria fuera del fragmento está libre.
Para extender la ejecución a un contexto de memoria arbitrario se utiliza,
en cambio,
\lstinline|r_frame|,
que exige que la memoria agregada sea disjunta de los heaps inicial y final.

%%%%%%%%%%%%%%%%%%%%%%%%%%%%%%%%%%%%%%%%%%%%%%%%%%%%%%%%%%%%%%%%%%%%%%%%%%%%%%%%%%%%
%%%%%%%%%%%%%%%%%%%%%%%%%%%%%%%%%%%% CAPÍTULO 5 %%%%%%%%%%%%%%%%%%%%%%%%%%%%%%%%%%%%
%%%%%%%%%%%%%%%%%%%%%%%%%%%%%%%%%%%%%%%%%%%%%%%%%%%%%%%%%%%%%%%%%%%%%%%%%%%%%%%%%%%%

\chapter{Evaluación}
%\guido{En todos los ejemplos acá abajo, está claro que la post es satisfacible?}

\label{cap:evaluacion}
\begin{quote}
``The purpose of computing is insight, not numbers.'' \\
\raggedleft --- Richard W. Hamming
\end{quote}

En el \autoref{cap:formalizacion-isl} pudimos construir,
mecanizar y demostrar la solidez matemática de la ISL utilizando \fstar{}.
El objetivo de este capítulo es evaluar la expresividad de la formalización mediante una serie de casos de estudio representativos.

En particular,
se estudian un error de uso de memoria después de su liberación,
un puntero colgante generado por una reasignación,
un bucle con elecciones no deterministas y el recorrido de una lista enlazada.
Estos casos muestran que la formalización puede expresar y verificar derivaciones no triviales.
Además,
se presenta un experimento exploratorio en el que un agente basado en un modelo de lenguaje utiliza la formalización para construir una prueba de error.

Para cada caso,
además de la derivación,
se proporciona un estado final testigo y se demuestra la existencia del error,
no solo la validez de la tripleta.

%\guido{``La formalización y ejemplos se desarrollaron sobre \fstar{} version XXX
%y se encuentran en el repositorio XXX, commit XXX''}

\section{Certificación de errores temporales básicos (\emph{use after free})}
\label{sec:uaf}
Para el primer caso de evaluación,
analizaremos la capacidad del sistema para detectar un error temporal de memoria en lenguajes de bajo nivel,
el \emph{use after free}.
Este error ocurre cuando un programa intenta leer o escribir en una dirección de memoria después de que esta fue liberada y dejó de estar activa para el programa.

En la formalización,
una ubicación liberada no desaparece del \emph{heap}.
Como se mencionó en el \autoref{cap:formalizacion-isl},
la operación \lstinline|Free| transforma una celda \lstinline|Full v| en una celda \lstinline|Empty|.
Este recurso negativo permite conservar la información de que la dirección fue liberada y usarla después para justificar una ejecución errónea.

El ejemplo consiste en reservar una celda,
liberarla y luego intentar escribir sobre ella.
En la notación matemática de la ISL,
estos pasos se describen mediante las siguientes tripletas:

\[
\begin{array}{c}
  \itriple{\mathsf{emp}}{x := \texttt{alloc()}}{ok : x \mapsto -}, \\[2pt]

  \itriple{x \mapsto -}{\texttt{free(}x\texttt{)}}{ok : x \nmapsto}, \\[2pt]
  
  \itriple{x \nmapsto}{[x] := 1}{er : x \nmapsto}.
\end{array}
\]
La alocación produce una celda activa y la liberación la transforma en una celda negativa.
Esta última permite justificar que la escritura posterior termina con error,
conservando la información de que la dirección fue liberada.

En nuestra sintaxis,
este programa se representa mediante:

\begin{lstlisting}
let prog_uaf =
  (Alloc "x") `Seq`
  ((Free (Var "x")) `Seq`
  (Store (Var "x") (Const 1)))
\end{lstlisting}

La presunción inicial es \lstinline|is_ok emp|.
Esta aserción describe un estado en modo \lstinline|Ok| cuyo fragmento local del \emph{heap} no contiene ubicaciones conocidas.
La propiedad que se busca verificar tiene la forma
$$\vdash_\mathrm{ISL} \itriple{\operatorname{is\_ok} \; emp}{\operatorname{prog\_uaf}}{\operatorname{post\_uaf\_err}}.$$
En este caso,
la postcondición final describe un estado con modo \lstinline|Er| en el que la dirección almacenada en \lstinline|x| continúa representada con un \emph{heap} negativo.

El primer paso se justifica mediante la regla \lstinline|ISL_Alloc1|,
que permite reservar una dirección nueva respecto del fragmento local.
La postcondición elegida para este paso es:

\begin{lstlisting}
let post_alloc (s : state) : prop =
  exists (l : loc) (v : value).
    s._3 == Ok /\ s._1 "x" == Loc l /\
    l =!= 0 /\ points_to l v s /\
    (exists x_old. emp (override s._1 "x" x_old, override s._2 l Unknown, Ok))
\end{lstlisting}
Esta postcondición garantiza que la variable \lstinline|x| contiene una dirección no nula y que el fragmento local tiene una celda activa en esa ubicación.
El cuantificador sobre \lstinline|x_old| conserva la información necesaria para reconstruir el \emph{store} anterior a la asignación.

El segundo paso del programa es la liberación de la memoria apuntada por \lstinline|x|.
Al aplicar \lstinline|ISL_Free|,
se obtiene una postcondición que describe la celda liberada y conserva la información necesaria para reconstruir el estado anterior:

\begin{lstlisting}
let post_free (s : state) : prop =
  exists (v : value).
    s._3 == Ok /\ points_to_empty (eval_expr s (Var "x")) s /\
    post_alloc (s._1, override s._2 (eval_expr s (Var "x")) (Full v), Ok)
\end{lstlisting}
En este punto,
la variable \lstinline|x| sigue apuntando a la misma dirección física.
La diferencia está en el \emph{heap}, ya que la celda ya no se describe como \lstinline|Full v|, sino como \lstinline|Empty|.

El fallo ocurre en la tercera instrucción,
cuando el programa intenta realizar una escritura en la dirección indicada por el puntero invalidado.
La regla \lstinline|ISL_StoreEr| permite justificar este fallo porque \lstinline|post_free| garantiza
que la dirección indicada por \lstinline|x| está descrita mediante \lstinline|points_to_empty|.
%\guido{no podemos mostrar primero la derivación a alto nivel,
%en sintaxis matematica? algo como
%$emp \rightarrow x -> algo \rightarrow x -/-> \rightarrow Error$}
Por lo tanto,
la postcondición final se define como:

\begin{lstlisting}
let post_uaf_err (s : state) : prop =
  s._3 == Er /\ points_to_empty (eval_expr s (Var "x")) s /\
  post_free (s._1, s._2, Ok)
\end{lstlisting}
La ejecución cambia al modo \lstinline|Er|,
pero conserva el \emph{heap} negativo que justifica el error.
De esta manera,
la derivación no solo afirma que hubo un fallo,
sino que contiene la información espacial necesaria para mostrar que se debe a una escritura sobre una ubicación liberada.

La construcción de esta demostración en \fstar{} se realiza componiendo los axiomas de nuestra lógica
mediante constructores secuenciales y la regla de consecuencia.
El árbol de derivación completo para este programa se codifica de la siguiente manera:

\begin{lstlisting}
let proof_uaf : isl_triple (is_ok emp) prog_uaf post_uaf_err =
  let p_alloc_raw = ISL_Alloc1 #emp "x" in
  let p_alloc = ISL_Consequence emp post_alloc p_alloc_raw () () in
  let p_free_raw = ISL_Free #post_alloc (Var "x") in
  let p_free = ISL_Consequence post_alloc post_free p_free_raw () () in
  let p_store = ISL_StoreEr #post_free (Var "x") (Const 1) in
  let p_seq1 = ISL_Seq p_free p_store in
  let p_seq2 = ISL_Seq p_alloc p_seq1 in
  ISL_Consequence emp post_uaf_err p_seq2 () ()
\end{lstlisting}
Las reglas \lstinline|ISL_Alloc1|,
\lstinline|ISL_Free| e \lstinline|ISL_StoreEr| devuelven las derivaciones correspondientes a las operaciones primitivas.
\lstinline|ISL_Seq| las compone de acuerdo con la estructura del programa.
Las aplicaciones de \lstinline|ISL_Consequence| adaptan las postcondiciones generadas directamente por las reglas
a las aserciones intermedias elegidas para este caso.

La derivación \lstinline|proof_uaf| establece la tripleta anterior.
Para concluir además que existe una ejecución errónea,
se proporciona un estado final concreto:
la variable \lstinline|"x"| contiene \lstinline|Loc 1|,
la ubicación \lstinline|1| está marcada como \lstinline|Empty| y el modo es \lstinline|Er|.
El lema \lstinline|uaf_reachable| utiliza este testigo y \lstinline|certify| para demostrar \lstinline|error_reachable (is_ok emp) prog_uaf|:

\begin{lstlisting}
let st_uaf = on var (fun y -> if y = "x" then Loc 1 else Nat 0)

let uaf_reachable () : Lemma (error_reachable (is_ok emp) prog_uaf) =
  let hp_full = override (singleton_empty_heap 1) 1 (Full (Nat 0)) in
  assert (points_to 1 (Nat 0) (st_uaf, hp_full, Ok));
  certify proof_uaf (st_uaf, singleton_empty_heap 1, Er)
\end{lstlisting}
El \emph{heap} auxiliar \lstinline|hp_full| reconstruye la celda antes de su liberación con el contenido \lstinline|Nat 0|,
que sirve como testigo para justificar la postcondición de la derivación.

\section{Reasignación de memoria y punteros colgantes}
\label{sec:pushback}
Un desafío en el análisis estático de software de bajo nivel,
como C o C++,
es el manejo de estructuras de datos dinámicas,
tales como los vectores elásticos (\lstinline|std::vector|).
Cuando un programa agrega un elemento a una colección mediante una función como \lstinline|push_back|,
la biblioteca subyacente decide en tiempo de ejecución si posee capacidad suficiente o si debe reasignar el bloque de memoria para acomodar los nuevos elementos.
Esta reasignación implica reservar un nuevo bloque y liberar la memoria original.

El problema surge si un cliente del código conserva un puntero hacia un elemento de la estructura original y luego invoca la función de agregado,
ya que corre el riesgo de que la memoria sea reasignada.
En ese escenario, su puntero local se convierte en un puntero colgante.
Este fue un caso de estudio motivacional para el diseño de la ISL original~\citep{isl}.

El siguiente fragmento,
escrito en un lenguaje similar a C,
presenta una versión simplificada de este comportamiento:

\begin{lstlisting}
int* x = *v;
push_back(v); // Puede no hacer nada, o hacer free(*v) y reasignar
*x = 88;      // Posible uso de puntero colgante
\end{lstlisting}
%\guido{aclarar que v tiene tipo $int**$?}
Aquí,
\lstinline|v| tiene tipo \lstinline|int**| y apunta a una celda que almacena el puntero a la memoria de la estructura.
El cliente guarda una copia de ese puntero en \lstinline|x|,
que no se actualiza si \lstinline|push_back| modifica el contenido de \lstinline|*v|.
El ejemplo se centra en este cambio de referencia,
sin modelar los elementos,
la longitud ni la capacidad de un vector.

En nuestra formalización,
la memoria referenciada se representa mediante una única celda.
El parámetro \lstinline|ptr| identifica la variable del programa que cumple el papel de \lstinline|v|.
El comportamiento del cliente se codifica mediante:

\begin{lstlisting}
let client (ptr : var) : stmt =
  (Load "x" (Var ptr)) `Seq`
  ((push_back ptr) `Seq` (Store (Var "x") (Const 88)))
\end{lstlisting}
%\guido{aca no estamos modelando realmente el vector, no?
%deberiamos decir algo de que es un ejemplo simplificado}

La secuencia de operaciones del cliente se divide en tres etapas:
\begin{enumerate}
  \item \textbf{Captura de la referencia:}
  el cliente accede a la estructura original para guardar un puntero directo a su memoria activa mediante la instrucción de lectura \lstinline|Load|.
  \item \textbf{Interacción con la biblioteca:}
  el programa cliente invoca la función de agregado llamando a \lstinline|push_back|.
  Esta función de biblioteca encapsula un comportamiento no determinista interno,
  por lo que el cliente desconoce si la estructura mantendrá su bloque de memoria original o si será movida y reasignada en un bloque nuevo:
\begin{lstlisting}
let push_back (ptr : var) : stmt =
  Choice
    ((Load "y" (Var ptr)) `Seq`
    ((Free (Var "y")) `Seq`
    ((Alloc "y") `Seq`
    (Store (Var ptr) (Var "y")))))
    Skip
\end{lstlisting}
  \item \textbf{Escritura mediante el alias:}
  asumiendo erróneamente que su puntero local sigue siendo válido,
  el cliente intenta escribir en la dirección que contiene mediante \lstinline|Store|.
\end{enumerate}

La presunción \lstinline|pre_client ptr| describe dos celdas separadas:
una contiene la referencia de la estructura y la otra,
su contenido inicial.
Para construir la derivación se fija un contexto lógico \lstinline|c : ctx|,
que reúne el nombre \lstinline|c.ptr|,
las direcciones de la estructura (\lstinline|c.l_v|),
de la celda original (\lstinline|c.l_a|) y de la nueva celda (\lstinline|c.l_b|),
y sus contenidos.
Las tres direcciones son no nulas y distintas,
y \lstinline|c.ptr| es diferente de \lstinline|"x"| y de \lstinline|"y"|.
Estos parámetros describen el caso considerado y no son variables adicionales del programa:

\begin{lstlisting}
type ctx = {
  ptr   : p : var { p <> "x" /\ p <> "y" };
  l_v   : l : loc { l <> 0 };
  l_a   : l : loc { l <> 0 /\ l <> l_v };
  l_b   : l : loc { l <> 0 /\ l <> l_v /\ l <> l_a };
  val_a : value;
  val_b : value;
}
\end{lstlisting}

A partir de este contexto,
\lstinline|heap3 c c1 c2 c3| construye un \emph{heap} con esos contenidos en las tres ubicaciones y \lstinline|Unknown| en las restantes.
La aserción \lstinline|cells c c1 c2 c3| compara el \emph{heap} del estado con esa representación:

\begin{lstlisting}
let heap3 (c : ctx) (c1 c2 c3 : cell) : heap =
  on loc (fun l ->
    if l = c.l_v then c1 else if l = c.l_a then c2
    else if l = c.l_b then c3 else Unknown)

let cells (c : ctx) (c1 c2 c3 : cell) : cond =
  fun s -> forall l. s._2 l == heap3 c c1 c2 c3 l
\end{lstlisting}

La derivación parte de \lstinline|pre_fixed c|,
que fija la referencia de la estructura y de las dos celdas iniciales,
dejando \lstinline|c.l_b| como \lstinline|Unknown|.
La primera lectura almacena en \lstinline|"x"| la dirección original.
La aserción \lstinline|post_load_x c| conserva ese \emph{heap} y registra el valor leído.
La lectura posterior de \lstinline|"y"| agrega la misma dirección en \lstinline|post_load_y c|.

\begin{lstlisting}
let pre_fixed (c : ctx) : cond = fun s ->
  s._1 c.ptr == Loc c.l_v /\
  cells c (Full (Loc c.l_a)) (Full c.val_a) Unknown s

let post_load_x (c : ctx) : cond = fun s ->
  is_ok (pre_fixed c) s /\ s._1 "x" == Loc c.l_a

let post_load_y (c : ctx) : cond = fun s ->
  post_load_x c s /\ s._1 "y" == Loc c.l_a
\end{lstlisting}

Para demostrar el error no es necesario analizar las dos ramas de \lstinline|push_back|.
La interpretación subaproximada permite seleccionar la rama que conduce al estado de interés.
En este caso,
la derivación selecciona la rama izquierda mediante \lstinline|ISL_ChoiceL|.
Esto no afirma que la rama derecha sea imposible,
sino que la rama izquierda tiene una ejecución válida que alcanza la postcondición buscada.

Para justificar el error,
se consideran las ejecuciones de esta rama en las que la alocación elige una dirección distinta de la liberada.
Así,
la estructura pasa a referenciar la nueva celda,
mientras que la dirección conservada en \lstinline|"x"| permanece marcada como liberada.

Esta información se registra en las aserciones intermedias,
que describen el \emph{heap} en el caso considerado después de cada operación.
La regla local se aplica a la celda utilizada y después \lstinline|ISL_Frame| incorpora los recursos restantes.
En particular,
\lstinline|post_load_y c| describe el estado después de cargar la dirección original en \lstinline|"y"|.
Tras liberarla se obtiene:

\begin{lstlisting}
let post_free_y (c : ctx) : cond = fun s ->
  s._3 == Ok /\ s._1 c.ptr == Loc c.l_v /\
  s._1 "x" == Loc c.l_a /\ s._1 "y" == Loc c.l_a /\
  cells c (Full (Loc c.l_a)) Empty Unknown s
\end{lstlisting}
La ubicación \lstinline|c.l_a| queda marcada como \lstinline|Empty|,
mientras que la celda de la estructura conserva la referencia original.

Para justificar esta transición mediante la regla local,
se separa la celda que se libera del resto del \emph{heap}.
Las aserciones \lstinline|free_pre c| y \lstinline|free_post c| describen esa celda antes y después de la operación.
Por otro lado,
\lstinline|frame_free c| conserva la celda de la estructura y los valores de \lstinline|c.ptr| y de \lstinline|"x"|:

\begin{lstlisting}
let free_pre (c : ctx) : cond = fun s ->
  s._3 == Ok /\ s._1 "y" == Loc c.l_a /\ points_to c.l_a c.val_a s

let free_post (c : ctx) : cond = fun s ->
  s._1 "y" == Loc c.l_a /\ points_to_empty c.l_a s

let frame_free (c : ctx) : cond = fun s ->
  s._1 c.ptr == Loc c.l_v /\ s._1 "x" == Loc c.l_a /\
  cells c (Full (Loc c.l_a)) Unknown Unknown s
\end{lstlisting}

La regla \lstinline|ISL_Free| exige que,
al restaurar algún contenido en la celda liberada,
se obtenga un estado que satisfaga la presunción.
Para adaptar esa postcondición a \lstinline|free_post c|,
el lema \lstinline|lemma_free_back| elige \lstinline|c.val_a| como contenido anterior:

\begin{lstlisting}
let lemma_free_back (c : ctx) (s : state)
  : Lemma (requires is_ok (free_post c) s)
          (ensures is_ok (fun t -> exists w.
            points_to_empty (eval_expr t (Var "y")) t /\
            free_pre c (t._1, override t._2 (eval_expr t (Var "y"))
              (Full w), Ok)) s) =
  exists_intro (fun w -> points_to_empty (eval_expr s (Var "y")) s /\
    free_pre c (s._1, override s._2 (eval_expr s (Var "y")) (Full w), Ok)) c.val_a
\end{lstlisting}
Como \lstinline|free_post c| fija el valor de \lstinline|"y"| en \lstinline|Loc c.l_a|,
la actualización reconstruye la celda \lstinline|c.l_a| con contenido \lstinline|Full c.val_a|.
El estado reconstruido satisface \lstinline|free_pre c|,
y \lstinline|exists_intro| aporta ese contenido como testigo.
Esta reconstrucción pertenece a la prueba y no agrega una escritura al programa.

Una vez justificada la operación local,
hay que relacionar su resultado con el estado descrito por \lstinline|post_free_y c|.
Para ello,
\lstinline|split_free| exhibe dos \emph{heaps}:
uno que contiene solo la celda liberada y otro que conserva la celda de la estructura:

\begin{lstlisting}
let split_free (c : ctx) (s : state)
  : Lemma (requires post_free_y c s)
          (ensures (is_ok (free_post c) ** frame_free c) s) =
  introduce_frame (is_ok (free_post c)) (frame_free c) s
    (singleton_empty_heap c.l_a)
    (heap3 c (Full (Loc c.l_a)) Unknown Unknown)
\end{lstlisting}
\lstinline|singleton_empty_heap c.l_a| asigna \lstinline|Empty| a \lstinline|c.l_a| y \lstinline|Unknown| a las demás ubicaciones.
El segundo \emph{heap} contiene la celda \lstinline|c.l_v| y deja \lstinline|c.l_a| como \lstinline|Unknown|.
Como ambas direcciones son distintas,
los fragmentos son disjuntos y su unión coincide con el \emph{heap} de \lstinline|post_free_y c|.
El auxiliar \lstinline|introduce_frame| reúne esos testigos y las aserciones que satisfacen para concluir la conjunción separadora.

Con estas dos implicaciones,
se construye la derivación de la liberación:

\begin{lstlisting}
let step_free (c : ctx)
  : isl_triple (is_ok (post_load_y c)) (Free (Var "y")) (post_free_y c) =
  let local = ISL_Consequence (free_pre c) (is_ok (free_post c))
    (ISL_Free #(free_pre c) (Var "y")) ()
    (forall_intro (move_requires (lemma_free_back c))) in
  ISL_Consequence (post_load_y c) (post_free_y c)
    (ISL_Frame (frame_free c) local ()) ()
    (forall_intro (move_requires (split_free c)))
\end{lstlisting}
La primera aplicación de \lstinline|ISL_Consequence| utiliza \lstinline|lemma_free_back| para obtener la derivación local
con resultado \lstinline|is_ok (free_post c)|.
Luego,
\lstinline|ISL_Frame| incorpora el marco,
que se conserva porque \lstinline|Free| no modifica variables del \emph{store}.
La última aplicación de \lstinline|ISL_Consequence| amplía la presunción a \lstinline|post_load_y c| y restringe el resultado a \lstinline|post_free_y c|,
cuya descomposición justifica \lstinline|split_free|.
Las combinaciones de \lstinline|move_requires| y \lstinline|forall_intro| presentan los lemas como implicaciones para todo estado.
Los demás pasos utilizan las reglas y los marcos correspondientes a cada instrucción.

Después de esta liberación,
la rama reserva otra celda y escribe su dirección en la celda de la estructura,
alcanzando \lstinline|post_store_v c|.
El cliente intenta entonces escribir mediante \lstinline|"x"|,
que conserva la dirección original liberada.
La derivación \lstinline|step_error| aplica \lstinline|ISL_StoreEr| a esa celda y utiliza el marco para conservar los otros recursos.
La postcondición final se expresa como:

\begin{lstlisting}
let post_err (ptr : var) (l_v l_a l_b : loc) (val_b : value) : cond =
  fun s -> s._3 == Er /\
    s._1 ptr == Loc l_v /\ s._1 "x" == Loc l_a /\ s._1 "y" == Loc l_b /\
    (points_to_empty l_a **
      (points_to l_v (Loc l_b) ** points_to l_b val_b)) s
\end{lstlisting}
La postcondición conserva tres recursos separados:
la celda original liberada,
la celda de la estructura actualizada y la nueva celda reservada.
La variable \lstinline|"x"| mantiene la dirección original,
por lo que la escritura sobre esa ubicación produce \lstinline|Er|.

Cada operación se justifica mediante una función \lstinline|step_*|.
Las cuatro operaciones de la rama de reasignación se componen mediante \lstinline|ISL_Seq| en \lstinline|proof_push_back|,
y \lstinline|ISL_ChoiceL| selecciona esa rama:

\begin{lstlisting}
let proof_push_back (c : ctx)
  : isl_triple (is_ok (post_load_x c)) (push_back c.ptr) (post_store_v c) =
  let branch = ISL_Seq (step_load_y c) (ISL_Seq (step_free c)
    (ISL_Seq (step_alloc c) (step_store c))) in
  ISL_Consequence (post_load_x c) (post_store_v c) (ISL_ChoiceL branch) () ()
\end{lstlisting}

Finalmente,
\lstinline|proof_client_ctx| compone esta derivación con \lstinline|step_load_x| y \lstinline|step_error|.
Antes del ajuste final,
la construcción parte de \lstinline|pre_fixed c| y concluye \lstinline|error_fixed c|,
que describe los mismos valores del \emph{store} y las mismas tres celdas que \lstinline|post_err c.ptr c.l_v c.l_a c.l_b c.val_b| con modo \lstinline|Er|.
Mediante \lstinline|initial_implies| y la regla \lstinline|ISL_Consequence|,
se amplía la presunción a \lstinline|pre_client c.ptr| y se expresa el resultado con \lstinline|post_err c.ptr c.l_v c.l_a c.l_b c.val_b|.
La función \lstinline|proof_client| recibe los parámetros por separado y construye el contexto.

Para obtener un certificado de error,
se instancia la derivación con la variable \lstinline|"v"| y las direcciones \lstinline|1|,
\lstinline|2| y \lstinline|3| para la estructura,
la celda original y la nueva,
respectivamente.
El estado final conserva \lstinline|"x"| apuntando a la dirección liberada,
mientras que la estructura e \lstinline|"y"| referencian la nueva celda:

\begin{lstlisting}
let st_pb = on var (fun z ->
  if z = "v" then Loc 1 else if z = "x" then Loc 2
  else if z = "y" then Loc 3 else Nat 0)

let h_new = on loc (fun l ->
  if l = 1 then Full (Loc 3) else if l = 3 then Full (Nat 0)
  else Unknown)

let h_err = on loc (fun l ->
  if l = 2 then Empty else h_new l)

let pb_reachable ()
  : Lemma (error_reachable (is_ok (pre_client "v")) (client "v")) =
  let s1 = (st_pb, h_err, Er) in
  introduce_frame (points_to 1 (Loc 3)) (points_to 3 (Nat 0)) (st_pb, h_new, Er)
    (singleton_full_heap 1 (Loc 3)) (singleton_full_heap 3 (Nat 0));
  introduce_frame (points_to_empty 2) (points_to 1 (Loc 3) ** points_to 3 (Nat 0))
    s1 (singleton_empty_heap 2) h_new;
  certify (proof_client "v" 1 2 3 (Nat 0) (Nat 0)) s1
\end{lstlisting}
Las llamadas a \lstinline|introduce_frame| justifican las dos conjunciones separadoras de la postcondición:
las ubicaciones \lstinline|1| y \lstinline|3| contienen la estructura actualizada y la nueva celda,
y la ubicación \lstinline|2| conserva la marca \lstinline|Empty| de la celda liberada.
Con este estado final,
\lstinline|certify| demuestra que el error es alcanzable y que la ejecución también puede trasladarse a los estados completados.

Este caso muestra que la formalización puede relacionar un error temporal de memoria con una modificación intermedia realizada a través de otro alias.
La celda negativa conserva la historia del bloque original,
mientras que el no determinismo permite seleccionar la ejecución de reasignación.

\section{No determinismo y razonamiento paramétrico sobre bucles}
\label{sec:bucle}
%A lo largo del desarrollo de software,
%la presencia de bucles combinados con comportamientos no deterministas supone una barrera para las técnicas de verificación tradicionales.
%Para comprender los límites prácticos de dichas metodologías frente a la subaproximación,
%evaluaremos un programa semirrealista con una alta cantidad de iteraciones y decisiones no deterministas.
%El caso busca mostrar dos propiedades del sistema deductivo:
%\begin{enumerate}
%  \item una derivación subaproximada puede seleccionar un único camino entre varias alternativas (en la \autoref{sec:pushback} se vio en menor magnitud);
%  \item la regla de variante de Kleene permite representar paramétricamente una familia de iteraciones sin construir una derivación distinta para cada una.
%\end{enumerate}
El ejemplo anterior utilizó una elección no determinista para seleccionar la rama de reasignación.
En este caso,
esa elección se repite dentro de un bucle con una gran cantidad de iteraciones,
lo que permite estudiar cómo construir una derivación paramétrica del camino que conduce al error.

Consideremos el siguiente programa:

\begin{lstlisting}
n := 1000000;
i, j := 0;
while (i < n){
  i++;
  if (random()) j++;
}
assert  (j != n);
\end{lstlisting}
El bucle se ejecuta un millón de veces.
En cada iteración,
\lstinline|i| se incrementa siempre,
mientras que el incremento de \lstinline|j| depende de una decisión no determinista.
Si la rama de incremento se selecciona en todas las iteraciones,
al finalizar se cumple que tanto \lstinline|i| como \lstinline|j| son iguales a un millón.
En ese estado,
falla la comprobación \lstinline|assert (j != n)|.

En la semántica formalizada,
la llamada \lstinline|random()| se representa con un \lstinline|Choice|.
No se modelan probabilidades,
independencia estadística ni distribuciones.
Si lo interpretamos bajo la visión de un programa ejecutable en la que cada decisión es independiente y equiprobable,
la probabilidad de incrementar \lstinline|j| en el millón de iteraciones sería de $2^{-1000000}$.
Este cálculo no interviene en la demostración mecanizada.
La lógica solo necesita establecer que la secuencia de elecciones existe.

Esta elección no determinista se expresa en nuestra sintaxis mediante el operador \lstinline{|||},
el cual abrevia \lstinline|Choice|:

\begin{lstlisting}
Assign "j" (Plus (Var "j") (Const 1)) ||| Skip
\end{lstlisting}

Para construir la ejecución que lleva al error,
la derivación usa \lstinline|ISL_ChoiceL| y selecciona la rama de incremento en el cuerpo del bucle.
Una exploración que mantuviera por separado cada alternativa de ejecución duplicaría el número de caminos en cada iteración.

A pesar de poder elegir una sola rama,
construir un millón de aplicaciones de la regla de Kleene produciría una derivación enorme.
Para evitar ese desenrollado se usa \lstinline|ISL_KleeneVariant|.

La familia de aserciones se define como sigue:

\begin{lstlisting}
let variant (k : nat) : cond =
  fun (st, _, m) ->
    k <= 1000000 /\ st "n" == Nat 1000000 /\
    st "i" == Nat k /\ st "j" == Nat k /\ m == Ok
\end{lstlisting}
Este predicado asocia el número de iteración lógica \lstinline|k| con los valores concretos del \emph{store}.
Afirma que,
después de \lstinline|k| iteraciones del camino seleccionado,
tanto \lstinline|i| como \lstinline|j| son iguales a \lstinline|k|.

Para aplicar \lstinline|ISL_KleeneVariant| debe construirse,
para cada natural \lstinline|k|,
una derivación del cuerpo desde \lstinline|is_ok (variant k)| hasta \lstinline|is_ok (variant (k + 1))|.
La demostración se realiza paramétricamente en \lstinline|k|.
La rama izquierda de \lstinline|Choice| incrementa \lstinline|j|,
mientras que las asignaciones restantes incrementan \lstinline|i|.
Las obligaciones aritméticas locales son verificadas por el \emph{solver}.
Una vez demostrado el paso general,
la regla permite obtener estados correspondientes a una cantidad finita de iteraciones
sin construir manualmente una derivación diferente para cada valor del índice.

La salida del bucle se representa mediante \lstinline|Assume (Eq (Var "i") (Var "n"))|.
En este programa,
\lstinline|i| comienza en cero y aumenta de a uno mientras sea menor que \lstinline|n|,
por lo que nunca supera ese valor.
Así,
sobre los estados alcanzables del bucle,
la igualdad equivale a la negación de su condición de continuación.

La instrucción \lstinline|assert| se codifica como una elección entre el camino de error y el exitoso.
Como la sintaxis de expresiones no incluye un constructor de negación,
para este ejemplo se define el auxiliar \lstinline|enot| como \lstinline|Minus (Const 1) e|.
Como aquí se aplica a expresiones booleanas que evalúan a 0 o 1,
\lstinline|enot e| evalúa a 1 cuando \lstinline|e| evalúa a 0,
y a 0 cuando evalúa a 1.
De este modo se expresa la condición opuesta sin ampliar la sintaxis de \lstinline|expr|.

\begin{lstlisting}
let assert_stmt =
  Seq (Assume (Eq (Var "j") (Var "n"))) Error |||
  Assume (enot (Eq (Var "j") (Var "n")))
\end{lstlisting}
La rama izquierda representa el caso de \lstinline|j| y \lstinline|n| iguales,
donde el \lstinline|Assume| permite continuar y luego \lstinline|Error| cambia el modo de terminación.
La rama derecha representa el caso en que \lstinline|j| y \lstinline|n| son distintos.

La variante construida para el camino seleccionado permite alcanzar el estado en el que \lstinline|i|,
\lstinline|j| y \lstinline|n| son iguales a un millón.
Por lo tanto,
la igualdad evaluada por el primer \lstinline|Assume| es verdadera.
La derivación vuelve a usar \lstinline|ISL_ChoiceL| para seleccionar la rama que lleva a \lstinline|Error|.

La postcondición final es:

\begin{lstlisting}
let post_err : cond =
  fun (st, _, m) ->
    st "n" == Nat 1000000 /\ st "i" == Nat 1000000 /\
    st "j" == Nat 1000000 /\ m == Er
\end{lstlisting}

Una búsqueda exhaustiva que enumere las secuencias de decisiones tendría que considerar hasta $2^{1000000}$ combinaciones,
lo que hace inviable explorar por completo este espacio para localizar el error.
La derivación presentada selecciona el camino que incrementa \lstinline|j| en todas las iteraciones
y utiliza una familia de aserciones para justificarlo paramétricamente,
sin enumerar las demás combinaciones.
%\guido{las dos ultimas oraciones deberian ser una sola, y creo que en este
%parrafo queres justificar que una busqueda exhaustiva tendria dificultad
%para encontrar este error, decilo explicitamente}

La satisfacibilidad de la postcondición se establece con un \emph{store} que asigna \lstinline|Nat 1000000| a todas las variables,
un \emph{heap} vacío y modo \lstinline|Er|.
Solo los valores de \lstinline|"n"|, \lstinline|"i"| y \lstinline|"j"| son relevantes para \lstinline|post_err|.
La derivación completa se denomina \lstinline|proof_random_loop| y se aplica al programa \lstinline|prog_random_loop|.
El certificado es:

\begin{lstlisting}
let rl_reachable () : Lemma (error_reachable (is_ok pre_init) prog_random_loop) =
  let st = on var (fun _ -> Nat 1000000) in
  certify proof_random_loop (st, empty_heap, Er)
\end{lstlisting}

\section{Recorrido de una lista enlazada}
\label{sec:listas}
Este caso combina el razonamiento iterativo de la sección anterior con estructuras de datos almacenadas en memoria dinámica.
Para este escenario,
formalizamos un programa que recorre una lista enlazada representada mediante nodos de dos celdas,
acumulando la suma de sus elementos en la variable \lstinline|sum| y contando su longitud en la variable \lstinline|len|.
Al finalizar,
el programa verifica mediante una instrucción \lstinline|assert| que la suma total sea distinta a la longitud calculada.
Si la suma es igual,
la comprobación falla y el programa aborta.

Conceptualmente,
el código bajo análisis se describe de la siguiente manera:

\begin{lstlisting}
sum = 0;
len = 0;
while (ptr != 0) {
  v = *ptr;
  sum += v;
  len++;
  ptr = *(ptr + 1);
}
assert(sum != len);
\end{lstlisting}

Para este caso,
modelamos una lista enlazada de longitud arbitraria \lstinline|n|.
La demostración es paramétrica en esta longitud y no requiere escribir una prueba distinta para cada cantidad de nodos.
Además,
la presunción describe una lista cuyos elementos son todos iguales a 1.
Esta condición es suficiente para garantizar que,
al finalizar el recorrido,
la suma de los elementos coincide con la longitud de la lista y la comprobación \lstinline|assert| falla.
No es la única configuración posible que satisface dicha igualdad,
pero permite expresar el argumento mediante un predicado espacial paramétrico simple.

En lenguajes de bajo nivel como C,
un nodo de una lista enlazada se representa típicamente como una estructura que agrupa datos y referencias.
En nuestra formalización,
modelamos esta estructura desarmándola en sus componentes de memoria fundamentales.
Un nodo ocupa dos ubicaciones de memoria adyacentes:
\begin{enumerate}
  \item la dirección base \lstinline|curr_ptr|,
  que almacena el valor numérico del nodo;
  \item y la dirección consecutiva \lstinline|curr_ptr + 1|,
  que almacena el puntero hacia el siguiente nodo de la estructura.
\end{enumerate}

Para describir una lista con una longitud arbitraria se define el predicado \lstinline|list_seg|.
La aserción \lstinline|list_seg i j xs| describe un segmento de lista que empieza en \lstinline|i|,
termina en \lstinline|j| y tiene la secuencia de valores \lstinline|xs|.
El caso vacío representa un segmento sin nodos.
El caso inductivo describe las dos celdas del primer nodo y el segmento restante con la conjunción separadora:

\begin{lstlisting}
let rec list_seg (i j : loc) (xs : list value) : Tot cond (decreases xs) =
  match xs with
  | [] -> fun s -> emp s /\ i == j
  | x :: xss -> fun s -> 
    exists (next : loc).
      (points_to i x ** points_to (i + 1) (Loc next) ** list_seg next j xss) s
\end{lstlisting}
Este predicado permite expresar la estructura completa sin enumerar cada dirección de forma individual.
La conjunción separadora garantiza que las celdas de nodos diferentes pertenezcan a fragmentos disjuntos.

A partir de esta representación,
se define el predicado \lstinline|history|
para relacionar la posición actual del puntero con el inicio de la lista.
Este predicado describe un recorrido de \lstinline|k| pasos desde \lstinline|start_ptr| hasta \lstinline|curr_ptr| sobre un mismo \emph{heap}.
La lista de valores utilizada en su caso base se construye mediante \lstinline|ones n|,
que contiene \lstinline|n| elementos iguales a \lstinline|Nat 1|.
La aserción se define de la siguiente manera:

\begin{lstlisting}
let heap_only (hp : heap) : state = on var (fun _ -> Nat 0), hp, Ok

let rec history (start_ptr curr_ptr : loc) (n : nat) (k : nat) (hp : heap)
  : GTot prop (decreases k) =
  let s = heap_only hp in
  if k = 0 then
    curr_ptr == start_ptr /\
    list_seg start_ptr 0 (ones n) s
  else exists (prev_ptr : loc).
    points_to_flex prev_ptr (Nat 1) s /\
    points_to_flex (prev_ptr + 1) (Loc curr_ptr) s /\
    history start_ptr prev_ptr n (k - 1) hp
\end{lstlisting}

La función auxiliar \lstinline|heap_only| permite evaluar las aserciones espaciales sobre el
\emph{heap} usando el \emph{store} fijo,
cuyos valores no intervienen en ellas.
El caso base de \lstinline|history| exige una lista de \lstinline|n| unos y ubica el puntero en su inicio.
El caso recursivo identifica un nodo anterior que contiene el valor uno y apunta a la posición actual.
Las consultas mediante \lstinline|points_to_flex| se hacen sobre el mismo \emph{heap} de la lista,
por lo que se combinan con la conjunción clásica ($\land$).
La separación entre celdas de la lista queda establecida por \lstinline|list_seg| en el caso base.

\begin{figure}[htbp]
\centering
\begin{tikzpicture}[
  cell/.style={
    rectangle,
    draw=black,
    thick,
    minimum width=0.8cm,
    minimum height=0.8cm,
    inner sep=0pt,
    outer sep=0pt
  },
  ptr/.style={->, thick, >=stealth}
]

% Primer nodo
\node[cell] (n1v) at (0,0) {$1$};
\node[cell, below=-\pgflinewidth of n1v]
  (n1p) {$\bullet$};

% Segundo nodo
\node[cell, right=0.8cm of n1v] (n2v) {$1$};
\node[cell, below=-\pgflinewidth of n2v]
  (n2p) {$\bullet$};

% Nodo alcanzado después de dos avances
\node[cell, fill=blue!15, right=0.8cm of n2v]
  (currv) {$1$};
\node[cell, fill=blue!15, below=-\pgflinewidth of currv]
  (currp) {$\bullet$};

% Cuarto nodo
\node[cell, right=0.8cm of currv] (n4v) {$1$};
\node[cell, below=-\pgflinewidth of n4v]
  (n4p) {$\bullet$};

% Último nodo: enlace nulo
\node[cell, right=0.8cm of n4v] (n5v) {$1$};
\node[cell, below=-\pgflinewidth of n5v]
  (n5p) {$0$};

% Enlaces entre nodos
\draw[ptr]
  (n1p.center) to[out=0,in=180] (n2v.west);
\draw[ptr]
  (n2p.center) to[out=0,in=180] (currv.west);
\draw[ptr]
  (currp.center) to[out=0,in=180] (n4v.west);
\draw[ptr]
  (n4p.center) to[out=0,in=180] (n5v.west);

% Direcciones inicial y actual
\node[above=0.7cm of n1v]
  (startlabel) {\small\texttt{start\_ptr}};
\draw[ptr] (startlabel.south) -- (n1v.north);

\node[above=0.7cm of currv]
  (currlabel) {\small\texttt{curr\_ptr}};
\draw[ptr] (currlabel.south) -- (currv.north);

% Cantidad de avances desde el inicio
\draw[
  decorate,
  decoration={brace,amplitude=6pt,mirror,raise=5pt}
]
  (n1p.south) -- (currp.south)
  node[midway,yshift=-25pt]
  {\small $k = 2$ avances};

\end{tikzpicture}
\caption{Ejemplo del recorrido descrito por
\lstinline|history|: desde \lstinline|start_ptr|
se alcanza \lstinline|curr_ptr| después de dos avances
en una lista de cinco unos.}
\label{fig:particion-lista}
\end{figure}

La \autoref{fig:particion-lista} ilustra un caso de \lstinline|history| con \lstinline|n = 5| y \lstinline|k = 2|:
el puntero actual se encuentra en el tercer nodo,
luego de realizar dos avances desde el inicio de la lista.

La presunción ubica \lstinline|ptr| al inicio de la lista y exige que el estado tenga modo \lstinline|Ok|.
Al utilizar \lstinline|history| con cero pasos,
también establece que el \emph{heap} describe una lista de \lstinline|n| elementos iguales a uno:

\begin{lstlisting}
let pre_init (start_ptr : loc) (n : nat) : cond =
  fun s -> s._3 == Ok /\
    s._1 "ptr" == Loc start_ptr /\
    history start_ptr start_ptr n 0 s._2
\end{lstlisting}

La familia de aserciones utilizada para el bucle relaciona la cantidad de iteraciones con los acumuladores
y con la posición alcanzada en la lista:

\begin{lstlisting}
let variant (start_ptr : loc) (n : nat) (k : nat) : cond =
  fun s ->
    k <= n /\ s._3 == Ok /\
    s._1 "sum" == Nat k /\ s._1 "len" == Nat k /\
    (k > 0 ==> s._1 "v" == Nat 1) /\
    exists (curr_ptr : loc).
      s._1 "ptr" == Loc curr_ptr /\
      history start_ptr curr_ptr n k s._2
\end{lstlisting}
La inicialización establece que \lstinline|sum| y \lstinline|len| valen cero y conserva el puntero y el \emph{heap} descritos por \lstinline|pre_init|.
Así se obtiene el caso \lstinline|k = 0| de la variante,
como demuestra \lstinline|proof_init|.
La condición \lstinline|s._1 "v" == Nat 1| se exige solo cuando \lstinline|k > 0|,
porque \lstinline|v| recibe su valor durante la primera lectura del cuerpo.

Para justificar el paso entre dos índices consecutivos de la variante,
se construye una derivación para cada instrucción del cuerpo.
Las dos lecturas requieren acceder a una celda sin perder la información sobre el resto de la lista.
Para ello se utiliza la derivación auxiliar \lstinline|framed_load|,
que combina \lstinline|ISL_Load| con la Regla del Marco.
Para construir el marco,
se utilizan variables auxiliares que conservan el valor anterior de la variable modificada,
la dirección leída y el contenido de la celda.
Así,
el marco puede describir el resto de la lista sin depender del valor actual de la variable que modifica la lectura.
Estas variables se ocultan mediante \lstinline|ISL_Exists| y no forman parte del programa.
Para que estas variables sean exclusivamente auxiliares,
\lstinline|load_fresh| exige que no aparezcan en la lectura y \lstinline|ghost_free pre| garantiza que la presunción no dependa de sus valores.
La postcondición resultante relaciona el valor leído con un estado anterior obtenido al restablecer la variable modificada:

\begin{lstlisting}
let load_post (pre : cond) (x : var) (e : expr) : cond =
  fun s -> s._3 == Ok /\
    exists (x_old : value) (l : loc) (v : value).
      points_to_flex l v s /\ s._1 x == v /\
      eval_expr (override s._1 x x_old, s._2, Ok) e == l /\
      pre (override s._1 x x_old, s._2, Ok)
\end{lstlisting}

La primera lectura establece \lstinline|post_load_v|,
donde ambos acumuladores valen \lstinline|Nat k| y \lstinline|v| contiene \lstinline|Nat 1|,
junto con las condiciones del recorrido.
A continuación,
el programa actualiza los acumuladores mediante dos asignaciones.
Durante estas actualizaciones,
el puntero permanece en el mismo nodo y el \emph{heap} no cambia.
Para describir los estados intermedios,
se define \lstinline|step_state|,
cuyos parámetros \lstinline|sum| y \lstinline|len| indican los valores de ambos acumuladores:

\begin{lstlisting}
let step_state (start_ptr : loc) (n : nat) (k : nat) (sum len : nat) : cond =
  fun s ->
    k < n /\ s._3 == Ok /\
    s._1 "sum" == Nat sum /\ s._1 "len" == Nat len /\
    s._1 "v" == Nat 1 /\
    exists (curr_ptr : loc).
      s._1 "ptr" == Loc curr_ptr /\
      history start_ptr curr_ptr n k s._2 /\
      points_to_flex curr_ptr (Nat 1) s
\end{lstlisting}

El índice \lstinline|k| indica la cantidad de avances del puntero realizados hasta el nodo actual,
mientras que los acumuladores se actualizan por separado.
La aserción \lstinline|post_sum| describe el estado posterior a incrementar la suma,
y \lstinline|post_len|,
el estado posterior a incrementar también la longitud:

\begin{lstlisting}
let post_sum (start_ptr : loc) (n : nat) (k : nat) : cond =
  step_state start_ptr n k (k + 1) k

let post_len (start_ptr : loc) (n : nat) (k : nat) : cond =
  step_state start_ptr n k (k + 1) (k + 1)
\end{lstlisting}

Para justificar la actualización de la suma mediante \lstinline|ISL_Assign|,
el siguiente lema reconstruye,
a partir de un estado que satisface \lstinline|post_sum|,
un \emph{store} anterior que satisface \lstinline|post_load_v|:

\begin{lstlisting}
let lemma_sum (start_ptr : loc) (n : nat) (k : nat) (s : state)
  : Lemma (requires (post_sum start_ptr n k s))
          (ensures (is_ok (fun s0 -> exists (st_init : store).
            post_load_v start_ptr n k (st_init, s0._2, s0._3) /\
            s0._1 "sum" == eval_value st_init (Plus (Var "sum") (Var "v")) /\
            (forall (y : var). y <> "sum" ==> s0._1 y == st_init y)) s)) =
  let st_init = override s._1 "sum" (Nat k) in
  exists_intro (fun (st_init : store) ->
    post_load_v start_ptr n k (st_init, s._2, s._3) /\
    s._1 "sum" == eval_value st_init (Plus (Var "sum") (Var "v")) /\
    (forall (y : var). y <> "sum" ==> s._1 y == st_init y)) st_init
\end{lstlisting}
El \emph{store} anterior se obtiene restableciendo únicamente \lstinline|sum| a \lstinline|Nat k|.
Como \lstinline|v| contiene \lstinline|Nat 1|,
la expresión de la asignación produce el valor \lstinline|Nat (k + 1)| del estado recibido.
Las demás variables y el \emph{heap} se conservan.
La llamada a \lstinline|exists_intro| proporciona este \emph{store} como testigo de la postcondición de \lstinline|ISL_Assign|.
Omitimos el lema correspondiente a \lstinline|len|,
ya que sigue el mismo patrón.

Una vez actualizados estos acumuladores,
el estado satisface \lstinline|post_len|:
\lstinline|sum| y \lstinline|len| valen \lstinline|Nat (k + 1)|,
pero el puntero todavía se encuentra en el nodo actual.
La última instrucción del cuerpo lee la dirección almacenada en \lstinline|ptr + 1| y la asigna a \lstinline|ptr|.
Este paso permite alcanzar un estado que satisface \lstinline|variant start_ptr n (k + 1)|.

En esta formalización,
si \lstinline|ptr| contiene \lstinline|Loc l|,
la expresión \lstinline|ptr + 1| se evalúa numéricamente como $l + 1$.
La regla \lstinline|Load| usa ese resultado como dirección del \emph{heap};
no requiere que la expresión produzca un valor de tipo \lstinline|Loc|.
Por eso la lectura en \lstinline|ptr + 1| queda bien definida en el modelo.
Si ese resultado se almacenara mediante \lstinline|eval_value|,
se representaría como \lstinline|Nat (l + 1)|.

Para justificar esta lectura,
el lema \lstinline|lemma_load_ptr| parte del estado posterior al avance.
Al desplegar \lstinline|history| para \lstinline|k + 1| pasos,
se obtiene un nodo anterior cuya celda de enlace contiene la dirección actual.
Esta información permite reconstruir el valor previo de \lstinline|ptr| y establecer la postcondición \lstinline|load_post| de la lectura.
Como la consulta de la celda mediante \lstinline|points_to_flex| no depende del \emph{store},
la información se conserva al restablecer el valor de \lstinline|"ptr"|.
El auxiliar \lstinline|eval_pointer| justifica la dirección evaluada antes de la lectura,
y \lstinline|load_post_intro| reúne los datos necesarios para establecer su postcondición:

\begin{lstlisting}
let lemma_load_ptr (start_ptr : loc) (n : nat) (k : nat) (s : state)
  : Lemma (requires variant start_ptr n (k + 1) s)
          (ensures load_post (post_len start_ptr n k) "ptr" (Plus (Var "ptr") (Const 1)) s) =
  let curr_ptr = indefinite_description_ghost
    loc (fun curr_ptr -> s._1 "ptr" == Loc curr_ptr /\ history start_ptr curr_ptr n (k + 1) s._2)
  in
  let hs = heap_only s._2 in
  let prev_ptr = indefinite_description_ghost loc (fun prev_ptr ->
    points_to_flex prev_ptr (Nat 1) hs /\
    points_to_flex (prev_ptr + 1) (Loc curr_ptr) hs /\
    history start_ptr prev_ptr n k s._2)
  in
  exists_intro (fun p ->
    (override s._1 "ptr" (Loc prev_ptr)) "ptr" == Loc p /\
    history start_ptr p n k s._2 /\
    points_to_flex p (Nat 1) (override s._1 "ptr" (Loc prev_ptr), s._2, Ok)) prev_ptr;
  eval_pointer s._1 s._2 prev_ptr;
  load_post_intro (post_len start_ptr n k) "ptr" (Plus (Var "ptr") (Const 1)) s (Loc prev_ptr) (prev_ptr + 1) (Loc curr_ptr)
\end{lstlisting}
En este caso,
el testigo \lstinline|prev_ptr| proporciona el valor anterior del puntero.
Al restablecerlo en el \emph{store},
se obtiene un estado que satisface \lstinline|post_len|,
con los acumuladores ya actualizados y el recorrido todavía en el índice \lstinline|k|.

El lema anterior se utiliza en \lstinline|step_load_ptr|,
que obtiene la derivación de la lectura mediante \lstinline|framed_load| y adapta su postcondición mediante \lstinline|ISL_Consequence|.
Para aplicar esta derivación auxiliar,
\lstinline|ghost_free_post_len| demuestra que \lstinline|post_len| no depende de las variables auxiliares,
y \lstinline|ghost_free_intro| expresa esta propiedad en la forma requerida por \lstinline|framed_load|:

\begin{lstlisting}
let step_load_ptr (start_ptr : loc) (n : nat) (k : nat)
  : isl_triple (is_ok (post_len start_ptr n k)) (Load "ptr" (Plus (Var "ptr") (Const 1)))
      (variant start_ptr n (k + 1)) =
  let load = framed_load (post_len start_ptr n k) "ptr" (Plus (Var "ptr") (Const 1)) #()
    #(ghost_free_intro (post_len start_ptr n k) (ghost_free_post_len start_ptr n k)) in
  ISL_Consequence (post_len start_ptr n k) (variant start_ptr n (k + 1)) load ()
    (forall_intro (move_requires (lemma_load_ptr start_ptr n k)))
\end{lstlisting}

Con esta última lectura se completa la justificación de una iteración del bucle.
La función \lstinline|proof_loop_step| compone mediante \lstinline|ISL_Seq| las derivaciones \lstinline|step_assume|,
\lstinline|step_load_v|, \lstinline|step_sum|, \lstinline|step_len| y \lstinline|step_load_ptr|.
Cada una adapta las aserciones de su instrucción mediante \lstinline|ISL_Consequence|.
Así se obtiene,
para cada \lstinline|k|,
una derivación del cuerpo desde \lstinline|is_ok (variant start_ptr n k)| hasta \lstinline|is_ok (variant start_ptr n (k + 1))|.
Esta familia proporciona la premisa de \lstinline|ISL_KleeneVariant|:

\begin{lstlisting}
let proof_loop (start_ptr : loc) (n : nat)
  : isl_triple (kleene_pre (variant start_ptr n))
  (Kleene loop_body) (kleene_post (variant start_ptr n)) =
  ISL_KleeneVariant (proof_loop_step start_ptr n)
\end{lstlisting}

Después del bucle,
la aserción \lstinline|post_loop start_ptr n| describe los estados que satisfacen la variante para algún índice entre cero y \lstinline|n|.
El programa comprueba entonces que \lstinline|ptr| contiene la dirección nula.
La derivación \lstinline|step_exit| selecciona como resultado \lstinline|post_exit start_ptr n|,
que describe el recorrido de \lstinline|n| pasos con el puntero nulo.
Desde esta condición,
\lstinline|step_assert| justifica la rama de error de la comprobación sobre los acumuladores.
La aserción \lstinline|post_err| conserva el \emph{store} y el \emph{heap} descritos por \lstinline|post_exit|,
pero exige el modo de terminación \lstinline|Er|:

\begin{lstlisting}
let post_exit (start_ptr : loc) (n : nat) : cond =
  fun s ->
    variant start_ptr n n s /\ s._1 "ptr" == Loc 0 /\
    eval_expr s (Eq (Var "sum") (Var "len")) <> 0

let post_err (start_ptr : loc) (n : nat) : cond =
  fun s -> s._3 == Er /\ post_exit start_ptr n (s._1, s._2, Ok)
\end{lstlisting}
La igualdad entre los acumuladores se deduce de \lstinline|variant start_ptr n n|,
que exige que ambos tengan el valor \lstinline|Nat n|.
Esta igualdad satisface la condición de la rama de error en la representación de la instrucción \lstinline|assert| del programa:

\begin{lstlisting}
let assert_stmt =
  Seq (Assume (Eq (Var "sum") (Var "len"))) Error |||
  Assume (enot (Eq (Var "sum") (Var "len")))
\end{lstlisting}

Como \lstinline|post_exit| establece la igualdad entre \lstinline|sum| y \lstinline|len|,
la derivación selecciona la rama izquierda mediante \lstinline|ISL_ChoiceL|.
El \lstinline|Assume| conserva el estado y \lstinline|Error| cambia el modo de terminación a \lstinline|Er|,
obteniendo así \lstinline|post_err|.

La función \lstinline|proof_assert| compone la comprobación del puntero nulo con la derivación de la rama de error.
Finalmente,
\lstinline|proof_list_sum_len| reúne la inicialización,
el bucle y esta comprobación final:

\begin{lstlisting}
let proof_list_sum_len (start_ptr : loc) (n : nat)
  : isl_triple (is_ok (pre_init start_ptr n)) prog_list_sum_len (post_err start_ptr n) =
  let p_assert = ISL_Consequence (kleene_post (variant start_ptr n)) (post_err start_ptr n)
    (proof_assert start_ptr n) () () in
  let p_loop = ISL_Consequence (kleene_pre (variant start_ptr n)) (post_err start_ptr n)
    (ISL_Seq (proof_loop start_ptr n) p_assert) () () in
  ISL_Consequence (pre_init start_ptr n) (post_err start_ptr n)
    (ISL_Seq (proof_init start_ptr n) p_loop) () ()
\end{lstlisting}
La derivación \lstinline|p_assert| adapta la presunción de \lstinline|proof_assert| a la postcondición del bucle.
Luego,
\lstinline|p_loop| compone ambas derivaciones,
y la última aplicación de \lstinline|ISL_Seq| incorpora la inicialización.
La tripleta justifica que cada estado descrito por \lstinline|post_err| es alcanzable mediante el programa desde algún estado que satisface
\lstinline|is_ok (pre_init start_ptr n)|.

La derivación anterior es paramétrica en la longitud y la dirección inicial.
Para demostrar explícitamente la existencia de un error,
se construye además un testigo para una lista de dos nodos,
ubicados en las direcciones \lstinline|1| y \lstinline|3|.
Ambos contienen el valor \lstinline|Nat 1|,
y los punteros al siguiente nodo se almacenan en las direcciones \lstinline|2| y \lstinline|4|, respectivamente.
El \emph{store} final sitúa el puntero en \lstinline|Loc 0| y deja los dos acumuladores con el valor \lstinline|Nat 2|:

\begin{lstlisting}
let h_two = on loc (fun l ->
  if l = 1 then Full (Nat 1)
  else if l = 2 then Full (Loc 3)
  else if l = 3 then Full (Nat 1)
  else if l = 4 then Full (Loc 0)
  else Unknown)

let st_two = on var (fun z ->
  if z = "ptr" then Loc 0
  else if z = "v" then Nat 1
  else Nat 2)
\end{lstlisting}

El lema \lstinline|two_node_list| demuestra \lstinline|history 1 0 2 2 h_two| para este \emph{heap}.
Primero construye las separaciones que describen la lista mediante \lstinline|list_seg|
y luego justifica las consultas \lstinline|points_to_flex| utilizadas en los dos pasos de \lstinline|history|.
Los \emph{heaps} auxiliares \lstinline|h_two_tl| y \lstinline|h_two_rest| describen,
respectivamente,
el segundo nodo y la celda que almacena el puntero del primer nodo junto con ese segundo nodo.
Con esta propiedad se obtiene el certificado final:

\begin{lstlisting}
let list_reachable ()
  : Lemma (error_reachable (is_ok (pre_init 1 2)) prog_list_sum_len) =
  two_node_list ();
  certify (proof_list_sum_len 1 2) (st_two, h_two, Er)
\end{lstlisting}

Este caso de estudio combina los predicados espaciales recursivos y la Regla del Marco con las asignaciones,
las lecturas de memoria y el razonamiento paramétrico mediante la regla de variante de Kleene.
La derivación muestra que el sistema puede razonar sobre una estructura enlazada de longitud paramétrica
sin escribir manualmente una prueba distinta para cada nodo.

\section{Experimento exploratorio}
\label{sec:experimento-agente}
Como complemento de los casos anteriores,
se realizó un experimento exploratorio,
motivado por Mizzle~\citep{mizzle},
para evaluar si un modelo de lenguaje podía usar la formalización desarrollada en esta tesina para verificar la presencia de errores en los programas.
Mizzle emplea agentes basados en modelos de lenguaje para construir demostraciones comprobables por un asistente de pruebas
y se presenta en la \autoref{sec:relacionado}.

Para ello,
en primer lugar se usó Claude Sonnet 5 para generar un programa en C que contuviera algún error de memoria.
En el programa creado,
hay una función \lstinline|apply_catchup_bonus| que conserva en una variable \lstinline|lowest|
un puntero a un elemento del arreglo dinámico \lstinline|q->items| y luego invoca \lstinline|queue_push|.
Cuando la cola se llena,
\lstinline|queue_push| ejecuta \lstinline|realloc|.
Si este último mueve el arreglo a una nueva ubicación,
\lstinline|lowest| queda apuntando a memoria liberada y la próxima escritura en \lstinline|lowest->score| es un caso de \emph{use after free}.
%\guido{por ahí podemos poner esto en la evaluación, es un caso de uso de tu herramienta. Y acá puede quedar
%sólo la comparación con Mizzle}

En segundo lugar,
se le proporcionó a un agente basado en GPT-5.5 el repositorio que contenía la formalización desarrollada junto con sus casos de uso,
y se le solicitó verificar el programa generado.
El agente identificó el camino de ejecución que conduce al error y produjo un archivo de \fstar{},
donde abstrajo el efecto relevante de \lstinline|realloc|.
Construyó una derivación de la ISL usando las reglas de la formalización
y obtuvo una postcondición de error que conserva el \emph{heap} negativo correspondiente a la dirección invalidada.
La demostración generada fue aceptada por \fstar{}.

Se agregó además un estado final testigo para este programa abstracto.
Las variables \lstinline|"q_size"| y \lstinline|"q_capacity"| conservan el valor cuatro,
utilizado para seleccionar el camino en el que la cola estaba llena.
El valor de \lstinline|"avg"| se abstrae mediante \lstinline|Nat 16|:
no se representa el cálculo en punto flotante,
cuyo resultado en el programa C es $16,25$,
ya que el error depende de la dirección de escritura.
La variable \lstinline|"lowest"| contiene \lstinline|Loc 1|.
La ubicación \lstinline|1| aparece como \lstinline|Empty| y el modo es \lstinline|Er|.
El lema \lstinline|tq_reachable| aplica \lstinline|certify| a la derivación y a este estado,
por lo que establece la existencia del error en el modelo y de la correspondiente ejecución sobre estados completados.

Esta demostración corresponde al programa abstracto que modela la invalidación de la celda referenciada por \lstinline|lowest|
y la escritura posterior sobre ella.
No se mecanizó la traducción desde C ni su correspondencia completa con \lstinline|realloc|.

Este experimento constituye una prueba preliminar de uso y no una evaluación sistemática de la capacidad de los modelos de lenguaje para construir pruebas formales.
En este caso,
muestra que la formalización puede usarse como un lenguaje objetivo en el cual un agente automático
construya certificados de incorrectitud que luego puedan ser comprobados por \fstar{}.

\section{Tamaño del desarrollo y base de confianza}
Las garantías que ofrece una formalización dependen de qué componentes deben aceptarse sin demostración.
En esta sección se detallan esos componentes y se resume el tamaño del desarrollo y el costo de su verificación.

\subsection*{Base de confianza}
El desarrollo no contiene usos de \lstinline|admit| ni de \lstinline|assume|:
todas las definiciones y lemas,
incluidos los de los casos de estudio,
son verificados completamente por \fstar{}.
Las apariciones de \lstinline|assume| en el código corresponden al comando \lstinline|Assume| del lenguaje formalizado y no al mecanismo de \fstar{}.

La formalización usa solo dos axiomas,
ambos provistos por la biblioteca estándar de \fstar{}:
\begin{itemize}
  \item el operador de elección \lstinline|indefinite_description_ghost|,
  que se encuentra en el módulo \lstinline|FStar.IndefiniteDescription|,
  que permite obtener un testigo a partir de una prueba de existencia.
  Se usa en las pruebas de solidez de la IL y de la ISL,
  en el Teorema de la Huella y en algunos casos de estudio,
  para elegir estados iniciales,
  valores anteriores de variables, ubicaciones y particiones del \emph{heap};
  \item la extensionalidad funcional,
  del módulo \lstinline|FStar.FunctionalExtensionality|,
  que se usa para concluir la igualdad de dos \emph{stores} o dos \emph{heaps} a partir de su igualdad punto a punto.
\end{itemize}
Ambos son axiomas estándar y su uso queda restringido al nivel fantasma,
por lo que no afectan al contenido computacional del desarrollo.

En consecuencia,
la base de confianza está formada por el verificador de tipos de \fstar{} (versión \emph{nightly} del 11 de octubre de 2026),
incluida su codificación de las obligaciones de prueba a SMT;
los dos axiomas mencionados;
y el solucionador Z3 (versión 4.13.3),
que descarga dichas obligaciones.

\subsection*{Tamaño y tiempo de verificación}
La \autoref{tab:desarrollo} muestra la cantidad de líneas de cada parte del desarrollo,
sin contar líneas en blanco ni comentarios,
y el tiempo de verificación de cada módulo.
Los tiempos se midieron verificando cada módulo desde cero,
en orden de dependencias,
de modo que el tiempo de un caso de estudio no incluye la verificación de \lstinline|IncSepLogic|.
Cada valor es el promedio de tres ejecuciones en una computadora con procesador Intel Core i5-1135G7 (2,40 GHz) y 16 GB de memoria.

\begin{table}[ht]
  \centering
  \begin{tabular}{llcc}
    \hline
    Módulo & Contenido & Líneas & Tiempo (s) \\
    \hline
    \lstinline|IncLogic|    & IL: lenguaje, semántica, reglas y soundness & 247 & 1,2  \\
    \hline
    \lstinline|IncSepLogic| & Sintaxis, heaps, semántica y metateoría & 285 & \multirow{5}{*}{11,1} \\
                            & Reglas de la ISL                        & 191 & \\
                            & Soundness de la ISL                     & 343 & \\
                            & Teorema de la Huella                    & 514 & \\
                            & Heaps completos y certificados de error &  57 & \\
    \hline
    \lstinline|AllocFree| & Caso de la \autoref{sec:uaf} & 42 & 0,4 \\
    \lstinline|PushBack| & Caso de la \autoref{sec:pushback} & 263 & 3,6 \\
    \lstinline|RandomLoop| & Caso de la \autoref{sec:bucle} & 127 & 1,0 \\
    \lstinline|ListSumLen| & Caso de la \autoref{sec:listas} & 413 & 16,5 \\
    \lstinline|TaskQueue| & Experimento de la \autoref{sec:experimento-agente} & 93 & 0,5 \\
    \hline
    \multicolumn{2}{l}{Total} & 2575 & 34,3 \\
    \hline
  \end{tabular}
  \caption{Tamaño y tiempo de verificación del desarrollo.}
  \label{tab:desarrollo}
\end{table}

El núcleo de la formalización,
compuesto por \lstinline|IncLogic| e \lstinline|IncSepLogic|,
ocupa 1637 líneas.
Dentro de la ISL,
la parte más extensa es la mecanización del Teorema de la Huella (\autoref{sec:isl-huella}),
con 514 líneas.
Los casos de estudio suman 938 líneas.
El más extenso es el recorrido de la lista enlazada,
ya que requiere definir el predicado espacial recursivo y los lemas que relacionan las aserciones de la variante con las lecturas de cada iteración.
Este es también el único módulo que modifica los parámetros por defecto del verificador:
en una de sus pruebas se aumentan el \emph{fuel} y el límite de recursos de Z3 (\lstinline|--fuel 3 --ifuel 1 --z3rlimit 100|).

\section{Resumen de los resultados}
Los cinco casos muestran distintas capacidades del núcleo lógico.
El primero valida el mecanismo básico de \emph{heaps} negativos,
mostrando que la liberación de una celda conserva la información necesaria para justificar un acceso posterior erróneo.

El segundo muestra que esta información puede propagarse a través de una secuencia más compleja que involucra alias y una llamada no determinista.
La derivación selecciona el comportamiento de reasignación y conserva la dirección histórica que usa después el cliente.

El tercer caso muestra que la subaproximación permite concentrarse en una única secuencia de decisiones.
La regla de variante de Kleene permite expresar el efecto de una cantidad arbitraria de iteraciones con una familia de aserciones,
sin materializar un árbol diferente para cada valor concreto.

El cuarto caso,
el recorrido de la lista enlazada,
combina el razonamiento paramétrico con predicados espaciales recursivos.
El predicado \lstinline|history| relaciona las posiciones sucesivas del puntero,
mientras que la Regla del Marco permite justificar las lecturas conservando la información sobre el resto de la lista.

Finalmente,
el último caso consiste en el experimento,
donde el agente utilizó la formalización para construir una prueba de error que luego fue comprobada por \fstar{}. 

En los cinco casos se obtiene además un certificado de existencia de error.

Esta evaluación reporta tiempos de verificación de cada módulo,
pero no establece completitud respecto de todos los programas posibles y no demuestra que las aserciones puedan sintetizarse automáticamente.
Su alcance consiste en validar que las construcciones desarrolladas en el \autoref{cap:formalizacion-isl}
puedan usarse para demostrar la alcanzabilidad de errores no triviales.

%%%%%%%%%%%%%%%%%%%%%%%%%%%%%%%%%%%%%%%%%%%%%%%%%%%%%%%%%%%%%%%%%%%%%%%%%%%%%%%%%%%%
%%%%%%%%%%%%%%%%%%%%%%%%%%%%%%%%%%%% CAPÍTULO 7 %%%%%%%%%%%%%%%%%%%%%%%%%%%%%%%%%%%%
%%%%%%%%%%%%%%%%%%%%%%%%%%%%%%%%%%%%%%%%%%%%%%%%%%%%%%%%%%%%%%%%%%%%%%%%%%%%%%%%%%%%

\chapter{Conclusiones}
\label{cap:conclusiones}
\begin{quote}
``We can only see a short distance ahead, but we can see plenty there that needs to be done.'' \\
\raggedleft --- Alan M. Turing
\end{quote}

\section{Resumen de aportes}
\label{sec:aportes}
En esta tesina se mecanizaron la Lógica de Incorrectitud y su extensión con razonamiento espacial,
la Lógica de Incorrectitud de Separación, utilizando el asistente de pruebas \fstar{}.
El desarrollo permite representar programas,
ejecuciones y derivaciones deductivas dentro del sistema de tipos,
y verificar formalmente que los resultados obtenidos mediante las reglas de ambas lógicas se encuentran respaldados por la semántica operacional.

En una primera etapa se formalizó la IL sobre un lenguaje imperativo simplificado.
Se definió una semántica operacional con el tipo inductivo \lstinline|runsto| y se representaron las reglas deductivas
mediante el tipo inductivo indexado \lstinline|il_triple|.
De esta manera,
la semántica del lenguaje y la derivabilidad dentro de la lógica quedaron expresadas como relaciones diferentes.
Sobre estas definiciones se mecanizó el Teorema de Soundness de la IL.
Este resultado establece que,
para cada estado perteneciente a la postcondición de una tripleta derivable,
existe un estado inicial que satisface la presunción y una ejecución operacional del comando que conecta los dos estados.

La formalización de la IL incluye las reglas correspondientes a asignación,
secuencia, elección no determinista, filtrado mediante \lstinline|Assume|, error e iteración.
También se representó la regla de variante de Kleene,
que permite razonar sobre una cantidad finita pero arbitraria de iteraciones con una familia de aserciones indexada por números naturales.

En una segunda etapa se extendió el desarrollo para formalizar la ISL.
Para ello se incorporó un \emph{heap} como componente independiente del estado y se definieron tres clases de celdas,
\lstinline|Full|, \lstinline|Empty| y \lstinline|Unknown|.
Las celdas \lstinline|Full| representan ubicaciones activas,
las celdas \lstinline|Empty| representan ubicaciones conocidas como no asignadas,
incluidas las que fueron liberadas,
y \lstinline|Unknown| indica que una ubicación no forma parte del fragmento local en el que se está trabajando.
Esta distinción permite conservar información sobre las ubicaciones liberadas y utilizarla para después justificar accesos erróneos.
Sobre este modelo se definieron las aserciones espaciales \lstinline|emp|,
\lstinline|points_to|, \lstinline|points_to_flex| y \lstinline|points_to_empty|,
junto con la conjunción separadora.
También se extendieron la sintaxis y la semántica operacional con instrucciones de reserva,
lectura, escritura y liberación de memoria.

Las reglas deductivas de la ISL se representaron con el tipo inductivo \lstinline|isl_triple|.
Además de las reglas para las operaciones primitivas,
se formalizó la Regla del Marco,
con la condición de independencia entre la aserción usada como marco y las variables modificadas por el comando.

Sobre esta formalización se mecanizó el Teorema de Soundness de la ISL.
Al igual que en el caso anterior,
el resultado relaciona cada estado final descrito por su derivación con un estado inicial y una ejecución operacional concreta.
La demostración del caso de la Regla del Marco requirió probar que una ejecución local puede extenderse con un \emph{heap} disjunto que no es modificado.

También se mecanizó el Teorema de la Huella,
que caracteriza las ejecuciones como extensiones de transiciones minimales mediante marcos disjuntos.
Además,
se demostró que las ejecuciones se preservan al completar los \emph{heaps} y se incorporaron certificados de existencia de errores.

Finalmente,
la formalización fue evaluada con cuatro casos de estudio y un experimento exploratorio.
Se construyeron derivaciones para un error \emph{use after free},
un puntero colgante producido por una reasignación no determinista,
una ejecución con una gran cantidad de iteraciones y decisiones no deterministas,
y el recorrido de una lista enlazada.
Además,
se verificó la abstracción del camino de error del programa generado en C en el experimento exploratorio.
Cada caso cuenta con un certificado de existencia de error.

El desarrollo obtenido constituye un núcleo lógico mecanizado y la formalización garantiza el respaldo operacional de las derivaciones construidas.

\section{Trabajo relacionado}
\label{sec:relacionado}
Existen distintos trabajos relacionados con esta tesina.
Los presentamos a continuación,
comparándolos con el desarrollo realizado en esta tesina.

\subsection*{Mecanización de la HL para código máquina}
\citet{mchoare} presentan un framework mecanizado de estilo Hoare para razonar sobre programas escritos en código máquina.
La lógica se construye sobre una semántica operacional que representa características propias de máquinas reales,
como la memoria finita,
la coexistencia de código y datos en un mismo espacio de direcciones,
los registros de estado y el riesgo de modificar registros como el contador de programa,
el puntero de pila y el registro de retorno.
El desarrollo fue mecanizado en HOL4~\citep{hol4}.
%\guido{yo diria solamente formalizado o mecanizado en HOL4}

El objetivo de ese trabajo es permitir la especificación y verificación de propiedades funcionales
y de uso de recursos directamente sobre programas de bajo nivel.
Aunque el modelo conserva detalles de la máquina,
las especificaciones se expresan mediante una representación flexible del estado inspirada en el razonamiento local de la SL.

El uso de una semántica operacional explícita y de un asistente de pruebas para justificar las reglas de razonamiento
es un enfoque que comparte ese trabajo con la tesina actual.
En los dos casos,
las propiedades del sistema lógico no dependen solo de una argumentación informal,
sino que se encuentran respaldadas por una mecanización.

La diferencia principal se encuentra en la interpretación de las tripletas.
El trabajo de Myreen y Gordon se orienta al razonamiento de correctitud.
En esta tesina, en cambio,
las tripletas de IL e ISL tienen una interpretación subaproximada y se usan para demostrar que los estados descritos por la postcondición son alcanzables.

También difiere el nivel de abstracción.
Myreen y Gordon trabajan directamente con código máquina y tienen que representar varios detalles de la arquitectura.
La presente formalización usa un lenguaje imperativo reducido,
lo que permite concentrar el desarrollo en la estructura de la IL y de la ISL y las demostraciones de sus propiedades metateóricas.

\subsection*{SL para la semántica de paso pequeño de Cminor}
\citet{cminor} presentan una formalización de la SL para Cminor,
un lenguaje imperativo intermedio usado dentro de la infraestructura de CompCert.
El trabajo redefine Cminor para facilitar el razonamiento con la HL y da una semántica operacional de paso pequeño,
en lugar de solo usar la semántica de paso grande del compilador.
Esta elección se encuentra motivada por la posibilidad de extender luego el lenguaje con concurrencia.

El desarrollo incluye una demostración mecanizada en Coq~\citep{coq} de la solidez de la SL respecto de esa semántica.
Los autores presentan el resultado como la primera demostración mecanizada de una SL a gran escala respecto de una semántica de paso pequeño.

Este trabajo es cercano a esta tesina por su metodología.
En los dos desarrollos se define una semántica operacional,
se formaliza un sistema de razonamiento espacial y se demuestra que las reglas de la lógica son sólidas respecto de dicha semántica.
También comparten el uso de una representación explícita del \emph{heap},
aserciones espaciales y razonamiento local con marcos.

Sin embargo,
las lógicas formalizadas tienen interpretaciones diferentes.
Appel y Blazy usan la SL para demostrar propiedades de correctitud y seguridad de programas Cminor.
Esta tesina formaliza la ISL,
cuyo objetivo es demostrar la existencia de ejecuciones que alcanzan determinadas postcondiciones, incluidos estados de error.
Esta diferencia se refleja en los teoremas de solidez.
En el razonamiento de correctitud,
la prueba debe garantizar que toda ejecución respete la postcondición.
En la interpretación subaproximada,
debe demostrarse que cada estado descrito por la postcondición tiene al menos una ejecución que lo alcanza desde algún estado que satisface la presunción.

Otra diferencia es la incorporación de \emph{heaps} negativos.
La formalización desarrollada en esta tesina distingue entre celdas activas,
celdas liberadas y direcciones que no pertenecen al fragmento local.
Esta representación permite justificar accesos erróneos a memoria liberada,
un fin diferente al perseguido por la SL de Cminor.

\subsection*{IsaBIL: verificación de correctitud e incorrectitud de binarios}
\citet{isabil} presentan IsaBIL,
un framework desarrollado en Isabelle/HOL para verificar propiedades de programas binarios.
La herramienta usa \emph{Binary Analysis Platform} (BAP)~\citep{bap} para traducir los binarios al lenguaje intermedio BIL
y genera un entorno dentro de Isabelle/HOL en el que pueden construirse pruebas basadas tanto en HL como en IL.
IsaBIL incluye una formalización de la sintaxis y la semántica operacional de BIL.
Las lógicas de correctitud e incorrectitud se definen de forma paramétrica respecto de una relación operacional,
lo que permite reutilizar sus reglas con diferentes modelos de estado.
Esta estructura se implementa mediante los \emph{locales} de Isabelle/HOL,
contextos parametrizados que agrupan hipótesis,
definiciones y resultados reutilizables.

Además,
el framework incorpora mecanismos de automatización destinados a reducir el esfuerzo de prueba.
Los autores aplican estas técnicas a distintos ejemplos y demuestran, con la IL,
la presencia de vulnerabilidades de \emph{double free},
incluida una vulnerabilidad de la biblioteca cURL.

Este trabajo es muy cercano a esta tesina en cuanto al sistema lógico estudiado.
Los dos desarrollos mecanizan la IL,
definen una semántica operacional y usan un asistente de pruebas para verificar que los razonamientos deductivos tienen respaldo semántico.

Su distinción está en el objetivo y nivel de abstracción.
IsaBIL busca proporcionar un entorno semiautomatizado para verificar ejecuciones reales,
incluso cuando no se tiene el código fuente.
Por esa razón, tiene que representar instrucciones,
registros, palabras de máquina, saltos y detalles de arquitecturas concretas.
Esta tesina usa un lenguaje más chico para estudiar directamente la estructura metateórica de la IL y la ISL.

La diferencia más importante se encuentra en el razonamiento espacial.
La versión de IsaBIL presentada en ese trabajo no incorpora la SL ni la ISL.
Los autores explican que la memoria de BIL se representa con un historial de actualizaciones y que esa representación
no es directamente compatible con la separación del \emph{heap},
por lo que dejan la ISL como una extensión futura.

\subsection*{Automatización del razonamiento espacial con VeriFast}
\citet{verifast} presentan VeriFast,
una herramienta de verificación para programas C y Java basada en SL y ejecución simbólica.
El usuario anota los programas con contratos,
predicados espaciales y otros elementos auxiliares,
mientras que la herramienta ejecuta simbólicamente los comandos y verifica las obligaciones de prueba.
VeriFast admite programas secuenciales y concurrentes y puede utilizarse para demostrar la ausencia de diferentes errores de ejecución,
como accesos inválidos a memoria,
desreferencias de punteros nulos y accesos fuera de los límites de arreglos.
Su sistema de especificación permite representar permisos y predicados definidos por el usuario para describir estructuras dinámicas.

Este trabajo se relaciona con esta tesina porque los dos usan aserciones espaciales y conjunción separadora para razonar sobre fragmentos disjuntos de memoria.
En ambos casos,
los predicados inductivos permiten representar estructuras enlazadas sin enumerar individualmente todas sus celdas.

La diferencia principal reside nuevamente en la interpretación lógica.
VeriFast usa SL para verificar que los programas respeten sus contratos y no produzcan determinados errores.
La ISL formalizada en esta tesina usa subaproximaciones para demostrar que existen ejecuciones que alcanzan estados determinados,
incluidos estados que representan accesos inválidos.

También difieren en el grado de automatización.
VeriFast selecciona automáticamente los pasos de ejecución simbólica y aplica las reglas necesarias para procesar la estructura del programa.
Aunque el usuario debe dar contratos,
invariantes y predicados auxiliares,
no construye manualmente un árbol de derivación para cada instrucción.
En esta tesina,
las derivaciones de \lstinline|isl_triple| se construyen de forma explícita con los constructores del tipo inductivo.
\fstar{} y Z3 verifican los refinamientos y condiciones laterales,
pero no eligen de manera autónoma los caminos de ejecución,
las reglas ni las aserciones intermedias.

Por otra parte,
VeriFast es una herramienta destinada a la verificación de programas concretos,
mientras que el objetivo de la tesina es mecanizar el núcleo de una lógica y demostrar sus propiedades metateóricas.
La comparación permite mostrar una posible línea de investigación futura:
usar las definiciones y los teoremas mecanizados como fundamento de un procedimiento que automatice la construcción de derivaciones de ISL.

\subsection*{Mizzle: mecanización de una lógica de incorrectitud concurrente}
\citet{mizzle} presentan Mizzle, una lógica de incorrectitud con separación para un subconjunto de OCaml 5,
modelando concurrencia bajo un modelo de consistencia secuencial.
El desarrollo se encuentra mecanizado en Rocq,
anteriormente denominado Coq,
sobre el framework Iris y cuenta con demostraciones de solidez y completitud.
Mientras que la solidez garantiza que una derivación de Mizzle está respaldada por una ejecución real,
la completitud establece que toda ejecución incorrecta considerada por el modelo puede representarse mediante una derivación de la lógica.
Durante el desarrollo de esta tesina,
los autores publicaron una segunda versión del artículo,
en octubre de 2026,
que amplía la evaluación de Mizzle.

Este trabajo es el más cercano a esta tesina porque los dos mecanizan razonamiento subaproximado
junto con herramientas de SL y relacionan las derivaciones lógicas con una semántica operacional.
Sin embargo,
los dos desarrollos persiguen objetivos diferentes.
Mizzle se enfoca en un lenguaje más expresivo,
con funciones de orden superior y concurrencia,
mientras que en esta tesina se usa un lenguaje imperativo simplificado para estudiar la estructura metateórica de la ISL.
Por otro lado,
Mizzle no incorpora liberación de memoria,
mientras que la formalización desarrollada en esta tesina modela la liberación con \emph{heaps} negativos
y usa esta información para demostrar errores de \emph{use after free} o \emph{double free}.

Un aspecto interesante de Mizzle es su utilización como lenguaje de prueba para agentes basados en modelos de lenguaje.
La propuesta consiste en no confiar en el razonamiento interno de un modelo de lenguaje de gran tamaño (LLM, por sus siglas en inglés)
para decidir si un error es real,
sino exigir que el modelo genere una demostración que pueda ser comprobada por Rocq.
En su evaluación,
los autores usan Claude Opus 5.5 junto con una guía de prueba,
una descripción de la arquitectura del proyecto y una biblioteca de tácticas.

Esta propuesta motivó el experimento exploratorio presentado en la \autoref{sec:experimento-agente}.

\subsection*{Resumen}
Al igual que varios de los desarrollos anteriores,
esta formalización se concentra en dar una representación mecanizada de la semántica,
las reglas lógicas y sus propiedades.
Al igual que VeriFast,
incorpora aserciones espaciales y predicados sobre memoria mutable.
Lo que caracteriza a este trabajo es la mecanización en \fstar{} de la interpretación subaproximada de la IL
junto con el razonamiento local de la SL,
incluyendo el tratamiento de la liberación de memoria mediante \emph{heaps} negativos.

A diferencia de IsaBIL y VeriFast,
el desarrollo no incorpora un procedimiento integrado de automatización de las derivaciones ni pretende analizar un lenguaje industrial.
Su aporte consiste en dar un núcleo formal en el que las derivaciones de la IL e ISL son objetos explícitos,
se verifican sus propiedades de solidez y localidad y se mecaniza el Teorema de la Huella dentro de \fstar{}.

\section{Trabajo futuro}
\label{sec:futuro}
Detallamos a continuación algunas líneas de investigación que permitirían extender el desarrollo realizado.

\subsection*{Automatización de las derivaciones}
Aunque el experimento motivado por Mizzle,
presentado en la \autoref{sec:experimento-agente},
mostró que un agente puede construir una derivación,
la formalización no incorpora un procedimiento de búsqueda:
las reglas que se aplican,
las aserciones intermedias y los constructores deben proporcionarse explícitamente.
Una extensión consiste en desarrollar un procedimiento que recorra la sintaxis del programa y construya automáticamente derivaciones de \lstinline|isl_triple|.
Este procedimiento podría calcular postcondiciones,
seleccionar ramas de elección no determinista,
descubrir caminos que llevan a errores y generar las obligaciones lógicas.

En el caso espacial,
técnicas de biabducción \citep{biabd} podrían usarse para inferir los recursos requeridos por una operación y los marcos que deben conservarse.
Esto permitiría aproximar la formalización a una herramienta de análisis estático sin abandonar las garantías mecanizadas del núcleo lógico.

\subsection*{Extensión del lenguaje y del modelo de memoria}
El lenguaje formalizado tiene las construcciones necesarias para estudiar las reglas centrales de la lógica,
pero no incorpora llamadas recursivas,
estructuras,
arreglos ni aritmética general de punteros.
Agregar estas características permitiría representar programas más cercanos a lenguajes de uso real.
También podría desarrollarse un modelo de memoria con bloques de distintos tamaños,
direcciones interiores,
alineamiento y una representación más detallada de los valores almacenados.

Cada extensión requeriría agregar las reglas operacionales y deductivas correspondientes,
y demostrar la solidez de las nuevas reglas deductivas respecto de la semántica ampliada.

\subsection*{Generación de trazas de error}
La función \lstinline|soundness| proporciona evidencia de que cada estado final derivado en error es alcanzable desde algún estado inicial.
Esta construcción se realiza en el contexto fantasma de la prueba,
por lo que no constituye directamente un generador ejecutable de trazas.
Una herramienta práctica podría complementar esta evidencia con un procedimiento que construya explícitamente los testigos y produzca una traza legible.

La traza podría indicar las instrucciones ejecutadas, las ramas no deterministas seleccionadas,
los valores relevantes del \emph{store} y las modificaciones realizadas sobre el \emph{heap}.
De esta manera,
la demostración formal podría usarse también para producir una explicación concreta del error.

\subsection*{Evaluación experimental}
La evaluación realizada en esta tesina es cualitativa y se concentra en la expresividad de la formalización.
Una implementación con más automatización permitiría estudiar experimentalmente su comportamiento.

Podrían medirse los tiempos de verificación,
el tamaño de las condiciones generadas, el consumo de memoria,
la participación del solucionador SMT y la variación de los valores frente a programas de distintas dimensiones.
También sería posible comparar diferentes estrategias de exploración y estudiar hasta qué punto
la construcción de derivaciones subaproximadas permite controlar el crecimiento producido por las elecciones no deterministas y los bucles.

\subsection*{Concurrencia}
Otra posible extensión es incorporar programas concurrentes.
Esto requeriría incorporar construcciones de composición paralela y una semántica operacional
que represente las posibles intercalaciones entre los pasos de los distintos hilos.

\citet{cisl} presentan la Lógica de Incorrectitud de Separación Concurrente (CISL),
que extiende la ISL al razonamiento subaproximado sobre programas concurrentes.
Mecanizar las propiedades de concurrencia en \fstar{} permitiría estudiar errores que no aparecen en programas secuenciales,
como \emph{race conditions}, \emph{deadlocks} y fallas dependientes del orden de ejecución.

% \appendix
% \label{sec:appendix}

%\nocite{*}
% Variante local de plainnat: usa "y" entre autores y editores.
\bibliographystyle{plainnat-es}
\bibliography{tesina}

\end{document}
